% Lesson 46 — Grammar: Review phrasal verbs in simple expressions of purpose — Anglais Tle SH — Cours signé OnBuch+
% Programme officiel MINESEC. Compiler avec : tectonic EN46-grammar-review-phrasal-verbs-in-simple-expressions-of-p.tex
\newcommand{\DOCMATIERE}{Anglais}
\newcommand{\DOCNIVEAU}{Tle SH}
\newcommand{\DOCMODULE}{Chapter 4 : Citizenship and Human Rights}
\newcommand{\DOCLECON}{EN46}
\newcommand{\DOCTITRE}{Lesson 46 — Grammar: Review phrasal verbs in simple expressions of purpose}
% OnBuch+ — préambule « cours signés » ENRICHI (compilé avec tectonic / XeLaTeX)
% Système visuel « Pop » d'OnBuch+ : fond crème (jamais blanc), bordures encre
% épaisses, ombres « dures » décalées, titres Archivo Black en capitales.
% Version MATHÉMATIQUES : graphes (pgfplots/tikz), boîtes pédagogiques
% (définition, propriété, méthode, attention, le savais-tu, exercice résolu,
% exercice, TP, bilan), page de garde et signature OnBuch+ ancrée en bas.
%
% Macros attendues dans main.tex AVANT \input{preamble} :
%   \DOCMATIERE  \DOCNIVEAU  \DOCMODULE  \DOCLECON (numéro)  \DOCTITRE
\documentclass[11pt]{article}

\usepackage{fontspec}
\usepackage[english,french]{babel} % français = langue principale ; l'anglais via \eng{..} / textebox
\usepackage[a4paper,margin=2cm,top=3.4cm,bottom=2.8cm,headheight=1.4cm,footskip=1.3cm]{geometry}
\usepackage{amsmath,amssymb,amsfonts}
\usepackage{mathtools} % \xmapsto, \coloneqq... souvent utilisés par les modèles
\usepackage{cancel} % simplifications barrées (bilans de forces)
% \abs{z} (module, valeur absolue) et \norm{u} (norme), utilisés par les modèles
\providecommand{\abs}{}\let\abs\relax\DeclarePairedDelimiter{\abs}{\lvert}{\rvert}
\providecommand{\norm}{}\let\norm\relax\DeclarePairedDelimiter{\norm}{\lVert}{\rVert}
\usepackage[table]{xcolor} % [table] : fournit aussi \rowcolors (alternance de lignes)
\usepackage{pagecolor}
\usepackage{tikz}
\usetikzlibrary{arrows.meta,positioning,calc,shapes.geometric,shapes.misc,shapes.arrows,decorations.pathmorphing,decorations.pathreplacing,decorations.markings,patterns,shadows,mindmap,trees,fit,backgrounds,matrix,intersections,angles,quotes}
\usepackage{pgfplots}
\usepackage[version=4]{mhchem} % équations nucléaires \ce{^{235}_{92}U}
\usepackage[european,siunitx]{circuitikz} % schémas électriques
\pgfplotsset{compat=1.18}
\usepgfplotslibrary{fillbetween,groupplots} % aires entre courbes (calcul intégral)
\usepackage{enumitem}
\usepackage{fancyhdr}
% [most] charge toutes les bibliothèques tcolorbox (listings, minted, raster,
% documentation...) et gonfle inutilement la mémoire TeX sur un document long
% (~300 boîtes) : on ne charge que ce qui est réellement utilisé.
\usepackage[skins,breakable]{tcolorbox}
\usepackage{graphicx}
\usepackage{colortbl}
\usepackage{booktabs}
\usepackage{tabularx}
\usepackage{diagbox} % case d'en-tête diagonale des tableaux à double entrée
% Colonnes extensibles centrée / gauche / droite (C, L, R), souvent utilisées par les modèles
\newcolumntype{C}{>{\centering\arraybackslash}X}
\newcolumntype{L}{>{\raggedright\arraybackslash}X}
\newcolumntype{R}{>{\raggedleft\arraybackslash}X}
\usepackage{array}
\usepackage{multicol}
\usepackage{multirow}
\usepackage{siunitx}
% parse-numbers=false : accepte aussi \SI{2,5 \times 10^{-3}}{...}
\sisetup{locale=FR,output-decimal-marker={,},group-minimum-digits=4,parse-numbers=false}
\DeclareSIUnit\ohm{\ensuremath{\symup{\Omega}}} % U+2126 absent de la police
\DeclareSIUnit\division{div} % réglages d'oscilloscope (ms/div, V/div)
\DeclareSIUnit\tr{tr} % tours (vitesse de rotation, ex. \SI{1500}{\tr\per\minute})
\DeclareSIUnit\molar{M} % concentration molaire (ex. \SI{3}{\molar}, \SI{1}{\micro\molar})
\DeclareSIUnit\an{an} % année (le modèle confond parfois avec \year, un compteur TeX) — ex. \SI{5}{\kilo\meter\per\an}
\DeclareSIUnit\FCFA{FCFA} % franc CFA (ex. \SI{500}{\FCFA\per\kilo\gram})
\DeclareSIUnit\degreeBrix{\textdegree Brix} % teneur en sucre d'un sirop/jus (réfractométrie)
\DeclareSIUnit\month{mois} % le modèle utilise parfois \month, un compteur TeX, comme unité de durée
\usepackage{qrcode}
% \degree : commande fréquemment utilisée par les modèles pour un angle ou une
% température en dehors de \SI{}{} (non fournie par les paquets chargés).
\providecommand{\degree}{^{\circ}}
% \qed (fin de démonstration), utilisé par les modèles sans amsthm
\providecommand{\qed}{\hfill$\square$}
% \barre : double barre des valeurs interdites dans un tableau de variations (tkz-tab)
\providecommand{\barre}{\ensuremath{\|}}
% environnement proof (amsthm non chargé) utilisé par les modèles
\ifdefined\proof\else\newenvironment{proof}[1][Démonstration]{\par\noindent\textit{#1.}\ }{\qed\par}\fi
\usepackage{lastpage}
\usepackage{hyperref}
\hypersetup{hidelinks}

% ============================================================
% Palette « Pop » OnBuch+ — mêmes hex que la classe P (plus_ui.dart)
% ============================================================
\definecolor{poporange}{HTML}{F59321}
\definecolor{popdark}{HTML}{E07A0C}
\definecolor{popcream}{HTML}{FFF6E8}
\definecolor{popink}{HTML}{1C1714}
\definecolor{poppurple}{HTML}{7A5AE0}
\definecolor{popgreen}{HTML}{2E9E6A}
\definecolor{poppink}{HTML}{E0457B}
\definecolor{popblue}{HTML}{2F7DE1}
\definecolor{popgold}{HTML}{C98A12}
\definecolor{popmuted}{HTML}{8A7F76}
\definecolor{popline}{HTML}{EADFD2}
\definecolor{popgray}{HTML}{8A7F76}
\colorlet{popgrey}{popgray}
% Alias tolérants (le modèle nomme parfois les couleurs différemment)
\colorlet{popwhite}{white}
\colorlet{popblack}{popink}
\colorlet{popred}{poppink}
\colorlet{popyellow}{popgold}
% Teintes claires (fonds de tableaux/encadrés) — le modèle nomme parfois
% les couleurs avec un suffixe L (ex. popblueL) sans les avoir définies.
\colorlet{poporangeL}{poporange!20}
\colorlet{popdarkL}{popdark!20}
\colorlet{poppurpleL}{poppurple!20}
\colorlet{popgreenL}{popgreen!20}
\colorlet{poppinkL}{poppink!20}
\colorlet{popblueL}{popblue!20}
\colorlet{popgoldL}{popgold!20}
\colorlet{popredL}{popred!20}
\colorlet{popgrayL}{popgray!20}
% Teintes claires pour fonds de tableaux / zones
\colorlet{popblueL}{popblue!10!white}
\colorlet{poporangeL}{poporange!14!white}
\colorlet{popgreenL}{popgreen!12!white}
\colorlet{poppurpleL}{poppurple!12!white}
\colorlet{poppinkL}{poppink!12!white}
\colorlet{popgoldL}{popgold!14!white}

\pagecolor{popcream}

% ============================================================
% Typographie
% ============================================================
\defaultfontfeatures{Ligatures=TeX}
\setmainfont{PlusJakartaSans-Medium.ttf}[Path=../../../fonts/,BoldFont=PlusJakartaSans-Bold.ttf]
% Maths en sans-serif (Fira Math) avec chiffres et lettres droites de Plus
% Jakarta, pour que les formules se fondent dans le texte.
\usepackage[math-style=ISO,bold-style=ISO]{unicode-math}
\setmathfont{FiraMath-Regular.otf}[Path=../../../fonts/]
\setmathfont{PlusJakartaSans-Medium.ttf}[Path=../../../fonts/,range={up/num,up/{Latin,latin},\mathup}]
\usepackage{icomma} % virgule décimale sans espace en mode math
% Caractères mathématiques Unicode absents des polices de texte (Plus Jakarta,
% Archivo Black) : rendus par leur équivalent mathématique, en texte comme en
% formule — sinon ils s'affichent en carré vide (ex. « Arithmétique dans ℤ »).
\usepackage{newunicodechar}
\newunicodechar{ℝ}{\ensuremath{\mathbb{R}}}
\newunicodechar{ℤ}{\ensuremath{\mathbb{Z}}}
\newunicodechar{ℂ}{\ensuremath{\mathbb{C}}}
\newunicodechar{ℕ}{\ensuremath{\mathbb{N}}}
\newunicodechar{ℚ}{\ensuremath{\mathbb{Q}}}
\newunicodechar{𝔻}{\ensuremath{\mathbb{D}}}
\newunicodechar{α}{\ensuremath{\alpha}}
\newunicodechar{β}{\ensuremath{\beta}}
\newunicodechar{γ}{\ensuremath{\gamma}}
\newunicodechar{δ}{\ensuremath{\delta}}
\newunicodechar{ε}{\ensuremath{\varepsilon}}
\newunicodechar{θ}{\ensuremath{\theta}}
\newunicodechar{λ}{\ensuremath{\lambda}}
\newunicodechar{μ}{\ensuremath{\mu}}
\newunicodechar{σ}{\ensuremath{\sigma}}
\newunicodechar{τ}{\ensuremath{\tau}}
\newunicodechar{φ}{\ensuremath{\varphi}}
\newunicodechar{ψ}{\ensuremath{\psi}}
\newunicodechar{ω}{\ensuremath{\omega}}
\newunicodechar{Σ}{\ensuremath{\Sigma}}
% majuscules produites par \MakeUppercase dans les titres (α-aminés -> Α)
\newunicodechar{Α}{\ensuremath{\alpha}}
\newunicodechar{Β}{\ensuremath{\beta}}
\newunicodechar{Γ}{\ensuremath{\Gamma}}
\newunicodechar{Δ}{\ensuremath{\Delta}}
\newunicodechar{Ε}{\ensuremath{\varepsilon}}
\newunicodechar{Θ}{\ensuremath{\Theta}}
\newunicodechar{Λ}{\ensuremath{\Lambda}}
\newunicodechar{Μ}{\ensuremath{\mu}}
\newunicodechar{Τ}{\ensuremath{\tau}}
\newunicodechar{Φ}{\ensuremath{\Phi}}
\newunicodechar{Ψ}{\ensuremath{\Psi}}
\newunicodechar{Ω}{\ensuremath{\Omega}}
\newunicodechar{↦}{\ensuremath{\mapsto}}
\newunicodechar{∈}{\ensuremath{\in}}
\newunicodechar{∉}{\ensuremath{\notin}}
\newunicodechar{∧}{\ensuremath{\wedge}}
\newunicodechar{⇔}{\ensuremath{\Leftrightarrow}}
\newunicodechar{⟺}{\ensuremath{\Longleftrightarrow}}
\newunicodechar{⇒}{\ensuremath{\Rightarrow}}
\newunicodechar{⟹}{\ensuremath{\Longrightarrow}}
\newunicodechar{∘}{\ensuremath{\circ}}
\newunicodechar{∪}{\ensuremath{\cup}}
\newunicodechar{∩}{\ensuremath{\cap}}
\newunicodechar{≡}{\ensuremath{\equiv}}
\newunicodechar{⊂}{\ensuremath{\subset}}
\newunicodechar{⊥}{\ensuremath{\perp}}
\newunicodechar{∃}{\ensuremath{\exists}}
\newunicodechar{∀}{\ensuremath{\forall}}
\newunicodechar{∛}{\ensuremath{\sqrt[3]{\,}}}
\newunicodechar{∜}{\ensuremath{\sqrt[4]{\,}}}
\newunicodechar{⃗}{\ensuremath{{}^{\to}}}
\newunicodechar{ⁿ}{\ifmmode^{n}\else\textsuperscript{\ensuremath{n}}\fi}
\newunicodechar{ˣ}{\ifmmode^{x}\else\textsuperscript{\ensuremath{x}}\fi}
\newunicodechar{ᵇ}{\ifmmode^{b}\else\textsuperscript{\ensuremath{b}}\fi}
\newunicodechar{ʳ}{\ifmmode^{r}\else\textsuperscript{\ensuremath{r}}\fi}
\newunicodechar{ᵘ}{\ifmmode^{u}\else\textsuperscript{\ensuremath{u}}\fi}
\newunicodechar{ʸ}{\ifmmode^{y}\else\textsuperscript{\ensuremath{y}}\fi}
\newunicodechar{ᵃ}{\ifmmode^{a}\else\textsuperscript{\ensuremath{a}}\fi}
\newunicodechar{ᵉ}{\ifmmode^{e}\else\textsuperscript{\ensuremath{e}}\fi}
\newunicodechar{ᵏ}{\ifmmode^{k}\else\textsuperscript{\ensuremath{k}}\fi}
\newunicodechar{ᵖ}{\ifmmode^{p}\else\textsuperscript{\ensuremath{p}}\fi}
\newunicodechar{ᴺ}{\ifmmode^{N}\else\textsuperscript{\ensuremath{N}}\fi}
\newunicodechar{ᶜ}{\ifmmode^{c}\else\textsuperscript{\ensuremath{c}}\fi}
\newunicodechar{ʰ}{\ifmmode^{h}\else\textsuperscript{\ensuremath{h}}\fi}
\newunicodechar{ᵠ}{\ifmmode^{q}\else\textsuperscript{\ensuremath{q}}\fi}
\newunicodechar{ₙ}{\ifmmode_{n}\else\textsubscript{\ensuremath{n}}\fi}
\newunicodechar{ₐ}{\ifmmode_{a}\else\textsubscript{\ensuremath{a}}\fi}
\newunicodechar{ᵢ}{\ifmmode_{i}\else\textsubscript{\ensuremath{i}}\fi}
\newunicodechar{ᵣ}{\ifmmode_{r}\else\textsubscript{\ensuremath{r}}\fi}
\newunicodechar{ₖ}{\ifmmode_{k}\else\textsubscript{\ensuremath{k}}\fi}
\newunicodechar{ᵦ}{\ifmmode_{\beta}\else\textsubscript{\ensuremath{\beta}}\fi}
\newunicodechar{ₓ}{\ifmmode_{x}\else\textsubscript{\ensuremath{x}}\fi}
\newfontfamily\popdisplay{ArchivoBlack-Regular.ttf}[Path=../../../fonts/]
\newfontfamily\poplogo{SpaceGrotesk-Bold.ttf}[Path=../../../fonts/]
\newfontfamily\popsemi{PlusJakartaSans-SemiBold.ttf}[Path=../../../fonts/]
\newfontfamily\popxbold{PlusJakartaSans-ExtraBold.ttf}[Path=../../../fonts/]
\newfontfamily\popxboldls{PlusJakartaSans-ExtraBold.ttf}[Path=../../../fonts/,LetterSpace=3.0]

\setlist[enumerate]{leftmargin=1.7em,itemsep=5pt,topsep=4pt}
\setlist[itemize]{leftmargin=1.7em,itemsep=4pt,topsep=4pt,label={\color{poporange}\textbullet}}
\setlength{\parindent}{0pt}
\setlength{\parskip}{5pt}
\renewcommand{\arraystretch}{1.25}

% pgfplots : style Pop par défaut
\pgfplotsset{
  every axis/.append style={axis line style={popink,line width=1pt},tick style={popink},
    grid style={popline,line width=0.5pt},label style={font=\small},tick label style={font=\footnotesize}},
  popcurve/.style={poporange,line width=1.8pt,no markers,smooth},
}
% Même style disponible dans un \draw TikZ simple (hors pgfplots)
\tikzset{popcurve/.style={poporange,line width=1.8pt}}
\tikzset{popgrid/.style={popline,very thin}}
\colorlet{popcurve}{poporange} % le nom du style est parfois utilisé comme couleur

% ============================================================
% Pill / squiggle
% ============================================================
\newcommand{\poppill}[3]{%
  \tikz[baseline=-0.55ex]\node[draw=popink,line width=1.2pt,rounded corners=2.6pt,%
    fill=#2,inner xsep=6.5pt,inner ysep=3pt]{%
    \popxboldls\fontsize{7}{7}\selectfont\color{#3}\MakeUppercase{#1}};%
}
\newcommand{\popsquiggle}[1]{%
  \tikz[baseline=-0.4ex]{
    \draw[#1,line width=2.6pt,line cap=round]
      plot [smooth] coordinates {(0,0.09) (0.34,0.01) (0.68,0.21) (1.02,0.01) (1.36,0.10)};
  }%
}
% Mot-clé surligné (marqueur orange sous le texte)
\newcommand{\cle}[1]{{\popxbold\color{popdark}#1}}

% ============================================================
% En-tête / pied
% ============================================================
\pagestyle{fancy}
\fancyhf{}
\renewcommand{\headrulewidth}{0pt}
\renewcommand{\footrulewidth}{0pt}
\lhead{\raisebox{-0.15ex}{\poplogo\fontsize{13}{13}\selectfont\color{popink}OnBuch\color{poporange}+}}
\rhead{\poppill{\DOCMATIERE\ \textbullet\ \DOCNIVEAU\ \textbullet\ Leçon \DOCLECON}{white}{popink}}
\lfoot{\popxbold\fontsize{7.6}{7.6}\selectfont\color{popmuted}\MakeUppercase{Cours signé OnBuch+}\ \textbullet\ {\popsemi\DOCTITRE}}
\rfoot{\tikz[baseline=-0.6ex]\node[rounded corners=8.5pt,fill=popink,minimum height=17pt,minimum width=17pt,inner xsep=5pt,inner sep=0]{\popxbold\fontsize{8}{8}\selectfont\color{popcream}\thepage\,/\,\pageref*{LastPage}};}

% ============================================================
% Titres
% ============================================================
\newcommand{\chaptitre}[2]{%
  \begin{tcolorbox}[enhanced,colback=poporange,colframe=popink,boxrule=2.6pt,arc=11pt,
      drop shadow=popink,left=15pt,right=15pt,top=12pt,bottom=11pt,before skip=3pt,after skip=18pt]
    \poppill{#2}{popink}{popcream}\par\vspace{9pt}
    {\raggedright\hyphenpenalty=10000\exhyphenpenalty=10000\popdisplay\fontsize{22}{24}\selectfont\color{popcream}\MakeUppercase{#1}\par}
  \end{tcolorbox}
}
% Section numérotée : pastille orange avec le numéro + titre Archivo
\newcounter{popsec}
\newcommand{\coursec}[1]{%
  \stepcounter{popsec}\setcounter{popsub}{0}%
  \par\vspace{16pt}\noindent
  \begin{minipage}{\linewidth}
  \tikz[baseline=-0.7ex]\node[draw=popink,line width=1.6pt,fill=poporange,rounded corners=4pt,minimum size=22pt,inner sep=0]{\popdisplay\fontsize{12}{12}\selectfont\color{popcream}\thepopsec};%
  \hspace{8pt}{\popdisplay\fontsize{15}{17}\selectfont\color{popink}\MakeUppercase{#1}}\par
  \vspace{4pt}\noindent{\color{poporange}\rule{36pt}{3.6pt}}
  \end{minipage}\par\nopagebreak\vspace{8pt}
}
\newcounter{popsub}
\newcommand{\courssub}[1]{%
  \stepcounter{popsub}%
  \par\vspace{9pt}\noindent{\popxbold\fontsize{11.5}{13}\selectfont\color{popink}{\color{poporange}\thepopsec.\thepopsub}\hspace{6pt}#1}\par\nopagebreak\vspace{4pt}
}

% ============================================================
% Boîtes pédagogiques — même recette : fond blanc, bordure épaisse colorée,
% ombre dure encre, badge plein. Titre optionnel [#1].
% ============================================================
\tcbset{popbase/.style={breakable,enhanced,colback=white,boxrule=2.3pt,arc=9pt,
  shadow={3pt}{-3pt}{0pt}{popink},left=14pt,right=14pt,top=11pt,bottom=11pt,before skip=11pt,after skip=15pt,
  fontupper=\popsemi\fontsize{10.6}{14.5}\selectfont\color{popink},
  fontlower=\popsemi\fontsize{10.6}{14.5}\selectfont\color{popink}}}
% Coche dessinée (le glyphe ✓ n'existe pas dans les polices du document)
\newcommand{\popcheck}[1]{\tikz[baseline=-0.5ex]\draw[#1,line width=1.6pt,line cap=round,line join=round](-0.13,0.0)--(-0.04,-0.09)--(0.13,0.1);}
\providecommand{\coche}{\popcheck{popgreen}}
\providecommand{\resultat}[1]{\par\smallskip\noindent\fcolorbox{popgreen}{popgreenL}{\parbox{\dimexpr\linewidth-2\fboxsep-2\fboxrule\relax}{\textbf{Résultat.} #1}}\par\smallskip}
\newcommand{\popboxhead}[3]{\poppill{#1}{#2}{white}\ifx\relax#3\relax\else\hspace{6pt}{\popxbold\fontsize{10.6}{12}\selectfont\color{popink}#3}\fi\par\vspace{7pt}}

\newtcolorbox{aretenir}[1][]{popbase,colframe=popgreen,before upper={\popboxhead{À retenir}{popgreen}{#1}}}
% Texte en anglais (document de lecture, dialogue, extrait) : cadre
% d'étude. Un titre optionnel (source, auteur) s'affiche dans l'en-tête.
\newtcolorbox{textebox}[1][]{popbase,colframe=popblue,before upper={\popboxhead{Text}{popblue}{#1}\begin{otherlanguage}{english}},after upper={\end{otherlanguage}}}
% Marques ✓ / ✗ (forme correcte / fautive) souvent utilisées par les modèles
\providecommand{\cross}{\textcolor{poppink}{\ensuremath{\times}}}
\providecommand{\xmark}{\cross}\providecommand{\crossmark}{\cross}\providecommand{\cmark}{\ensuremath{\checkmark}}
% Ligne de réponse / justification à compléter (inventée par les modèles)
\providecommand{\justification}[1]{\par\noindent\textit{Justification~:} \dotfill\par}
% \textipa (package tipa non chargé) : la police ne couvre pas l'API -> simple passage
\providecommand{\textipa}[1]{#1}
% Soulignement (ulem non chargé) : \uline, \ul
\providecommand{\uline}[1]{\underline{#1}}\providecommand{\ul}[1]{\underline{#1}}
% \glossaire{\item ...} inventée par les modèles : liste de glossaire
\providecommand{\glossaire}[1]{\begin{itemize}#1\end{itemize}}
% \alert (beamer) inventée par les modèles : texte mis en évidence
\providecommand{\alert}[1]{\textbf{\textcolor{poppink}{#1}}}
% variantes ASCII d'ordinaux écrites en macro
\providecommand{\iere}{\textsuperscript{re}}\providecommand{\ieme}{\textsuperscript{e}}\providecommand{\ere}{\textsuperscript{re}}\providecommand{\eme}{\textsuperscript{e}}\providecommand{\er}{\textsuperscript{er}}\providecommand{\ie}{\textsuperscript{e}}
% Ordinaux français écrits en macro par les modèles : 1\ière, 3\ième
\newcommand{\ière}{\textsuperscript{re}}\newcommand{\ième}{\textsuperscript{e}}
% \exemple (inventée par les modèles) : étiquette « Exemple : »
\providecommand{\exemple}{\textbf{Exemple~:}\ }
% \square utilisable aussi hors mode math (cases à cocher)
\AtBeginDocument{\let\squareMath\square\renewcommand{\square}{\ensuremath{\squareMath}}}
% Mot ou phrase en anglais à l'intérieur d'un paragraphe français
\newcommand{\eng}[1]{\ifmmode\text{\textit{#1}}\else\begingroup\selectlanguage{english}\itshape #1\endgroup\fi}
\let\esp\eng % alias toléré
% flèches d'intonation inventées par les modèles
\providecommand{\textsearrow}{}\renewcommand{\textsearrow}{\ensuremath{\searrow}}\providecommand{\textnearrow}{}\renewcommand{\textnearrow}{\ensuremath{\nearrow}}
% \fr{..} : mot français (inventé par les modèles)
\providecommand{\fr}[1]{\textit{#1}}
\newtcolorbox{exemplebox}[1][]{popbase,colframe=poporange,before upper={\popboxhead{Exemple}{poporange}{#1}}}
\newtcolorbox{definition}[1][]{popbase,colframe=popblue,before upper={\popboxhead{Définition}{popblue}{#1}}}
\newtcolorbox{propriete}[1][]{popbase,colframe=poppurple,before upper={\popboxhead{Propriété}{poppurple}{#1}}}
\newtcolorbox{methode}[1][]{popbase,colframe=popgold,before upper={\popboxhead{Méthode}{popgold}{#1}}}
\newtcolorbox{attention}[1][]{popbase,colframe=poppink,before upper={\popboxhead{Attention}{poppink}{#1}}}
\newtcolorbox{experience}[1][]{popbase,colframe=popblue,colback=popblueL,before upper={\popboxhead{Expérience}{popblue}{#1}}}
% « Le savais-tu ? » : carte violette pleine (contexte, culture, Cameroun)
\newtcolorbox{savaistu}[1][]{popbase,colback=poppurple,colframe=popink,
  fontupper=\popsemi\fontsize{10.4}{14.2}\selectfont\color{white},
  before upper={\poppill{Le savais-tu\,?}{white}{poppurple}\ifx\relax#1\relax\else\hspace{6pt}{\popxbold\color{white}#1}\fi\par\vspace{7pt}}}
% Exercice résolu : énoncé en haut, \tcblower puis la solution sur fond crème
\newcounter{popexo}\newcounter{popexr}
\newtcolorbox{exoresolu}[1][]{popbase,colframe=poporange,colbacklower=popcream,
  segmentation style={popink,line width=1.2pt,dashed},
  before upper={\stepcounter{popexr}\popboxhead{Exercice résolu \thepopexr}{poporange}{#1}},
  before lower={\poppill{Solution}{popink}{popcream}\par\vspace{6pt}}}
% Exercice (non résolu en place) — la correction est dans la partie Corrigés
\newtcolorbox{exercice}[1][]{popbase,colframe=popink,
  before upper={\stepcounter{popexo}\popboxhead{Exercice \thepopexo}{popink}{#1}}}
% Corrigé
\newtcolorbox{corrige}[1][]{popbase,colframe=popgreen,colback=popgreenL,
  before upper={\popboxhead{Corrigé}{popgreen}{#1}}}
% Objectifs / prérequis / bilan
\newtcolorbox{objectifs}{popbase,colframe=popink,colback=poporangeL,
  before upper={\popboxhead{Objectifs de la leçon}{popink}{}}}
\newtcolorbox{prerequis}{popbase,colframe=popink,colback=popblueL,
  before upper={\popboxhead{Prérequis}{popblue}{}}}
\newtcolorbox{bilan}{popbase,colframe=popink,colback=white,boxrule=2.8pt,
  before upper={\popboxhead{Fiche bilan}{poporange}{L'essentiel à connaître pour l'examen}}}
% Figure encadrée avec légende
\newtcolorbox{popfigure}[1][]{enhanced,breakable=false,colback=white,colframe=popline,boxrule=1.4pt,arc=8pt,
  left=10pt,right=10pt,top=10pt,bottom=8pt,before skip=10pt,after skip=12pt,halign=center,
  lower separated=false,halign lower=center,
  fontlower=\popsemi\fontsize{9}{11.5}\selectfont\color{popmuted},
  #1}
\newcommand{\legende}[1]{\par\vspace{6pt}{\popsemi\fontsize{9}{11.5}\selectfont\color{popmuted}#1}}

% ============================================================
% Signature OnBuch+
% ============================================================
\newcommand{\poplogoinline}[1]{{\poplogo\fontsize{#1}{#1}\selectfont\color{popink}OnBuch\color{poporange}+}}
\newcommand{\signatureonbuch}{%
  \begin{tcolorbox}[enhanced,colback=popink,colframe=popink,boxrule=2.2pt,arc=10pt,
      drop shadow=poporange,left=14pt,right=14pt,top=10pt,bottom=10pt,before skip=0pt,after skip=0pt]
    \tikz[baseline=-0.7ex]\node[circle,fill=popgreen,draw=popcream,line width=1.3pt,minimum size=22pt,inner sep=0]{\popcheck{white}};%
    \hspace{9pt}\begin{minipage}[c]{0.62\linewidth}
      {\popxbold\fontsize{12}{13}\selectfont\color{popcream}Signé\ \textbullet\ {\poplogo OnBuch\color{poporange}+}}\par\vspace{2pt}
      {\raggedright\popsemi\fontsize{8}{10.5}\selectfont\color{popline}\MakeUppercase{Équipe pédagogique OnBuch+}\\\MakeUppercase{Conforme au programme officiel MINESEC}\par}
    \end{minipage}\hfill
    {\poplogo\fontsize{22}{22}\selectfont\color{popcream}OnBuch\color{poporange}+}
  \end{tcolorbox}%
}

% ============================================================
% Page de garde
% #1 titre · #2 matière · #3 niveau · #4 module · #5 n° leçon · #6 durée ·
% #7 description / objectifs (texte ou itemize)
% ============================================================
\newcommand{\pagegarde}[7]{%
  \thispagestyle{empty}
  \newgeometry{a4paper,margin=1.7cm,top=1.6cm,bottom=1.4cm}%
  \begin{tikzpicture}[remember picture,overlay]
    \fill[poporange!18] ([xshift=-3.2cm,yshift=-3.6cm]current page.north east) circle (4.6cm);
    \fill[poppurple!14] ([xshift=2.4cm,yshift=7.6cm]current page.south west) circle (3.8cm);
  \end{tikzpicture}%
  {\poplogo\fontsize{30}{30}\selectfont\color{popink}OnBuch\color{poporange}+}\hfill
  \raisebox{0.6ex}{\poppill{Cours signé \textbullet\ Premium}{popink}{popcream}}\par
  \vspace{1pt}\popsquiggle{poppurple}\par
  \vspace{8pt}
  \begin{tcolorbox}[enhanced,colback=poporange,colframe=popink,boxrule=2.8pt,arc=13pt,
      drop shadow=popink,left=18pt,right=18pt,top=12pt,bottom=13pt,after skip=10pt]
    \poppill{#2 \textbullet\ #3}{popink}{popcream}\hspace{5pt}\poppill{#4}{white}{popink}\par\vspace{8pt}
    {\popdisplay\fontsize{14}{14}\selectfont\color{popink}LEÇON #5}\par\vspace{4pt}
    {\raggedright\hyphenpenalty=10000\exhyphenpenalty=10000\popdisplay\fontsize{25}{27}\selectfont\color{popcream}\MakeUppercase{#1}\par}
    \vspace{8pt}
    {\popxbold\fontsize{9.5}{9.5}\selectfont\color{popink}\MakeUppercase{Durée conseillée : #6}}
  \end{tcolorbox}
  \begin{tcolorbox}[enhanced,colback=white,colframe=popink,boxrule=2.2pt,arc=10pt,
      drop shadow=popink,left=16pt,right=16pt,top=10pt,bottom=10pt,after skip=8pt]
    \poppill{Dans cette leçon}{poppurple}{white}\par\vspace{5pt}
    \popsemi\fontsize{9.3}{12.6}\selectfont\color{popink}#7
  \end{tcolorbox}
  \noindent\begin{minipage}[t]{0.487\linewidth}
    \begin{tcolorbox}[enhanced,colback=popgreen,colframe=popink,boxrule=2.2pt,arc=10pt,
        drop shadow=popink,left=11pt,right=11pt,top=10pt,bottom=10pt,height=3.1cm,valign=center]
      \tcbox[colback=white,colframe=popink,boxrule=1.3pt,arc=5pt,boxsep=0pt,nobeforeafter,
        left=3pt,right=3pt,top=3pt,bottom=3pt,valign=center]{\qrcode[height=1.9cm]{https://play.google.com/store/apps/details?id=com.luvvix.onbuch&hl=fr}}%
      \hspace{8pt}\begin{minipage}[c]{0.5\linewidth}
        {\popdisplay\fontsize{11}{12.5}\selectfont\color{white}\MakeUppercase{Télécharge OnBuch}}\par\vspace{4pt}
        {\raggedright\popsemi\fontsize{8.4}{11}\selectfont\color{white}Gratuit \textbullet\ Google Play\\Tuteur IA, annales, concours, résultats\par}
      \end{minipage}
    \end{tcolorbox}
  \end{minipage}\hfill
  \begin{minipage}[t]{0.487\linewidth}
    \begin{tcolorbox}[enhanced,colback=poppurple,colframe=popink,boxrule=2.2pt,arc=10pt,
        drop shadow=popink,left=11pt,right=11pt,top=10pt,bottom=10pt,height=3.1cm,valign=center]
      \tikz[baseline=-0.7ex]\node[circle,fill=white,draw=popink,line width=1.3pt,minimum size=1.6cm,inner sep=0]{\popdisplay\fontsize{24}{24}\selectfont\color{poppurple}+};%
      \hspace{8pt}\begin{minipage}[c]{0.6\linewidth}
        {\popdisplay\fontsize{11}{12.5}\selectfont\color{white}\MakeUppercase{Passe à OnBuch+}}\par\vspace{4pt}
        {\raggedright\popsemi\fontsize{8.4}{11}\selectfont\color{white}Tous les cours signés, Tuteur IA illimité, fiches PDF — onglet OnBuch+ de l'app\par}
      \end{minipage}
    \end{tcolorbox}
  \end{minipage}
  \vspace{8pt}
  \popcontenu
  \vfill
  \signatureonbuch
  \restoregeometry
  \newpage
}

% Rangée « Ce cours contient » (page de garde)
\newcommand{\poptile}[3]{% #1 couleur #2 gros texte #3 légende
  \begin{tcolorbox}[enhanced,colback=#1,colframe=popink,boxrule=2pt,arc=9pt,shadow={2.5pt}{-2.5pt}{0pt}{popink},
      left=6pt,right=6pt,top=9pt,bottom=9pt,halign=center,valign=center,height=1.75cm,width=\linewidth,nobeforeafter]
    {\popdisplay\fontsize{12.5}{13}\selectfont\color{white}\MakeUppercase{#2}}\par\vspace{3pt}
    {\centering\popsemi\fontsize{7.8}{9.5}\selectfont\color{white}#3\par}
  \end{tcolorbox}}
\newcommand{\popcontenu}{%
  \par\noindent{\popxboldls\fontsize{7.5}{7.5}\selectfont\color{popmuted}CE COURS SIGNÉ CONTIENT}\par\vspace{7pt}
  \noindent\begin{minipage}[t]{0.235\linewidth}\poptile{popblue}{Cours}{complet, illustré, pas à pas}\end{minipage}\hfill
  \begin{minipage}[t]{0.235\linewidth}\poptile{poporange}{Méthodes}{et exercices résolus}\end{minipage}\hfill
  \begin{minipage}[t]{0.235\linewidth}\poptile{poppink}{Exercices}{type examen + corrigés}\end{minipage}\hfill
  \begin{minipage}[t]{0.235\linewidth}\poptile{popgreen}{Fiche bilan}{l'essentiel à retenir}\end{minipage}\par}

% Sommaire
\newtcolorbox{sommaire}{enhanced,colback=white,colframe=popink,boxrule=2.2pt,arc=10pt,
  drop shadow=popink,left=18pt,right=18pt,top=13pt,bottom=13pt,before skip=6pt,after skip=20pt,
  before upper={\poppill{Sommaire}{poppurple}{white}\par\vspace{9pt}\popsemi\fontsize{10.6}{14.5}\selectfont\color{popink}}}

% Signature de fin de cours — poussée tout en bas de la dernière page
\newcommand{\findecours}{%
  \par\vfill\nopagebreak
  \begin{center}\popsquiggle{poporange}\end{center}\vspace{4pt}
  \signatureonbuch
}
\AtBeginDocument{\def\llbracket{\mathopen{[\mkern-3.5mu[}}\def\rrbracket{\mathclose{]\mkern-3.5mu]}}} % intervalles d'entiers (glyphe ⟦ absent de la police)
% Glyphes absents de Fira Math : redéfinis à partir de glyphes disponibles
\AtBeginDocument{%
  \def\setminus{\mathbin{\backslash}}\let\smallsetminus\setminus
  \def\vdots{\mathord{\vbox{\baselineskip3.2pt\lineskiplimit0pt\kern4pt\hbox{.}\hbox{.}\hbox{.}}}}%
  \def\ddots{\mathinner{\mkern1mu\raise6.5pt\hbox{.}\mkern2mu\raise3.6pt\hbox{.}\mkern2mu\raise0.7pt\hbox{.}\mkern1mu}}%
  \def\triangle{\mathord{\scalebox{0.9}{$\symup{\Delta}$}}}%
  \def\nabla{\mathord{\raisebox{\depth}{\scalebox{1}[-1]{$\symup{\Delta}$}}}}%
  \def\checkmark{\popcheck{popdark}}%
  \def\star{\mathbin{\ast}}\def\bigstar{\mathord{\ast}}%
  \def\diamond{\mathbin{\scalebox{0.6}{$\Diamond$}}}%
  \def\bigcup{\mathop{\mathchoice{\raisebox{-0.3em}{\scalebox{1.7}{$\cup$}}}{\scalebox{1.25}{$\cup$}}{\cup}{\cup}}\displaylimits}%
  \def\bigcap{\mathop{\mathchoice{\raisebox{-0.3em}{\scalebox{1.7}{$\cap$}}}{\scalebox{1.25}{$\cap$}}{\cap}{\cap}}\displaylimits}%
  \def\lesseqqgtr{\mathrel{\lessgtr}}\def\bigoplus{\mathop{\oplus}\displaylimits}%
}
\providecommand{\ssi}{si et seulement si} % abréviation fréquente
\AtBeginDocument{\providecommand{\wideparen}{\overparen}} % arc de cercle

\begin{document}

\pagegarde{Lesson 46 — Grammar: Review phrasal verbs in simple expressions of purpose}{Anglais}{Tle SH}{Chapter 4 : Citizenship and Human Rights}{EN46}{2 h}{Cette leçon de grammaire révise les phrasal verbs les plus utiles et montre comment les employer dans des expressions simples de but (to, in order to, so as to + infinitif). Ancrée dans le chapitre « Citizenship and Human Rights », elle prépare les élèves à exprimer leurs motivations et objectifs en anglais correct et naturel.
\vspace{4pt}
\begin{multicols}{2}\small
\begin{itemize}
\item À la fin de cette leçon, je sais définir ce qu'est un phrasal verb et reconnaître ses deux particules (adverbe et préposition)
\item Je sais distinguer les phrasal verbs séparables et inséparables et placer correctement le complément
\item Je sais employer les phrasal verbs courants liés à la citoyenneté et aux droits (stand up for, give up, carry out, look after, put forward...)
\item Je sais exprimer le but avec to + infinitif, in order to et so as to
\item Je sais exprimer le but négatif avec in order not to / so as not to
\item Je sais combiner un phrasal verb et une expression de but dans une même phrase
\end{itemize}\end{multicols}}

\begin{sommaire}
\begin{multicols}{2}
\begin{enumerate}
\item Observation : découvrir les phrasal verbs en contexte
\item Règle 1 : formation et types de phrasal verbs
\item Règle 2 : les phrasal verbs de la citoyenneté et des droits
\item Règle 3 : exprimer le but (expressions of purpose)
\item Comparaison français/anglais et pièges des francophones
\item Entraînement guidé : de la phrase au paragraphe
\item Activité d'intégration
\item Méthodes et exercices résolus
\item Exercices
\item Corrigés des exercices
\item Fiche bilan
\end{enumerate}
\end{multicols}
\end{sommaire}

\begin{objectifs}
\begin{itemize}
\item À la fin de cette leçon, je sais définir ce qu'est un phrasal verb et reconnaître ses deux particules (adverbe et préposition)
\item Je sais distinguer les phrasal verbs séparables et inséparables et placer correctement le complément
\item Je sais employer les phrasal verbs courants liés à la citoyenneté et aux droits (stand up for, give up, carry out, look after, put forward...)
\item Je sais exprimer le but avec to + infinitif, in order to et so as to
\item Je sais exprimer le but négatif avec in order not to / so as not to
\item Je sais combiner un phrasal verb et une expression de but dans une même phrase
\item Je sais éviter les erreurs typiques des francophones (traduction mot à mot, mauvais placement du pronom, oubli de la particule)
\item Je sais utiliser ces structures dans des phrases et un court paragraphe sur l'engagement citoyen
\end{itemize}
\end{objectifs}

\begin{prerequis}
\begin{itemize}
\item L'infinitif avec to après certains verbes (want to, decide to) — vu en Première
\item Les verbes de base et leurs compléments (verbe transitif / intransitif)
\item Les pronoms compléments anglais (it, them, him, her)
\item Le vocabulaire de base du chapitre : rights, duties, citizen, vote, law
\item La négation simple et la structure de la phrase anglaise (SVO)
\end{itemize}
\end{prerequis}

\newpage

\coursec{Observation : découvrir les phrasal verbs en contexte}

Imagine a young activist in Bamenda who wants to \eng{stand up for} children's rights. She says: “I joined the association \eng{to speak out against} injustice. I will never \eng{give up}, because I want \eng{to bring about} change in my community.” Without noticing it, she uses \cle{phrasal verbs} — those small \emph{verb + particle} combinations that English speakers love — to explain \textbf{why} she acts. Today, you will learn to do the same: express your purposes naturally, like a true citizen of the English-speaking world.

\textbf{Problématique :} comment l'anglais exprime-t-il des idées que le français rend par un seul verbe savant (\emph{abandonner}, \emph{défendre}, \emph{réaliser}, \emph{protéger}) ? Et comment dit-on clairement \emph{pourquoi} l'on agit ? Pour le découvrir, observons d'abord un texte authentique avant d'énoncer la moindre règle.

\courssub{Lecture : Amina, une volontaire engagée à Bamenda}

\begin{textebox}[A Young Volunteer in Bamenda]
Amina is a young volunteer in Bamenda, in the North-West Region of Cameroon. She is nineteen years old, and she studies law at the university. Two years ago, she joined an NGO \textbf{to stand up for} children's rights. Every Saturday, she \textbf{carries out} door-to-door visits in the neighbourhoods \textbf{in order to look after} street children. She talks to their families, brings them food and school materials, and listens to their problems.

Last year, Amina noticed that many children could not go to school because their parents were too poor. She decided to act. She organised meetings with local chiefs and business people \textbf{to raise money} for school fees. She also spoke on the community radio \textbf{so as to raise awareness} about child labour. Her message was simple: “Every child deserves a classroom, not a workshop.”

The work is not easy. Some people laugh at her, and some doors stay closed. But Amina wants \textbf{to bring about} real change, so as to give these children a better future. “When I feel tired,” she says, “I remember the smile of a child who goes back to school. That is why I will never \textbf{give up}!”

Thanks to her action, more than fifty children have returned to school this year. Amina believes that young citizens can transform their community — if they refuse to give up and if they stand up for what is right.
\end{textebox}

\textbf{Glossaire :}

\begin{center}
\begin{tabularx}{\linewidth}{|l|l|X|}
\hline
\rowcolor{popblueL} Mot anglais & Nature & Traduction française \\
\hline
volunteer & nom & volontaire, bénévole \\
\hline
NGO (non-governmental organisation) & nom & ONG \\
\hline
door-to-door visits & nom & visites à domicile (porte-à-porte) \\
\hline
street children & nom & enfants des rues \\
\hline
school fees & nom & frais de scolarité \\
\hline
to raise awareness & locution verbale & sensibiliser (litt. « élever la conscience ») \\
\hline
child labour & nom & travail des enfants \\
\hline
to deserve & verbe & mériter \\
\hline
thanks to & locution & grâce à \\
\hline
\end{tabularx}
\end{center}

\textbf{Questions de compréhension (reading) :}
\begin{enumerate}
\item Who is Amina and where does she live?
\item Why did she join an NGO?
\item What does she do every Saturday, and for what purpose?
\item How did she use the community radio, and why?
\item What sentence shows that Amina refuses to stop her action?
\end{enumerate}

\begin{corrige}[Questions de compréhension]
\begin{enumerate}
\item Amina is a nineteen-year-old law student and volunteer; she lives in Bamenda, in the North-West Region of Cameroon.
\item She joined an NGO to stand up for children's rights.
\item Every Saturday, she carries out door-to-door visits in order to look after street children.
\item She spoke on the community radio so as to raise awareness about child labour.
\item “I will never give up!” shows her determination.
\end{enumerate}
\end{corrige}

\courssub{Repérage guidé : les verbes à particule du texte}

Relisons le texte et soulignons les groupes verbaux en gras. On constate qu'ils sont tous formés de \textbf{deux ou trois mots} : un verbe suivi d'une ou deux particules (\emph{up, out, for, after, about}).

\begin{center}
\begin{tabularx}{\linewidth}{|l|X|X|}
\hline
\rowcolor{popblueL} Phrasal verb & Exemple du texte & Traduction \\
\hline
stand up for & She joined an NGO to stand up for children's rights. & défendre (les droits de l'enfant) \\
\hline
carry out & She carries out door-to-door visits. & effectuer, mener (des visites) \\
\hline
look after & \ldots in order to look after street children. & s'occuper de, prendre soin de \\
\hline
bring about & She wants to bring about real change. & provoquer, entraîner (un changement) \\
\hline
give up & I will never give up! & abandonner, renoncer \\
\hline
\end{tabularx}
\end{center}

\begin{definition}[Phrasal verb — verbe à particule]
A \cle{phrasal verb} is a verb combined with a \cle{particle} (an adverb or a preposition: \emph{up, out, for, after, about, on, off\ldots}) that \textbf{changes or completes its meaning}. Example: \eng{give} + \eng{up} = \eng{give up} = “stop trying” (abandonner). Le sens du phrasal verb est souvent \textbf{différent} du sens du verbe seul : il faut l'apprendre comme un mot de vocabulaire à part entière.
\end{definition}

\begin{exemplebox}[Le verbe seul ≠ le phrasal verb]
\eng{She will never give up.} « Elle n'abandonnera jamais. » (\eng{give} seul = « donner » !)

\eng{They stand up for their rights.} « Ils défendent leurs droits. » (\eng{stand} seul = « se tenir debout ».)

\eng{She looks after the children.} « Elle s'occupe des enfants. » (\eng{look} seul = « regarder ».)

\eng{They carried out a survey.} « Ils ont mené une enquête. » (\eng{carry} seul = « porter ».)
\end{exemplebox}

\begin{attention}[Piège : ne traduis pas mot à mot !]
Le francophone traduit souvent la particule séparément : \eng{give up} n'a rien à voir avec « donner vers le haut ». La particule \textbf{change le sens du verbe}. Retiens chaque phrasal verb avec sa traduction globale, comme un seul mot : \eng{give up} = \emph{abandonner} ; \eng{look after} = \emph{s'occuper de} ; \eng{bring about} = \emph{provoquer}.
\end{attention}

\courssub{Repérage guidé : les expressions de but dans le texte}

Le texte ne dit pas seulement \emph{ce que} fait Amina ; il explique \emph{pourquoi} elle agit. Observons les structures qui expriment le \cle{but} (purpose) :

\begin{center}
\begin{tabularx}{\linewidth}{|l|X|X|}
\hline
\rowcolor{popgreenL} Expression de but & Exemple du texte & Traduction \\
\hline
to + verbe & She joined an NGO \textbf{to stand up for} children's rights. & pour défendre \\
\hline
in order to + verbe & \ldots visits \textbf{in order to look after} street children. & afin de s'occuper de \\
\hline
so as to + verbe & She spoke on the radio \textbf{so as to raise awareness}. & de façon à sensibiliser \\
\hline
\end{tabularx}
\end{center}

\begin{aretenir}[Premier constat]
\begin{itemize}
\item \eng{to + base verbale} est la façon la plus simple et la plus fréquente d'exprimer le but : \eng{She studies law to defend children.} « Elle étudie le droit pour défendre les enfants. »
\item \eng{in order to} et \eng{so as to} sont plus explicites et légèrement plus formels ; ils insistent sur l'intention.
\item Les trois structures sont suivies de la \textbf{base verbale} (l'infinitif sans \eng{to}), jamais d'un verbe conjugué ni d'un \emph{-ing}.
\end{itemize}
\end{aretenir}

\begin{attention}[Erreur fréquente des francophones]
Ne dis jamais \eng{for to protect the children} ni \eng{for protect}. En anglais, le « pour » de but devant un verbe ne se traduit pas par \eng{for} : on utilise \eng{to}. \eng{She came to help us.} « Elle est venue pour nous aider. » (\eng{for} + nom seulement : \eng{She came for the meeting.})
\end{attention}

\coursec{Règle 1 : Formation et types de phrasal verbs}

Au-delà du sens, il faut maîtriser la \textbf{grammaire} des phrasal verbs pour les utiliser correctement à l'écrit comme à l'oral (Bac : expression écrite, traduction). On distingue quatre types selon leur transitivité et la place de l'objet.

\begin{propriete}[Type 1 — Intransitifs (sans objet)]
Ces phrasal verbs n'acceptent pas d'objet (COI ou COD). La particule est un adverbe.
\begin{itemize}
\item \eng{give up} (abandonner) — \eng{She never gives up.}
\item \eng{show up} (arriver, se présenter) — \eng{He didn't show up.}
\item \eng{grow up} (grandir) — \eng{They grew up in Douala.}
\item \eng{come back} (revenir) — \eng{When will you come back?}
\end{itemize}
\emph{Règle :} Pas d'objet possible, donc pas de problème de placement.
\end{propriete}

\begin{propriete}[Type 2 — Transitifs séparables (verbe + adverbe)]
L'objet (nom) peut se placer \textbf{après la particule} OU \textbf{entre le verbe et la particule}.
\begin{itemize}
\item \eng{carry out} (effectuer) — \eng{carry out a survey} / \eng{carry a survey out}
\item \eng{give up} (arrêter qqch) — \eng{give up smoking} / \eng{give smoking up}
\item \eng{put off} (reporter) — \eng{put off the meeting} / \eng{put the meeting off}
\item \eng{find out} (découvrir) — \eng{find out the truth} / \eng{find the truth out}
\end{itemize}
\emph{Règle d'or :} Si l'objet est un \textbf{pronom} (\eng{it, them, me, us\ldots}), il \textbf{doit} s'intercaler : \eng{carry it out} (jamais \eng{carry out it}).
\end{propriete}

\begin{propriete}[Type 3 — Transitifs inséparables (verbes prépositionnels)]
Le verbe est suivi d'une préposition ; l'objet (nom ou pronom) se place \textbf{toujours après la préposition}.
\begin{itemize}
\item \eng{look after} (s'occuper de) — \eng{look after the children} / \eng{look after them}
\item \eng{stand up for} (défendre) — \eng{stand up for rights} / \eng{stand up for them}
\item \eng{look for} (chercher) — \eng{look for a job} / \eng{look for it}
\item \eng{go through} (traverser, subir) — \eng{go through a crisis} / \eng{go through it}
\end{itemize}
\emph{Astuce :} La particule est une préposition (elle introduit l'objet) ; on ne peut pas la séparer du verbe.
\end{propriete}

\begin{propriete}[Type 4 — Verbes à trois mots (verbe + adverbe + préposition)]
Ils sont toujours transitifs et \textbf{inséparables}. L'objet suit la deuxième particule.
\begin{itemize}
\item \eng{stand up for} (défendre) — \eng{stand up for them}
\item \eng{put up with} (tolérer, supporter) — \eng{put up with the noise} / \eng{put up with it}
\item \eng{look forward to} (avoir hâte de) — \eng{look forward to it} (\eng{to} est préposition $\rightarrow$ \eng{-ing} : \eng{look forward to hearing from you})
\item \eng{get away with} (s'en tirer avec) — \eng{get away with it}
\end{itemize}
\end{propriete}

\begin{exemplebox}[Test de séparation : nom vs pronom]
\begin{tabular}{lll}
\hline
\rowcolor{popblueL} Phrase correcte & Phrase incorrecte & Type \\
\hline
\eng{She carried out the plan.} & — & Séparable (nom) \\
\eng{She carried it out.} & \eng{*She carried out it.} & Séparable (pronom) \\
\eng{She looks after them.} & \eng{*She looks them after.} & Inséparable \\
\eng{She stands up for them.} & \eng{*She stands them up for.} & 3 mots (inséparable) \\
\hline
\end{tabular}
\end{exemplebox}

\begin{exercice}[Identification du type]
Classe ces phrasal verbs du texte (et du thème citoyenneté) dans le bon type : \eng{stand up for, carry out, look after, bring about, give up, speak out against, sign up for, join in, reach out to}.
\begin{enumerate}
\item Intransitifs :
\item Transitifs séparables :
\item Transitifs inséparables (2 mots) :
\item Verbes à 3 mots (inséparables) :
\end{enumerate}
\end{exercice}

\begin{corrige}[Exercice « Identification du type »]
\begin{enumerate}
\item Intransitifs : \eng{give up} (dans le sens « abandonner » sans objet), \eng{join in} (participer).
\item Transitifs séparables : \eng{carry out}, \eng{bring about}, \eng{find out}, \eng{give up} (qqch, ex: \eng{give up smoking}).
\item Transitifs inséparables (2 mots) : \eng{look after}, \eng{reach out to}.
\item Verbes à 3 mots (inséparables) : \eng{stand up for}, \eng{speak out against}, \eng{sign up for}, \eng{put up with}.
\end{enumerate}
\emph{Note :} \eng{give up} est ambivalent : intransitif (\eng{She gave up}) ou transitif séparable (\eng{She gave up smoking / gave it up}).
\end{corrige}

\coursec{Règle 2 : Les phrasal verbs de la citoyenneté et des droits}

Le thème « Citizenship and Human Rights » impose un lexique précis. Voici les phrasal verbs incontournables pour le Bac, classés par fonction communicative.

\begin{definition}[Lexique citoyen — Phrasal verbs essentiels]
\begin{center}
\begin{tabularx}{\linewidth}{|l|l|X|X|}
\hline
\rowcolor{popblueL} Phrasal verb & Type & Sens français & Exemple contextuel \\
\hline
stand up for & 3 mots (insép.) & défendre (une cause, qqn) & \eng{Citizens must stand up for their rights.} \\
\hline
speak out against & 3 mots (insép.) & dénoncer, s'élever contre & \eng{She spoke out against corruption.} \\
\hline
look after & 2 mots (insép.) & s'occuper de, prendre soin de & \eng{NGOs look after vulnerable children.} \\
\hline
carry out & Séparable & mener, effectuer (action) & \eng{They carried out a survey.} \\
\hline
bring about & Séparable & provoquer, entraîner (changement) & \eng{The law brought about change.} \\
\hline
give up & Intrans. / Sépar. & abandonner (effort / qqch) & \eng{Don't give up! / Give up bad habits.} \\
\hline
find out & Séparable & découvrir, s'informer sur & \eng{Find out your rights.} \\
\hline
sign up for & 3 mots (insép.) & s'inscrire, adhérer à & \eng{He signed up for the volunteer program.} \\
\hline
join in & Intransitif & participer à & \eng{Everyone joined in the cleanup.} \\
\hline
reach out to & 2 mots (insép.) & tendre la main à, contacter & \eng{We reached out to the community.} \\
\hline
stand up to & 2 mots (insép.) & résister à, faire face à (intimidation) & \eng{She stood up to the bully.} \\
\hline
put up with & 3 mots (insép.) & tolérer, supporter (injustice) & \eng{We won't put up with discrimination.} \\
\hline
\end{tabularx}
\end{center}
\end{definition}

\begin{exemplebox}[Contexte camerounais]
\begin{itemize}
\item \eng{Students in Buea \textbf{signed up for} the election observation mission.} (Les étudiants de Buéa se sont inscrits pour la mission d'observation électorale.)
\item \eng{The association \textbf{reaches out to} displaced persons in the Far North.} (L'association tend la main aux déplacés de l'Extrême-Nord.)
\item \eng{We must not \textbf{put up with} illegal taxation at roadblocks.} (Nous ne devons pas tolérer la taxation illégale aux barrages.)
\item \eng{Young leaders \textbf{spoke out against} gender-based violence.} (De jeunes leaders ont dénoncé les violences basées sur le genre.)
\end{itemize}
\end{exemplebox}

\coursec{Règle 3 : Exprimer le but (Expressions of purpose)}

Les phrasal verbs s'intègrent naturellement dans les structures d'infinitif de but. C'est le cœur de la compétence « Grammar » de cette leçon.

\begin{propriete}[Les trois structures affirmatives]
\begin{center}
\begin{tabularx}{\linewidth}{|l|X|X|l|}
\hline
\rowcolor{popgreenL} Structure & Niveau de langue & Emploi & Exemple \\
\hline
\eng{to} + BV & Neutre / Courant & But général, le plus fréquent & \eng{She works hard \textbf{to help} others.} \\
\hline
\eng{in order to} + BV & Formel / Écrit & Insiste sur l'intention, la démarche & \eng{They met \textbf{in order to discuss} the project.} \\
\hline
\eng{so as to} + BV & Formel / Écrit & Souvent en début de phrase ou pour insister & \eng{\textbf{So as to} avoid conflict, they negotiated.} \\
\hline
\end{tabularx}
\end{center}
\cle{BV} = \cle{Base Verbale} (infinitif sans \eng{to}). Jamais de \emph{-ing}, jamais de conjugaison.
\end{propriete}

\begin{propriete}[La forme négative : \eng{in order not to} / \eng{so as not to}]
Pour exprimer un but négatif (« pour ne pas »), \textbf{on n'utilise pas \eng{not to} seul} en début de proposition de but (trop ambigu). On utilise obligatoirement :
\begin{itemize}
\item \eng{in order not to} + BV : \eng{She spoke softly \textbf{in order not to wake} the baby.}
\item \eng{so as not to} + BV : \eng{He left early \textbf{so as not to miss} the bus.}
\end{itemize}
\emph{Attention :} \eng{not to} seul est possible \emph{après} certains verbes (\eng{decide not to, choose not to, promise not to}) mais pas pour exprimer le but d'une action principale.
\end{propriete}

\begin{exemplebox}[But affirmatif vs négatif]
\eng{Amina uses the radio \textbf{to raise awareness}.} (Pour sensibiliser.)

\eng{Amina uses the radio \textbf{in order not to miss} any village.} (Pour ne manquer aucun village.)

\eng{\textbf{So as to} protect children, the law was passed.} (Début de phrase, formel.)
\end{exemplebox}

\begin{attention}[Pièges classiques des francophones]
\begin{itemize}
\item \textbf{\eng{For to} / \eng{For + -ing} :} \eng{*She came for to help.} $\rightarrow$ \eng{She came \textbf{to} help.} \eng{*This tool is for to cut.} $\rightarrow$ \eng{This tool is \textbf{for cutting} / \textbf{to cut}.} (\eng{for} + V-ing = fonction de l'objet).
\item \textbf{Oubli de la base verbale :} \eng{*She studies to \textbf{becomes} a lawyer.} $\rightarrow$ \eng{to \textbf{become}.}
\item \textbf{Confusion \eng{so that} :} \eng{so that} + \textbf{sujet + modal} (\eng{can, could, will, would}). \eng{She spoke loudly \textbf{so that everyone could hear}.} (Pas \eng{so that to hear}).
\end{itemize}
\end{attention}

\courssub{Comparaison français/anglais et pièges des francophones}

\begin{savaistu}[Phrasal verbs et verbes latins : une question de registre]
Les anglophones utilisent les \cle{phrasal verbs} à l'oral et dans les textes courants (presse, dialogues, chansons), tandis que les verbes d'origine latine (\eng{abandon, protect, accomplish, defend, observe}) — hérités du français normand après 1066 ! — sonnent plus formels et appartiennent plutôt à l'écrit administratif ou académique. Conséquence utile pour le Bac : employer correctement un phrasal verb dans ta composition montre que tu maîtrises l'anglais \emph{naturel} — c'est valorisé par les correcteurs ! Exemple : \eng{The government carried out reforms to bring about change.} sonne plus idiomatique que \eng{The government effected reforms to effectuate change.}
\end{savaistu}

\begin{center}
\begin{tabularx}{\linewidth}{|X|X|X|}
\hline
\rowcolor{popblueL} Français (verbe savant) & English (courant / phrasal) & English (formel / latin) \\
\hline
abandonner & give up & abandon, renounce \\
\hline
défendre (une cause) & stand up for & defend \\
\hline
s'occuper de & look after & care for \\
\hline
effectuer, mener & carry out & accomplish, perform \\
\hline
provoquer (un changement) & bring about & cause, generate \\
\hline
découvrir & find out & discover \\
\hline
continuer & carry on / go on & continue \\
\hline
tolérer & put up with & tolerate \\
\hline
dénoncer & speak out against & denounce \\
\hline
résister à & stand up to & resist \\
\hline
\end{tabularx}
\end{center}

\begin{attention}[Faux amis et erreurs de traduction ciblées]
\begin{itemize}
\item \eng{Look after} $\neq$ \emph{regarder après}. $\rightarrow$ \emph{S'occuper de}. (\eng{Look for} = chercher).
\item \eng{Bring about} $\neq$ \emph{apporter environ}. $\rightarrow$ \emph{Provoquer, entraîner} (un changement, une réaction). Ne pas confondre avec \eng{bring up} (élever un enfant / vomir / mentionner un sujet).
\item \eng{Stand up for} $\neq$ \emph{se lever pour}. $\rightarrow$ \emph{Défendre activement}. \eng{Stand up to} = \emph{faire face à, résister à} (une personne, un tyran).
\item \eng{Carry out} $\neq$ \emph{porter dehors}. $\rightarrow$ \emph{Exécuter, mener à bien} (une tâche, un ordre, une enquête).
\item \eng{Give up} $\neq$ \emph{donner en haut}. $\rightarrow$ \emph{Abandonner}. \eng{Give in} = \emph{céder} (sous la pression).
\end{itemize}
\end{attention}

\courssub{Carte mentale : synthèse de la leçon}

\begin{popfigure}
\begin{center}
\begin{tikzpicture}[mindmap, grow cyclic, every node/.style={concept, align=center, font=\small},
root concept/.append style={concept color=popblue, font=\bfseries},
level 1/.append style={level distance=3.4cm, sibling angle=90},
level 2/.append style={level distance=2.4cm, sibling angle=45}]
\node[root concept, align=center]{PHRASAL VERBS \\ \& PURPOSE}
  child[concept color=poporange]{ node {Formation}
      child { node[concept color=poporangeL, align=center]{Intransitifs\\ (give up) } }
      child { node[concept color=poporangeL, align=center]{Séparables\\ (carry out) } }
      child { node[concept color=poporangeL, align=center]{Inséparables 2 mots\\ (look after) } }
      child { node[concept color=poporangeL, align=center]{3 mots inséparables\\ (stand up for) } }
  }
  child[concept color=popgreen]{ node {Citoyenneté}
      child { node[concept color=popgreenL] {stand up for / speak out against} }
      child { node[concept color=popgreenL] {look after / reach out to} }
      child { node[concept color=popgreenL] {carry out / bring about} }
      child { node[concept color=popgreenL] {sign up for / join in} }
  }
  child[concept color=poppurple]{ node {Expressions de but}
      child { node[concept color=poppurpleL] {to / in order to / so as to + BV} }
      child { node[concept color=poppurpleL] {Négatif: in order not to / so as not to} }
      child { node[concept color=poppurpleL] {so that + sujet + modal} }
  }
  child[concept color=poppink]{ node {Pièges francophones}
      child { node[concept color=poppinkL] {for to / for + V-ing (but)} }
      child { node[concept color=poppinkL] {Pronom après particule séparable} }
      child { node[concept color=poppinkL] {Traduction mot à mot} }
  };
\end{tikzpicture}
\end{center}
\legende{Synthèse visuelle : 4 types grammaticaux, lexique citoyen, 3 structures de but (+ forme négative), 3 pièges majeurs.}
\end{popfigure}

\coursec{Entraînement guidé : de la phrase au paragraphe}

\courssub{Exercice résolu : repérer, analyser, transformer}

\begin{exoresolu}[Analyse grammaticale complète]
\textbf{Énoncé.} Pour la phrase ci-dessous : 
1. Souligne le(s) phrasal verb(s) et indique son (leur) type. 
2. Entoure l'expression de but. 
3. Traduis la phrase. 
4. Reformule avec \eng{in order to} puis \eng{so as to}.

\eng{Marius joined the human rights club to find out more about the law. He wants to bring about change in his village, so he carries out interviews in order to speak up for young people. He will never give up.}

\tcblower
\textbf{Solution détaillée.}
\begin{enumerate}
\item \textbf{Phrasal verbs et types :}
\begin{itemize}
\item \eng{find out} (transitif séparable) — objet \eng{more} (ou \eng{more about the law}) placé après.
\item \eng{bring about} (transitif séparable) — objet \eng{change}.
\item \eng{carries out} (transitif séparable) — objet \eng{interviews}.
\item \eng{speak up for} (3 mots, inséparable) — objet \eng{young people}.
\item \eng{give up} (intransitif ici) — pas d'objet.
\end{itemize}
\item \textbf{Expressions de but :} \eng{to find out} (but simple) ; \eng{in order to speak up for} (but formel, insistance sur la démarche).
\item \textbf{Traduction :} « Marius a rejoint le club des droits de l'homme pour en savoir plus sur la loi. Il veut provoquer le changement dans son village, alors il mène des interviews afin de défendre les jeunes. Il n'abandonnera jamais. »
\item \textbf{Reformulations :}
\begin{itemize}
\item \eng{Marius joined the human rights club \textbf{in order to find out} more about the law.}
\item \eng{Marius joined the human rights club \textbf{so as to find out} more about the law.}
\item \eng{He carries out interviews \textbf{so as to speak up for} young people.}
\end{itemize}
\end{enumerate}
\end{exoresolu}

\courssub{Exercices d'application (niveau Bac)}

\begin{exercice}[Transformation : phrasal verb + but]
Remplace le verbe savant entre parenthèses par le phrasal verb approprié (forme correcte) et relie les deux propositions par une expression de but (\eng{to / in order to / so as to}).
\begin{enumerate}
\item The NGO \_\_\_\_\_\_\_\_ (investigates) the case \_\_\_\_\_\_\_\_ \_\_\_\_\_\_\_\_ (discover) the truth. (\eng{carry out / find out})
\item We must \_\_\_\_\_\_\_\_ (defend) minority rights \_\_\_\_\_\_\_\_ \_\_\_\_\_\_\_\_ (ensure) equality. (\eng{stand up for})
\item She \_\_\_\_\_\_\_\_ (organised) a march \_\_\_\_\_\_\_\_ \_\_\_\_\_\_\_\_ (provoke) a reaction from the mayor. (\eng{carry out / bring about})
\item Volunteers \_\_\_\_\_\_\_\_ (take care of) refugees \_\_\_\_\_\_\_\_ \_\_\_\_\_\_\_\_ (help) them rebuild lives. (\eng{look after})
\item Don't \_\_\_\_\_\_\_\_ (tolerate) injustice ; \_\_\_\_\_\_\_\_ (resist) it! (\eng{put up with / stand up to})
\end{enumerate}
\end{exercice}

\begin{exercice}[Traduction thématique : Citoyenneté]
Traduis en anglais en utilisant obligatoirement un phrasal verb de la liste et une expression de but.
\begin{enumerate}
\item Les jeunes se sont inscrits pour participer au nettoyage du quartier. (\eng{sign up for / join in})
\item L'association a mené une enquête afin de découvrir les besoins des orphelins. (\eng{carry out / find out})
\item Nous devons dénoncer la corruption pour provoquer un vrai changement. (\eng{speak out against / bring about})
\item Elle a tendu la main aux victimes pour ne pas les abandonner. (\eng{reach out to / give up})
\item Le gouvernement a voté cette loi pour protéger les mineurs. (\eng{bring about / stand up for} — attention au sens)
\end{enumerate}
\end{exercice}

\begin{exercice}[Production écrite : Paragraphe argumenté]
\textbf{Consigne :} Écris un paragraphe de 80--100 mots sur le thème : \eng{``Young Cameroonians can bring about change in their communities.''}
\textbf{Contraintes :}
\begin{itemize}
\item Utilise au moins \textbf{4 phrasal verbs} différents de la leçon (surligne-les).
\item Utilise au moins \textbf{2 expressions de but} différentes (\eng{to}, \eng{in order to}, \eng{so as to}).
\item Respecte la structure : Thesis statement -- Arguments/Exemples -- Conclusion.
\item Voculaire suggéré : \eng{stand up for, speak out against, carry out, sign up for, look after, reach out to, give up, bring about}.
\end{itemize}
\end{exercice}

\begin{corrige}[Exercice « Transformation »]
\begin{enumerate}
\item The NGO \textbf{carries out} the investigation \textbf{in order to find out} the truth. (ou \eng{to find out})
\item We must \textbf{stand up for} minority rights \textbf{to ensure} equality. (ou \eng{in order to ensure})
\item She \textbf{carried out} a march \textbf{so as to bring about} a reaction from the mayor.
\item Volunteers \textbf{look after} refugees \textbf{to help} them rebuild \textbf{their} lives. (\eng{their} lives, pas \eng{the} lives).
\item Don't \textbf{put up with} injustice; \textbf{stand up to} it! (Impératif négatif \eng{Don't put up with} ; \eng{stand up to it} inséparable pronom).
\end{enumerate}
\end{corrige}

\begin{corrige}[Exercice « Traduction thématique »]
\begin{enumerate}
\item Young people \textbf{signed up for} the cleanup \textbf{to join in} the neighbourhood work. (ou \eng{in order to join in})
\item The association \textbf{carried out} a survey \textbf{in order to find out} the orphans' needs.
\item We must \textbf{speak out against} corruption \textbf{so as to bring about} real change.
\item She \textbf{reached out to} the victims \textbf{so as not to give up} on them. (But négatif : \eng{so as not to}).
\item The government passed this law \textbf{to stand up for} minors. (ou \eng{in order to protect} — \eng{stand up for} implique une action active de défense, \eng{protect} est le verbe savant ; ici phrasal verb exigé : \eng{stand up for}).
\end{enumerate}
\end{corrige}

\begin{corrige}[Exercice « Production écrite » — Proposition de modèle]
\eng{Young Cameroonians can \textbf{bring about} significant change in their communities if they decide to act. First, they must \textbf{sign up for} local associations \textbf{in order to} unite their forces. By doing so, they can \textbf{carry out} concrete projects like cleaning neighbourhoods or tutoring street children. Moreover, they should \textbf{speak out against} corruption and \textbf{stand up for} human rights \textbf{so as to} build a fairer society. Of course, the path is hard, but they must never \textbf{give up}. Every small action counts to \textbf{bring about} the Cameroon we dream of.}
\textbf{Analyse :} 6 phrasal verbs (\eng{bring about x2, sign up for, carry out, speak out against, stand up for, give up}) ; 3 expressions de but (\eng{in order to, so as to, to}). Structure claire. Registre formel mais naturel.
\end{corrige}

\begin{aretenir}[Bilan final de la leçon 46]
\begin{itemize}
\item \textbf{Forme :} Un phrasal verb = Verbe + Particule(s). 4 types grammaticaux (Intransitif, Séparable, Inséparable 2 mots, 3 mots). \textbf{Règle du pronom :} objet pronom $\rightarrow$ intercalé pour les séparables (\eng{give it up}), après pour les inséparables (\eng{look after it}).
\item \textbf{Sens :} Sens idiomatique $\neq$ somme des mots. Apprendre par blocs lexicaux (thème Citoyenneté : \eng{stand up for, speak out against, bring about\ldots}).
\item \textbf{But :} \eng{to / in order to / so as to} + \textbf{Base Verbale}. Négatif : \eng{in order not to / so as not to}. \eng{so that} + sujet + modal.
\item \textbf{Registre :} Phrasal verbs = anglais naturel, courant, valorisé à l'oral et à l'écrit (composition). Verbes latins = formel, administratif.
\item \textbf{Erreurs interdites :} \eng{for to V}, \eng{for V-ing} (but), séparation incorrecte (\eng{*look after them} $\rightarrow$ \eng{look them after}), traduction littérale.
\end{itemize}
\end{aretenir}

\coursec{Règle 1 : formation et types de phrasal verbs}

Dans la section précédente, nous avons observé des phrasal verbs en contexte, à travers un texte sur la citoyenneté et les droits humains : \eng{stand up for our rights}, \eng{carry out a campaign}, \eng{give up}, \eng{look after}… Nous avons remarqué que ces verbes sont formés de deux mots (parfois trois) et que leur sens ne se devine pas toujours à partir des mots séparés. Il est temps maintenant de dégager la règle : comment se forment les phrasal verbs, et combien de types existe-t-il ? C'est ce que nous allons établir rigoureusement dans cette section, car la maîtrise de ces verbes est indispensable à l'épreuve de grammaire du Baccalauréat.

\courssub{La structure générale : verbe + particule}

\begin{definition}[Phrasal verb]
Un \cle{phrasal verb} (verbe à particule) est un verbe composé de deux éléments : un \cle{verbe} de base (\eng{carry}, \eng{look}, \eng{give}, \eng{turn}) suivi d'une \cle{particule} (\eng{up}, \eng{down}, \eng{out}, \eng{off}, \eng{on}, \eng{in}, \eng{for}, \eng{after}, \eng{about}, \eng{away}). L'ensemble forme une unité de sens : \eng{carry out} ne veut pas dire « porter dehors » mais « mener, réaliser ». Certains phrasal verbs comportent une troisième particule : \eng{stand up for} (« défendre »), \eng{put up with} (« tolérer »), \eng{look forward to} (« attendre avec impatience »).
\end{definition}

La particule ressemble souvent à une préposition, mais elle n'en est pas une : elle modifie le sens du verbe lui-même. Comparez :

\eng{The plane took off.} « L'avion a décollé. » (\eng{take off} = un seul verbe, sens nouveau)

\eng{He took the book off the table.} « Il a pris le livre sur la table. » (ici \eng{off} est une vraie préposition : « de, dessus »)

\begin{propriete}[La nuance apportée par les particules fréquentes]
Chaque particule apporte souvent une nuance reconnaissable. Connaître ces nuances aide à deviner le sens d'un phrasal verb inconnu dans un texte de compréhension écrite.
\end{propriete}

\begin{center}
\begin{tabularx}{\linewidth}{|l|X|X|}
\hline
\rowcolor{popblueL} \textbf{Particule} & \textbf{Nuance fréquente} & \textbf{Exemples} \\
\hline
\eng{up} & achèvement, élévation, augmentation & \eng{give up} (abandonner), \eng{stand up} (se lever), \eng{speak up} (parler plus fort) \\
\hline
\eng{out} & extension, distribution, apparition & \eng{carry out} (mener à bien), \eng{hand out} (distribuer), \eng{break out} (éclater) \\
\hline
\eng{off} & départ, arrêt, séparation & \eng{take off} (décoller), \eng{turn off} (éteindre), \eng{call off} (annuler) \\
\hline
\eng{down} & réduction, refus, chute & \eng{turn down} (refuser), \eng{cut down} (réduire), \eng{break down} (tomber en panne) \\
\hline
\eng{on} & continuation, contact & \eng{carry on} (continuer), \eng{put on} (mettre un vêtement) \\
\hline
\eng{away} & éloignement, disparition & \eng{run away} (s'enfuir), \eng{throw away} (jeter) \\
\hline
\eng{after / for} & poursuite, recherche & \eng{look after} (s'occuper de), \eng{stand up for} (défendre) \\
\hline
\end{tabularx}
\end{center}

\begin{attention}[Ces nuances ne sont pas des lois]
Le tableau ci-dessus donne des tendances, pas des règles absolues : \eng{give up} (« abandonner ») n'a rien d'une « élévation ». La seule méthode sûre reste d'apprendre chaque phrasal verb comme un mot de vocabulaire à part entière, avec sa traduction et son type (séparable ou inséparable).
\end{attention}

\courssub{Les trois types de phrasal verbs}

Tous les phrasal verbs ne se comportent pas de la même façon vis-à-vis de leur complément. Il existe trois grands types, et les distinguer est indispensable pour construire des phrases correctes.

\textbf{Type 1 : les phrasal verbs intransitifs (sans complément)}\par

Ces verbes n'admettent \emph{aucun} complément d'objet. La particule est soudée au verbe et rien ne peut se placer entre les deux.

\begin{exemplebox}[Phrasal verbs intransitifs]
\eng{War broke out in the region.} « La guerre a éclaté dans la région. »

\eng{Please stand up when the judge enters.} « Levez-vous, s'il vous plaît, quand le juge entre. »

\eng{The activist never gave up.} « Le militant n'a jamais abandonné. »
\end{exemplebox}

\textbf{Type 2 : les phrasal verbs transitifs séparables}\par

Ces verbes admettent un complément, et ce complément peut se placer \emph{avant} ou \emph{après} la particule quand c'est un nom. C'est le type le plus délicat pour les francophones.

\begin{propriete}[La règle du verbe séparable]
Avec un phrasal verb séparable, un \cle{nom} complément peut aller avant ou après la particule : \eng{carry out the plan} = \eng{carry the plan out} (« mener le projet à bien »). Mais un \cle{pronom} (\eng{it}, \eng{them}, \eng{him}…) va \textbf{toujours} au milieu, entre le verbe et la particule : \eng{carry it out}, jamais \eng{carry out it}.
\end{propriete}

\begin{exemplebox}[Verbes séparables : nom et pronom]
\eng{They carried out the campaign.} = \eng{They carried the campaign out.} « Ils ont mené la campagne. »

\eng{They carried it out.} « Ils l'ont menée. »

\eng{The committee turned down our proposal.} = \eng{The committee turned our proposal down.} « Le comité a refusé notre proposition. »

\eng{The committee turned it down.} « Le comité l'a refusée. »
\end{exemplebox}

\begin{popfigure}
\begin{center}
\begin{tikzpicture}[
  box/.style={draw=popblue, fill=popblueL, rounded corners=3pt, minimum height=0.9cm, align=center, font=\bfseries},
  arr/.style={-{Stealth[length=3mm]}, thick, popdark}]
  \node[box, align=center] (v) at (0,0){VERB\\ \eng{carry}};
  \node[box, fill=poporangeL, draw=poporange, align=center] (p) at (4,0){PRONOUN\\ \eng{it}};
  \node[box, fill=popgreenL, draw=popgreen, align=center] (pa) at (8,0){PARTICLE\\ \eng{out}};
  \draw[arr] (v) -- (p);
  \draw[arr] (p) -- (pa);
  \node[below=0.35cm of p, align=center, font=\itshape] {ordre obligatoire avec un pronom};
\end{tikzpicture}
\end{center}
\legende{Structure d'un phrasal verb séparable avec pronom : \eng{carry it out}.}
\end{popfigure}

\textbf{Type 3 : les phrasal verbs transitifs inséparables}\par

Ces verbes admettent un complément, mais celui-ci se place \emph{toujours après} la particule, qu'il s'agisse d'un nom ou d'un pronom. On ne peut jamais insérer quoi que ce soit entre le verbe et la particule.

\begin{exemplebox}[Phrasal verbs inséparables]
\eng{She looks after the children.} « Elle s'occupe des enfants. » (jamais \eng{look the children after})

\eng{She looks after them.} « Elle s'occupe d'eux. » (jamais \eng{looks them after})

\eng{We must stand up for our rights.} « Nous devons défendre nos droits. »

\eng{The police are looking into the case.} « La police enquête sur l'affaire. »
\end{exemplebox}

\begin{attention}[Piège : \eng{stand up} et \eng{stand up for}]
\eng{Stand up} seul est intransitif : « se lever ». Pour dire « défendre », il faut ajouter la préposition \eng{for} : \eng{stand up for human rights}. On obtient alors un verbe inséparable à trois mots : \eng{We stand up for them.} « Nous les défendons. »
\end{attention}

\begin{popfigure}
\begin{center}
\begin{tikzpicture}[
  head/.style={draw=popdark, fill=popblue, text=white, rounded corners=3pt, align=center, font=\bfseries, minimum width=3.4cm, minimum height=1cm},
  ex/.style={align=center, font=\small, below=0.25cm}]
  \node[head, align=center] (a) at (0,0){Intransitif\\ (pas de complément)};
  \node[head, fill=poporange, align=center] (b) at (5.4,0){Séparable\\ (nom : 2 positions)};
  \node[head, fill=popgreen, align=center] (c) at (10.8,0){Inséparable\\ (complément après)};
  \node[ex] at (a.south) {\eng{War broke out.}};
  \node[ex] at (b.south) {\eng{carry it out}};
  \node[ex] at (c.south) {\eng{look after them}};
\end{tikzpicture}
\end{center}
\legende{Les trois types de phrasal verbs et un exemple de chaque.}
\end{popfigure}

\courssub{La règle d'or du pronom}

C'est la règle la plus testée à l'examen et la plus souvent violée par les francophones. Retenez-la comme une formule.

\begin{aretenir}[Règle d'or]
Verbe séparable + pronom : le pronom se \cle{coince} entre le verbe et la particule.

\eng{carry it out} — \eng{turn it down} — \eng{pick them up} — \eng{put it on}

Verbe inséparable + pronom : le pronom suit la particule.

\eng{look after them} — \eng{stand up for them} — \eng{look into it}
\end{aretenir}

\begin{attention}[Erreur fréquente des francophones]
Ne dites jamais \eng{carry out it} par calque du français « le réaliser » ou « réaliser-le ». En français, le pronom précède le verbe (« ils l'ont menée ») ; en anglais, avec un verbe séparable, le pronom s'intercale entre le verbe et la particule : \eng{they carried it out}. De même, on écrit \eng{she turned it down} (« elle l'a refusée »), jamais \eng{she turned down it}.
\end{attention}

\courssub{Comparaison avec le français}

Le français n'a pas de phrasal verbs : il utilise soit un verbe simple, soit un verbe pronominal, soit une expression entière. D'où l'importance d'apprendre chaque phrasal verb comme un mot de vocabulaire à part entière.

\begin{center}
\begin{tabularx}{\linewidth}{|X|X|X|}
\hline
\rowcolor{popblueL} \textbf{Français} & \textbf{English} & \textbf{Remarque} \\
\hline
mener, réaliser & \eng{carry out} & séparable : \eng{carry it out} \\
\hline
refuser (une offre) & \eng{turn down} & séparable : \eng{turn it down} \\
\hline
s'occuper de & \eng{look after} & inséparable : \eng{look after them} \\
\hline
défendre (une cause) & \eng{stand up for} & inséparable, 3 mots \\
\hline
abandonner & \eng{give up} & intransitif, ou suivi d'un verbe en \eng{-ing} : \eng{give up smoking} (« arrêter de fumer ») \\
\hline
éclater (guerre) & \eng{break out} & intransitif, jamais de complément \\
\hline
distribuer & \eng{hand out} & séparable : \eng{hand them out} \\
\hline
annuler & \eng{call off} & séparable : \eng{call it off} \\
\hline
\end{tabularx}
\end{center}

\courssub{Comment savoir si un verbe est séparable ?}

\begin{methode}[Le test du pronom \eng{it}]
Pour déterminer si un phrasal verb est séparable, procédez ainsi :
\begin{enumerate}
\item Prenez le verbe avec un complément nom : \eng{carry out the plan}.
\item Remplacez le complément par le pronom \eng{it}.
\item Placez \eng{it} entre le verbe et la particule : \eng{carry it out}.
\item Si la phrase sonne juste, le verbe est séparable. Si elle est impossible et qu'il faut dire \eng{look after it}, le verbe est inséparable.
\end{enumerate}
Test : \eng{look after the children} $\rightarrow$ \eng{look the children after} ? Impossible. Donc \eng{look after} est inséparable : \eng{look after them}.
\end{methode}

\begin{exoresolu}[Placer le complément au bon endroit]
Énoncé — Réécrivez chaque phrase en remplaçant le complément en gras par un pronom (\eng{it} ou \eng{them}).

1. The students carried out \textbf{the project}.

2. The mayor turned down \textbf{the request}.

3. My aunt looks after \textbf{the orphans}.

4. The volunteers handed out \textbf{the leaflets}.

\tcblower
Solution détaillée —

1. \eng{Carry out} est séparable (le test \eng{carry it out} fonctionne) : le pronom va au milieu. \eng{The students carried it out.} « Les élèves l'ont mené à bien. »

2. \eng{Turn down} est séparable : \eng{The mayor turned it down.} « Le maire l'a refusée. »

3. \eng{Look after} est inséparable : le pronom reste après la particule. \eng{My aunt looks after them.} « Ma tante s'occupe d'eux. »

4. \eng{Hand out} est séparable : \eng{The volunteers handed them out.} « Les volontaires les ont distribués. »
\end{exoresolu}

\begin{exercice}[À vous de jouer]
Traduisez en anglais en utilisant le phrasal verb indiqué et un pronom.

1. Ils ont mené la campagne. (\eng{carry out}, pronom)

2. Elle a refusé l'offre. (\eng{turn down}, pronom)

3. Il s'occupe des enfants. (\eng{look after}, pronom)

4. Nous devons défendre nos droits. (\eng{stand up for}, pronom)
\end{exercice}

\begin{corrige}[Exercice : placer le pronom]
1. \eng{They carried it out.} (séparable : pronom au milieu)

2. \eng{She turned it down.} (séparable : pronom au milieu)

3. \eng{He looks after them.} (inséparable : pronom après la particule)

4. \eng{We must stand up for them.} (inséparable à trois mots : pronom après \eng{for})
\end{corrige}

Maintenant que la mécanique des trois types est acquise — et surtout la règle d'or du pronom —, nous pouvons passer à la règle 2 : les phrasal verbs les plus utiles pour parler de citoyenneté et de droits humains.

\coursec{Règle 2 : les phrasal verbs de la citoyenneté et des droits}

Dans la section précédente, tu as appris à \emph{former} les phrasal verbs et à distinguer les verbes séparables (\eng{turn the offer down}), inséparables (\eng{look after the children}) et intransitifs (\eng{speak out}). Il est temps maintenant de mettre cette mécanique au service du thème du chapitre : la citoyenneté et les droits humains. En effet, les journaux anglophones, les discours politiques et les rapports des ONG utilisent constamment un petit groupe de phrasal verbs pour parler de défense des droits, d'engagement civique et de changement social. Les maîtriser, c'est pouvoir lire la presse anglophone et écrire des compositions riches et naturelles.

\courssub{Observation : un texte pour découvrir les verbes en contexte}

Lis attentivement le texte suivant. Repère les huit phrasal verbs en gras et essaie de deviner leur sens global grâce au contexte, avant de consulter le glossaire.

\begin{textebox}[Youth for Justice: a student association in Bamenda]
Last year, a group of students in Bamenda decided to \textbf{stand up for} the rights of children who could not afford school fees. They refused to \textbf{give up}, even when some people told them their project was unrealistic. First, they \textbf{carried out} a survey in three neighbourhoods to understand why so many children were out of school. The results were shocking: many parents simply could not pay.

The students then \textbf{put forward} a proposal to the local council: free evening classes for working children. At first, the council \textbf{turned down} the idea, saying there was no money. But the students did not stop. They \textbf{spoke out against} injustice on local radio stations and in community meetings. Little by little, their campaign \textbf{brought about} a real change: two churches offered their halls, and retired teachers volunteered to teach.

Today, the association also \textbf{looks after} fifty children every evening, giving them lessons, books and a hot meal. “Nobody will defend our community for us,” says Mirabel, the president. “We must act ourselves.”

\emph{Glossaire :} school fees (n. pl.) : frais de scolarité ; survey (n.) : enquête ; council (n.) : conseil municipal ; hall (n.) : salle ; to volunteer (v.) : se porter volontaire ; hot meal (n.) : repas chaud.
\end{textebox}

\textbf{Questions de compréhension :}
\begin{enumerate}
\item Why did the students carry out a survey?
\item What did the council do with the first proposal?
\item How did the students bring about a change?
\end{enumerate}

\courssub{Les huit phrasal verbs clés : sens, type et emploi}

\begin{definition}[Phrasal verbs de la citoyenneté et des droits]
Un \cle{phrasal verb} a un sens \emph{global} qui ne se devine pas en additionnant le verbe et la particule. Voici les huit verbes essentiels du thème « Citizenship and Human Rights » :
\begin{itemize}
\item \eng{stand up for} = défendre (des droits, une cause, une personne) — \textbf{inséparable} ;
\item \eng{give up} = abandonner, renoncer — \textbf{séparable} ou \textbf{intransitif} ;
\item \eng{carry out} = mener, réaliser (une enquête, un projet, une mission) — \textbf{séparable} ;
\item \eng{look after} = s'occuper de, prendre soin de — \textbf{inséparable} ;
\item \eng{bring about} = provoquer, entraîner (un changement) — \textbf{séparable} ;
\item \eng{put forward} = proposer, présenter (une idée, un plan, un candidat) — \textbf{séparable} ;
\item \eng{turn down} = refuser (une offre, une demande) — \textbf{séparable} ;
\item \eng{speak out (against)} = dénoncer publiquement, prendre la parole — \textbf{intransitif}.
\end{itemize}
\end{definition}

\begin{exemplebox}[Une phrase traduite pour chaque verbe, en contexte « droits humains »]
\begin{itemize}
\item \eng{Citizens must stand up for their rights.} « Les citoyens doivent défendre leurs droits. »
\item \eng{The activists never gave up the fight for justice.} « Les militants n'ont jamais abandonné le combat pour la justice. » (ou : \eng{They never gave up.} — intransitif)
\item \eng{The NGO carried out a survey in rural areas.} « L'ONG a mené une enquête dans les zones rurales. »
\item \eng{Volunteers look after refugees at the centre.} « Des bénévoles s'occupent des réfugiés au centre. »
\item \eng{Education can bring about real social change.} « L'éducation peut provoquer un véritable changement social. »
\item \eng{The committee put forward a new proposal.} « Le comité a présenté une nouvelle proposition. »
\item \eng{The minister turned down their request for a meeting.} « Le ministre a refusé leur demande de rendez-vous. »
\item \eng{Journalists spoke out against corruption.} « Des journalistes ont dénoncé publiquement la corruption. »
\end{itemize}
\end{exemplebox}

\begin{propriete}[Rappel de comportement syntaxique]
\begin{itemize}
\item Verbes \textbf{inséparables} : la particule reste collée au verbe, même avec un pronom : \eng{stand up for them} (jamais \emph{stand up them for}) ; \eng{look after them} (jamais \emph{look them after}).
\item Verbes \textbf{séparables} : avec un nom, les deux ordres sont possibles (\eng{carry out a survey} / \eng{carry a survey out}) ; avec un \textbf{pronom}, la séparation est \emph{obligatoire} : \eng{they carried it out}, \eng{she turned it down}, \eng{he put it forward}.
\item Verbe \textbf{intransitif} : \eng{speak out} ne prend pas de complément direct ; on ajoute \eng{against} + nom pour préciser la cible : \eng{speak out against injustice}.
\end{itemize}
\end{propriete}

\begin{center}
\begin{tabularx}{\linewidth}{|l|l|X|X|}
\hline
\rowcolor{popblueL} Phrasal verb & Type & Français & Exemple \\
\hline
stand up for & inséparable & défendre & \eng{Stand up for your rights!} \\
\hline
give up & sép. / intr. & abandonner & \eng{Don't give up hope.} \\
\hline
carry out & séparable & mener, réaliser & \eng{They carried out an inquiry.} \\
\hline
look after & inséparable & s'occuper de & \eng{She looks after orphans.} \\
\hline
bring about & séparable & provoquer & \eng{Laws bring about change.} \\
\hline
put forward & séparable & proposer, présenter & \eng{He put forward an idea.} \\
\hline
turn down & séparable & refuser & \eng{They turned down the offer.} \\
\hline
speak out (against) & intransitif & dénoncer & \eng{She spoke out against violence.} \\
\hline
\end{tabularx}
\end{center}

\begin{popfigure}
\begin{tikzpicture}[mindmap, grow cyclic,
  every node/.style={concept, align=center, font=\small},
  root concept/.append style={concept color=popblue, text=white, font=\bfseries},
  level 1 concept/.append style={sibling angle=90, level distance=3.4cm},
  level 2 concept/.append style={sibling angle=60, level distance=2.6cm}]
\node[root concept, align=center]{Phrasal verbs\\ of citizenship}
  child[concept color=popgreenL]{ node[align=center]{Defend\\ \& resist}
    child[concept color=popgreen!40]{ node[align=center]{stand up for\\ (insép.)} }
    child[concept color=popgreen!40]{ node[align=center]{speak out\\ (intr.)} }
    child[concept color=popgreen!40]{ node[align=center]{give up\\ (sép./intr.)} }
  }
  child[concept color=poporangeL]{ node[align=center]{Act\\ \& change}
    child[concept color=poporange!40]{ node[align=center]{carry out\\ (sép.)} }
    child[concept color=poporange!40]{ node[align=center]{bring about\\ (sép.)} }
    child[concept color=poporange!40]{ node[align=center]{look after\\ (insép.)} }
  }
  child[concept color=popgoldL]{ node[align=center]{Propose\\ \& refuse}
    child[concept color=popgold!40]{ node[align=center]{put forward\\ (sép.)} }
    child[concept color=popgold!40]{ node[align=center]{turn down\\ (sép.)} }
  };
\end{tikzpicture}
\legende{Carte mentale : classer les huit phrasal verbs par fonction communicative (défendre, agir, proposer) facilite la mémorisation.}
\end{popfigure}

\courssub{Comparaison avec le français et pièges}

En français, on utilise souvent un verbe unique et soutenu (\emph{défendre}, \emph{mener}, \emph{proposer}) ; l'anglais courant préfère le phrasal verb, surtout à l'oral et dans la presse. Attention donc à ne pas chercher un équivalent mot à mot.

\begin{center}
\begin{tabularx}{\linewidth}{|X|X|X|}
\hline
\rowcolor{popblueL} Français & English (naturel) & Remarque \\
\hline
défendre ses droits & stand up for one's rights & sens global, pas « se lever » \\
\hline
mener une enquête & carry out a survey & « carry » seul = porter ! \\
\hline
proposer une idée & put forward an idea & plus formel que « suggest » \\
\hline
refuser une offre & turn down an offer & « refuse » existe aussi \\
\hline
s'occuper des enfants & look after the children & inséparable : \eng{look after them} \\
\hline
provoquer un changement & bring about a change & ne pas confondre avec « bring back » \\
\hline
dénoncer la corruption & speak out against corruption & intransitif + « against » \\
\hline
abandonner & give up & intransitif possible : \eng{Never give up!} \\
\hline
\end{tabularx}
\end{center}

\begin{attention}[Ne traduis jamais mot à mot !]
Le sens d'un phrasal verb est \textbf{global}, pas la somme des mots. \eng{Stand up for} ne veut pas dire « se lever pour » mais « défendre » ; \eng{give up} ne veut pas dire « donner vers le haut » mais « abandonner ». Autres erreurs fréquentes des francophones :
\begin{itemize}
\item Séparer un verbe inséparable : \emph{look them after} est faux ; dis \eng{look after them}.
\item Ne pas séparer avec un pronom : \emph{turn down it} est faux ; dis \eng{turn it down}.
\item Oublier la particule : \emph{they carried a survey} signifie « ils ont porté une enquête » ! Dis \eng{they carried out a survey}.
\item Traduire « renoncer à » par \emph{renounce to} : dis \eng{give up} + nom ou \eng{-ing} : \eng{She gave up smoking.} « Elle a arrêté de fumer. »
\end{itemize}
\end{attention}

\begin{exoresolu}[Application : complète et transforme]
Énoncé — Complète avec le bon phrasal verb conjugué : (a) The students \ldots\ldots\ldots a petition to the mayor. (put forward) (b) We must \ldots\ldots\ldots human rights everywhere. (stand up for) (c) Remplace le nom par un pronom : \eng{The council turned down the proposal.}
\tcblower
Solution — (a) \eng{The students put forward a petition to the mayor.} « Les élèves ont présenté une pétition au maire. » (\eng{put} est irrégulier : put–put–put ; au prétérit, la forme ne change pas.) (b) \eng{We must stand up for human rights everywhere.} « Nous devons défendre les droits humains partout. » (inséparable : la particule \eng{for} reste collée.) (c) Le complément devient le pronom \eng{it} ; avec un verbe séparable, le pronom se place \emph{entre} le verbe et la particule : \eng{The council turned it down.} « Le conseil l'a refusée. »
\end{exoresolu}

\begin{methode}[Comment mémoriser durablement ces huit verbes]
\begin{enumerate}
\item N'apprends \textbf{jamais} un phrasal verb isolément dans une liste : retiens-le dans une \textbf{phrase complète} liée aux droits humains, par exemple \eng{Citizens must stand up for their rights.}
\item Classe-les par \textbf{fonction} (défendre / agir / proposer), comme dans la carte mentale ci-dessus : le cerveau retient mieux les familles de sens.
\item Marque à côté de chaque verbe son \textbf{type} (sép. / insép. / intr.) et teste-toi avec un pronom : \eng{carry out a survey} $\rightarrow$ \eng{carry it out}.
\item Réutilise-les dans \textbf{tes propres phrases} sur le Cameroun : \eng{Young people in Douala are speaking out against unemployment.}
\end{enumerate}
\end{methode}

\begin{savaistu}[Un atout pour le Bac]
Au baccalauréat, les correcteurs valorisent les copies qui utilisent un vocabulaire authentique et idiomatique. Employer correctement un phrasal verb comme \eng{stand up for}, \eng{carry out} ou \eng{bring about} dans ta composition sur les droits humains montre que tu maîtrises l'anglais réel, celui des journaux et des discours — et pas seulement celui des dictionnaires. Un seul phrasal verb bien placé peut faire la différence entre une copie correcte et une copie excellente !
\end{savaistu}

\begin{exercice}[À toi de jouer]
\begin{enumerate}
\item Traduis : « Les bénévoles s'occupent des enfants et n'abandonnent jamais. »
\item Mets le pronom à la bonne place : \eng{The NGO carried out the project.} (it)
\item Complète : \eng{The new law will \ldots\ldots\ldots a real change in people's lives.}
\item Trouve l'erreur et corrige : \eng{We must stand up them for.}
\end{enumerate}
\end{exercice}

\begin{corrige}[Exercice « À toi de jouer »]
\begin{enumerate}
\item \eng{The volunteers look after the children and never give up.} (inséparable : \eng{look after} ; \eng{give up} intransitif.)
\item \eng{The NGO carried it out.} (pronom obligatoirement entre verbe et particule.)
\item \eng{The new law will bring about a real change in people's lives.}
\item Faute : verbe inséparable séparé. Correction : \eng{We must stand up for them.}
\end{enumerate}
\end{corrige}

Tu connais désormais les huit phrasal verbs incontournables du vocabulaire citoyen, leur sens global et leur comportement syntaxique. Dans la section suivante, tu apprendras à les combiner avec les \emph{expressions of purpose} pour expliquer \emph{pourquoi} les citoyens s'engagent : \eng{They spoke out in order to defend freedom.}

\coursec{Règle 3 : exprimer le but (expressions of purpose)}

Dans la section précédente, nous avons appris à reconnaître et à utiliser les phrasal verbs de la citoyenneté et des droits : \eng{stand up for}, \eng{carry out}, \eng{bring about}, \eng{take part in}… Mais un citoyen n'agit jamais sans raison. Pourquoi vote-t-il ? Pourquoi manifeste-t-il ? Pour répondre à ces questions, l'anglais dispose de structures précises pour exprimer le \cle{but} (\eng{purpose}). C'est l'objet de cette règle 3 : dire \emph{pourquoi} l'on fait quelque chose, et combiner ces structures avec nos phrasal verbs.

\courssub{Observation : pourquoi les citoyens agissent-ils ?}

Lisons d'abord ce court texte. Soulignez mentalement toutes les expressions qui répondent à la question \eng{Why?}.

\begin{textebox}[Why citizens act — a short report from Yaoundé]
Every morning, Mrs. Etoundi leaves her house early \textbf{to} open her small shop near the Mokolo market. She works hard \textbf{to} support her family and \textbf{to} pay her children's school fees. Last year, she joined a women's association \textbf{in order to} learn about her rights as a trader. She attends every meeting \textbf{so as not to} miss any important information.

In the same neighbourhood, a group of students created a club \textbf{to} clean up the streets. They meet every Saturday \textbf{to} carry out their project, and they visit schools \textbf{in order to} educate younger children about the environment. “We act \textbf{so that} our city can become cleaner,” explains their leader, a nineteen-year-old boy. “We speak to the mayor \textbf{in order to} bring about real change.”

Not far away, Mr. Kamga votes at every election \textbf{to} choose honest leaders. He reads the newspapers every day \textbf{so as not to} be manipulated. “I keep myself informed \textbf{so that} my vote can really count,” he says. These ordinary citizens show one simple truth: every action has a purpose.
\end{textebox}

\textbf{Glossaire :} \emph{trader} (nom) : commerçant ; \emph{to clean up} (phrasal verb) : nettoyer ; \emph{to carry out} (phrasal verb) : mener à bien, réaliser ; \emph{to bring about} (phrasal verb) : provoquer, entraîner ; \emph{mayor} (nom) : maire ; \emph{to count} (verbe) : compter (avoir de l'importance) ; \emph{purpose} (nom) : but, objectif.

\textbf{Questions de compréhension :}
\begin{enumerate}
\item Why does Mrs. Etoundi work hard?
\item Why did she join a women's association?
\item Why do the students visit schools?
\item Why does Mr. Kamga read the newspapers every day?
\end{enumerate}

Observez les réponses possibles : \eng{She works hard to support her family.} / \eng{She joined the association in order to learn about her rights.} / \eng{They visit schools to educate younger children.} / \eng{He reads the newspapers so that his vote can count.} Trois structures apparaissent : \eng{to}, \eng{in order to}, \eng{so as (not) to} suivis d'un infinitif, et \eng{so that} suivi d'une proposition complète. Voyons maintenant la règle en détail.

\courssub{La structure de base : to + infinitif}

\begin{propriete}[To + infinitif pour exprimer le but]
Pour exprimer le \cle{but} d'une action (répondre à la question \eng{Why?} ou \eng{What... for?}), l'anglais utilise \textbf{to + infinitif} (la base verbale) :
\begin{center}
\textbf{sujet + verbe + (complément) + TO + base verbale}
\end{center}
\eng{She joined the NGO to help children.} « Elle a rejoint l'ONG pour aider les enfants. »

C'est la structure la plus fréquente, dans tous les registres de langue. L'infinitif de but se place généralement \textbf{après} le verbe principal, mais il peut aussi commencer la phrase pour mettre le but en valeur : \eng{To help children, she joined the NGO.} « Pour aider les enfants, elle a rejoint l'ONG. »
\end{propriete}

\begin{exemplebox}[To + infinitif en contexte citoyen]
\begin{itemize}
\item \eng{He votes to choose his leaders.} « Il vote pour choisir ses dirigeants. »
\item \eng{We pay taxes to finance public services.} « Nous payons des impôts pour financer les services publics. »
\item \eng{They went to court to defend their rights.} « Ils sont allés au tribunal pour défendre leurs droits. »
\item \eng{She studies law to become a human rights lawyer.} « Elle étudie le droit pour devenir avocate des droits de l'homme. »
\item \eng{To participate in the elections, you must register first.} « Pour participer aux élections, il faut d'abord s'inscrire. »
\end{itemize}
\end{exemplebox}

\begin{attention}[Répondre à la question “Why...?”]
Quand on répond à une question en \eng{Why}, la réponse courte commence par \eng{To + infinitif} — jamais par \eng{Because to} ni par \eng{For to} :

\eng{Why did she join the NGO? — To stand up for children's rights.} « Pourquoi a-t-elle rejoint l'ONG ? — Pour défendre les droits des enfants. »

\eng{Why do they protest? — To bring about change.} « Pourquoi manifestent-ils ? — Pour provoquer le changement. »

Notez la différence : \eng{because} répond par une \emph{cause} (une proposition complète), \eng{to} répond par un \emph{but} (un infinitif). \eng{Why is he studying? — Because he has an exam.} (cause) / \eng{— To pass his exam.} (but).
\end{attention}

\courssub{Variantes plus formelles : in order to et so as to}

\begin{propriete}[In order to / so as to + infinitif]
\eng{In order to} et \eng{so as to} + infinitif ont \textbf{exactement le même sens} que \eng{to} seul, mais appartiennent à un registre plus \cle{soutenu} : discours officiels, dissertations, lettres formelles, textes juridiques. Ils sont très utiles au Bac pour soigner l'expression écrite.
\begin{itemize}
\item \eng{The government created a commission \textbf{in order to} investigate the case.} « Le gouvernement a créé une commission afin d'enquêter sur l'affaire. »
\item \eng{Citizens must register \textbf{so as to} take part in the vote.} « Les citoyens doivent s'inscrire afin de prendre part au vote. »
\end{itemize}
\eng{In order to} peut commencer la phrase ; \eng{so as to} ne le peut pas : \eng{In order to vote, you must be eighteen.} (correct) ; on évitera de commencer une phrase par \eng{so as to}.
\end{propriete}

\begin{propriete}[Le but négatif : in order NOT to / so as NOT to]
Pour exprimer un but \textbf{négatif} (« pour ne pas… », « afin de ne pas… »), la négation \eng{not} se place \textbf{avant} le \eng{to} :
\begin{center}
\textbf{in order NOT to + infinitif} \quad ou \quad \textbf{so as NOT to + infinitif}
\end{center}
\eng{She left early so as not to miss the meeting.} « Elle est partie tôt pour ne pas manquer la réunion. »

\eng{He spoke slowly in order not to be misunderstood.} « Il a parlé lentement pour ne pas être mal compris. »

En anglais standard, on ne dit jamais \eng{to not miss} dans cette structure : la négation précède \eng{to}. (La forme \eng{not to} seule est possible à l'oral familier : \eng{I left early not to be late.}, mais préférez \eng{so as not to} à l'écrit.)
\end{propriete}

\begin{center}
\begin{tabularx}{\linewidth}{|X|X|X|}
\hline
\rowcolor{popblueL} Structure & Registre & Exemple \\
\hline
\eng{to + infinitif} & courant, tous contextes & \eng{He works hard to carry out his duties.} \\
\hline
\eng{in order to + infinitif} & soutenu, écrit formel & \eng{They protest in order to bring about change.} \\
\hline
\eng{so as to + infinitif} & soutenu, écrit formel & \eng{She trained so as to serve her community better.} \\
\hline
\eng{in order not to / so as not to} & but négatif & \eng{He left early so as not to miss the debate.} \\
\hline
\eng{so that + sujet + verbe} & sujets différents & \eng{She speaks loudly so that everyone can hear her.} \\
\hline
\end{tabularx}
\end{center}

\courssub{Combiner phrasal verbs et expressions de but}

C'est ici que les deux règles de cette leçon se rejoignent. Le phrasal verb exprime l'\emph{action} ; l'infinitif de but exprime la \emph{finalité}. La structure est : \textbf{sujet + verbe + to / in order to + phrasal verb}, ou l'inverse (phrasal verb d'abord, but ensuite).

\begin{exemplebox}[Phrasal verb + but]
\begin{itemize}
\item \eng{He works hard \textbf{to carry out} his duties.} « Il travaille dur pour remplir ses devoirs. »
\item \eng{They protest \textbf{in order to bring about} change.} « Ils manifestent afin de provoquer le changement. »
\item \eng{She joined the association \textbf{to stand up for} women's rights.} « Elle a rejoint l'association pour défendre les droits des femmes. »
\item \eng{We organised a meeting \textbf{in order to talk about} the new law.} « Nous avons organisé une réunion afin de discuter de la nouvelle loi. »
\item \eng{The students meet every week \textbf{to clean up} their school.} « Les élèves se réunissent chaque semaine pour nettoyer leur école. »
\item \eng{They educate citizens \textbf{in order to bring about} peace.} « Ils éduquent les citoyens afin d'apporter la paix. »
\end{itemize}
\end{exemplebox}

\courssub{Exception fondamentale : so that + proposition}

Toutes les structures vues jusqu'ici supposent que \textbf{le sujet des deux actions est le même} : dans \eng{She joined the NGO to help children}, c'est \emph{elle} qui a rejoint l'ONG et c'est \emph{elle} qui veut aider les enfants. Mais que faire quand les sujets sont \textbf{différents} ?

\begin{propriete}[So that + sujet + verbe (sujets différents)]
Quand la personne qui agit et la personne visée par le but ne sont \textbf{pas les mêmes}, on ne peut pas utiliser l'infinitif. On utilise alors :
\begin{center}
\textbf{so that + sujet + can/could/will/would + verbe}
\end{center}
\eng{She speaks loudly \textbf{so that} everyone \textbf{can hear} her.} « Elle parle fort pour que tout le monde puisse l'entendre. » (Elle parle, mais c'est \emph{tout le monde} qui entend.)

\eng{We act \textbf{so that} our children \textbf{can live} freely.} « Nous agissons pour que nos enfants puissent vivre librement. »

Au passé, on utilise \eng{could} ou \eng{would} : \eng{The teacher explained the law slowly so that the students could understand it.} « Le professeur a expliqué la loi lentement pour que les élèves puissent la comprendre. »

C'est l'équivalent du français « pour que + subjonctif ». Cette structure est un complément précieux pour le Bac, car elle permet d'éviter les infinitifs impossibles. À l'écrit très formel, on rencontre aussi \eng{in order that + sujet + may/might} : \eng{The law was published in order that every citizen might know it.} « La loi fut publiée afin que chaque citoyen puisse la connaître. »
\end{propriete}

\begin{attention}[Deux sujets différents = jamais l'infinitif]
Phrase impossible : \eng{She speaks loudly to everyone hear her.} Pourquoi ? Parce que le sujet de \eng{hear} n'est pas \eng{she} mais \eng{everyone}. Il faut une proposition complète : \eng{so that everyone can hear her.}

Test rapide : demandez-vous « \emph{qui} réalise l'action du but ? ». Si c'est la même personne que le sujet principal, utilisez \eng{to / in order to / so as to}. Sinon, utilisez \eng{so that}.
\end{attention}

\begin{popfigure}
\begin{tikzpicture}[
  node distance=0.9cm and 1.6cm,
  start/.style={rectangle, rounded corners, draw=popblue, fill=popblueL, thick, align=center, font=\bfseries, minimum height=1cm, inner sep=6pt},
  decision/.style={diamond, draw=poporange, fill=poporangeL, thick, align=center, aspect=2.2, inner sep=2pt},
  result/.style={rectangle, rounded corners, draw=popgreen, fill=popgreenL, thick, align=center, minimum height=1cm, inner sep=6pt},
  lbl/.style={font=\bfseries\small, fill=white, inner sep=2pt},
  arr/.style={-{Stealth[length=3mm]}, thick, popdark}
]
\node[start, align=center] (purpose){PURPOSE\\ exprimer le but};
\node[decision, below=of purpose, align=center] (same){Le sujet des deux\\ actions est-il le même ?};
\node[result, below left=1.1cm and 0.2cm of same, align=center] (yes){OUI : \textbf{to / in order to /}\\ \textbf{so as to} + infinitif\\ \eng{He votes to choose his leaders.}};
\node[result, below right=1.1cm and 0.2cm of same, align=center] (no){NON : \textbf{so that} + sujet + verbe\\ (can / could / will / would)\\ \eng{She speaks loudly so that}\\ \eng{everyone can hear her.}};
\draw[arr] (purpose) -- (same);
\draw[arr] (same) -| node[lbl, pos=0.25] {même sujet} (yes);
\draw[arr] (same) -| node[lbl, pos=0.25] {sujets différents} (no);
\end{tikzpicture}
\legende{Organigramme de choix : comment exprimer le but en anglais.}
\end{popfigure}

\courssub{Comparaison avec le français et pièges des francophones}

\begin{center}
\begin{tabularx}{\linewidth}{|X|X|X|}
\hline
\rowcolor{popblueL} Français & English & Remarque \\
\hline
Il vote pour choisir ses dirigeants. & \eng{He votes to choose his leaders.} & « pour » + infinitif = \eng{to} + infinitif \\
\hline
Elle est partie tôt pour ne pas manquer la réunion. & \eng{She left early so as not to miss the meeting.} & « pour ne pas » = \eng{so as not to} \\
\hline
Ils éduquent les citoyens afin d'apporter la paix. & \eng{They educate citizens in order to bring about peace.} & « afin de » = \eng{in order to} \\
\hline
Nous agissons pour que nos enfants puissent vivre librement. & \eng{We act so that our children can live freely.} & « pour que » + subjonctif = \eng{so that} + \eng{can} \\
\hline
\end{tabularx}
\end{center}

\begin{attention}[Les trois erreurs classiques des francophones]
\begin{enumerate}
\item \textbf{Le calque « for to ».} En français on dit « pour aider », donc beaucoup d'élèves écrivent \eng{for to help}. C'est \textbf{faux} : le but s'exprime avec \eng{to} + infinitif, jamais \eng{for to}. \eng{She joined the NGO to help children} (et non \eng{for to help}). Attention : \eng{for + nom} existe (\eng{I went to the office for information} « je suis allé au bureau pour des informations »), et \eng{for + -ing} sert à décrire la fonction d'un objet (\eng{A pen is used for writing.} « Un stylo sert à écrire. »), mais devant un \textbf{verbe à l'infinitif} exprimant le but d'une action, c'est toujours \eng{to}.
\item \textbf{La négation mal placée.} On entend souvent \eng{to not miss} par calque de « pour ne pas manquer ». En anglais standard : \eng{so as \textbf{not to} miss} / \eng{in order \textbf{not to} miss}. Le \eng{not} se place \textbf{avant} \eng{to}.
\item \textbf{L'infinitif avec deux sujets différents.} \eng{I explain to you understand} est impossible. Deux sujets différents exigent \eng{so that} : \eng{I explain slowly so that you can understand.}
\end{enumerate}
\end{attention}

\begin{exoresolu}[Transformez selon le modèle]
Énoncé : exprimez le but en reliant les deux phrases avec la structure indiquée entre parenthèses, comme dans le modèle.

Modèle : She joined the NGO. She wants to help children. (\eng{to}) → \eng{She joined the NGO to help children.}

\begin{enumerate}
\item He works hard. He wants to carry out his duties. (\eng{to})
\item They protest peacefully. They want to bring about change. (\eng{in order to})
\item She left the house early. She did not want to miss the meeting. (\eng{so as not to})
\item The teacher speaks loudly. Everyone can hear her. (\eng{so that})
\item We act today. Our children will live freely tomorrow. (\eng{so that})
\end{enumerate}

\tcblower
Solution détaillée :
\begin{enumerate}
\item Même sujet (\eng{he}) dans les deux phrases → infinitif de but simple : \eng{He works hard to carry out his duties.} « Il travaille dur pour remplir ses devoirs. »
\item Même sujet (\eng{they}), registre soutenu demandé : \eng{They protest peacefully in order to bring about change.} « Ils manifestent pacifiquement afin de provoquer le changement. »
\item But \textbf{négatif} : le \eng{not} se place avant \eng{to} : \eng{She left the house early so as not to miss the meeting.} « Elle a quitté la maison tôt pour ne pas manquer la réunion. » (Jamais \eng{so as to not miss}.)
\item Sujets différents (\eng{the teacher} / \eng{everyone}) → \eng{so that} + \eng{can} : \eng{The teacher speaks loudly so that everyone can hear her.} « Le professeur parle fort pour que tout le monde puisse l'entendre. »
\item Sujets différents (\eng{we} / \eng{our children}) et futur → \eng{so that} + \eng{will} ou \eng{can} : \eng{We act today so that our children can live freely tomorrow.} « Nous agissons aujourd'hui pour que nos enfants puissent vivre librement demain. »
\end{enumerate}
\end{exoresolu}

\begin{aretenir}[L'essentiel de la règle 3]
\begin{itemize}
\item But, même sujet : \textbf{to + infinitif} (courant) ; \textbf{in order to} / \textbf{so as to} + infinitif (soutenu).
\item But négatif : \textbf{in order NOT to} / \textbf{so as NOT to} + infinitif — jamais \eng{to not}.
\item Réponse à \eng{Why...?} : \eng{To stand up for children's rights.}
\item Sujets différents : \textbf{so that + sujet + can/could/will/would + verbe} (= « pour que »).
\item Interdit : \eng{for to + infinitif}. Devant un verbe, le but = \eng{to}, jamais \eng{for to}. (\eng{for + -ing} = fonction d'un objet : \eng{a tool for building}.)
\end{itemize}
\end{aretenir}

\begin{exercice}[À vous de jouer]
\begin{enumerate}
\item Traduisez : « Elle a rejoint l'ONG pour défendre les droits des enfants. » / « Il est parti tôt afin de ne pas manquer le débat. » / « Nous parlons fort pour que tout le monde puisse nous entendre. »
\item Complétez avec \eng{to}, \eng{in order not to} ou \eng{so that} : (a) Citizens vote \ldots\ choose their leaders. (b) He whispered \ldots\ wake the baby. (c) The mayor explained the project \ldots\ everyone could give an opinion.
\item Écrivez trois phrases sur la citoyenneté au Cameroun en combinant un phrasal verb (\eng{stand up for}, \eng{carry out}, \eng{bring about}, \eng{take part in}) et une expression de but différente à chaque fois.
\end{enumerate}
\end{exercice}

\begin{corrige}[Exercice « À vous de jouer »]
\begin{enumerate}
\item \eng{She joined the NGO to stand up for children's rights.} / \eng{He left early so as not to miss the debate.} (ou \eng{in order not to}) / \eng{We speak loudly so that everyone can hear us.}
\item (a) \eng{to} — même sujet ; (b) \eng{in order not to} — but négatif ; (c) \eng{so that} — sujets différents (\eng{the mayor} / \eng{everyone}).
\item Exemples de réponses : \eng{Young people take part in elections to choose their future.} / \eng{Journalists write articles in order to bring about justice.} / \eng{Teachers explain the law so that students can carry out their duties as citizens.}
\end{enumerate}
\end{corrige}

\coursec{Comparaison français/anglais et pièges des francophones}

Dans la section précédente, nous avons appris à exprimer le but en anglais (\eng{to vote}, \eng{in order to protect}, \eng{so that everyone can speak}). Mais une difficulté subsiste pour le francophone : le français et l'anglais ne « découpent » pas les idées de la même façon. Là où le français utilise \textbf{un seul verbe} (\emph{abandonner}, \emph{défendre}, \emph{réaliser}), l'anglais courant préfère souvent \textbf{un verbe + une particule} (\eng{give up}, \eng{stand up for}, \eng{carry out}). Cette section confronte les deux systèmes et vous donne une stratégie de vérification pour ne plus jamais tomber dans les cinq pièges classiques.

\courssub{Observation : deux langues, deux logiques}

\begin{textebox}[Two ways of saying the same thing]
\textbf{Version A (a French speaker's English):} “The citizens defended their rights. The government realized a survey. Many young people abandoned the struggle, but the association continued to protect the victims.”

\textbf{Version B (natural English):} “The citizens \cle{stood up for} their rights. The government \cle{carried out} a survey. Many young people \cle{gave up} the struggle, but the association \cle{went on} protecting the victims.”
\end{textebox}

Les deux versions sont grammaticalement correctes. Pourtant, un anglophone trouvera la version A un peu “livresque” et la version B beaucoup plus naturelle. Pourquoi ? Parce que l'anglais de tous les jours — conversations, presse populaire, discours — préfère les \cle{phrasal verbs}, alors que le français préfère les verbes uniques, souvent d'origine latine. Ironie de l'histoire : l'anglais possède aussi ces verbes latins (\eng{to defend}, \eng{to abandon}, \eng{to realise}), mais il les réserve au style soutenu ou juridique.

\begin{popfigure}
\begin{center}
\begin{tikzpicture}[
  boxfr/.style={rectangle, draw=poporange, fill=poporangeL, rounded corners=2pt, minimum width=4.2cm, minimum height=0.75cm, align=center, font=\small},
  boxen/.style={rectangle, draw=popblue, fill=popblueL, rounded corners=2pt, minimum width=4.2cm, minimum height=0.75cm, align=center, font=\small},
  arr/.style={-{Stealth[length=2.5mm]}, thick, popdark}
]
\node[boxfr, font=\small\bfseries] (t1) at (0,0) {Français (verbe unique)};
\node[boxen, font=\small\bfseries] (t2) at (6.5,0) {Anglais (phrasal verb)};
\node[boxfr] (f1) at (0,-1.0) {abandonner};
\node[boxen] (e1) at (6.5,-1.0) {to give up};
\node[boxfr] (f2) at (0,-1.9) {défendre};
\node[boxen] (e2) at (6.5,-1.9) {to stand up for};
\node[boxfr] (f3) at (0,-2.8) {réaliser (un projet)};
\node[boxen] (e3) at (6.5,-2.8) {to carry out};
\node[boxfr] (f4) at (0,-3.7) {chercher};
\node[boxen] (e4) at (6.5,-3.7) {to look for};
\node[boxfr] (f5) at (0,-4.6) {s'occuper de};
\node[boxen] (e5) at (6.5,-4.6) {to look after};
\draw[arr] (f1) -- (e1);
\draw[arr] (f2) -- (e2);
\draw[arr] (f3) -- (e3);
\draw[arr] (f4) -- (e4);
\draw[arr] (f5) -- (e5);
\end{tikzpicture}
\end{center}
\legende{Un verbe français = souvent un phrasal verb anglais}
\end{popfigure}

\courssub{Tableau récapitulatif : verbes français et phrasal verbs anglais}

\begin{center}
\begin{tabularx}{\linewidth}{|X|X|X|}
\hline
\rowcolor{popblueL} \textbf{Français} & \textbf{English (phrasal verb)} & \textbf{Exemple} \\
\hline
abandonner, renoncer à & to give up & \eng{He gave up smoking.} « Il a arrêté de fumer. » \\
\hline
défendre (une cause) & to stand up for & \eng{Stand up for your rights!} « Défends tes droits ! » \\
\hline
réaliser, mener (une action) & to carry out & \eng{They carried out a campaign.} « Ils ont mené une campagne. » \\
\hline
chercher & to look for & \eng{She is looking for a job.} « Elle cherche un emploi. » \\
\hline
s'occuper de & to look after & \eng{He looks after his grandmother.} « Il s'occupe de sa grand-mère. » \\
\hline
dénoncer, protester contre & to speak out against & \eng{They spoke out against corruption.} « Ils ont dénoncé la corruption. » \\
\hline
continuer & to go on / to carry on & \eng{Go on fighting!} « Continue à te battre ! » \\
\hline
découvrir, apprendre & to find out & \eng{We found out the truth.} « Nous avons découvert la vérité. » \\
\hline
participer à & to take part in & \eng{She took part in the march.} « Elle a participé à la marche. » \\
\hline
tolérer, supporter & to put up with & \eng{We cannot put up with injustice.} « Nous ne pouvons pas tolérer l'injustice. » \\
\hline
\end{tabularx}
\end{center}

\begin{attention}[Le français dit « défendre ses droits » en un mot]
Le français exprime en \textbf{un seul verbe} ce que l'anglais courant dit en \textbf{deux ou trois mots} : \emph{défendre ses droits} = \eng{stand up for one's rights}. À l'oral et dans l'écriture naturelle, ne cherche pas un verbe unique savant (\eng{to uphold}, \eng{to advocate}) : c'est le phrasal verb qui sonne juste. À l'inverse, dans un texte juridique très formel, le verbe latin (\eng{to defend}, \eng{to protect}) reste possible. Retiens l'échelle : \textbf{style courant = phrasal verb ; style soutenu = verbe unique}.
\end{attention}

\courssub{Les cinq pièges des francophones}

\textbf{Piège 1 : oublier la particule.}

En français, « abandonner » se suffit à lui-même ; le francophone traduit donc mot à mot et oublie la particule. Or sans elle, le sens change ou disparaît.

\begin{exemplebox}[La particule change tout]
\eng{He gave up smoking.} « Il a arrêté de fumer. » (et non \eng{He gave smoking}, qui ne veut rien dire)\par
\eng{She gave her seat to an old lady.} « Elle a cédé sa place à une dame âgée. » (ici \eng{give} seul = « donner » : autre sens !)\par
\eng{They carried out a survey in Yaoundé.} « Ils ont réalisé une enquête à Yaoundé. » (et non \eng{They carried a survey})
\end{exemplebox}

\textbf{Piège 2 : mauvais placement du pronom.}

Règle absolue : avec un phrasal verb \textbf{séparable}, un \textbf{pronom} complément se place \textbf{obligatoirement entre le verbe et la particule}. Le francophone, habitué à « je le réalise », veut coller le pronom après tout le groupe.

\begin{exemplebox}[Le pronom au milieu]
\eng{They carried it out successfully.} « Ils l'ont réalisé avec succès. » (jamais \eng{they carried out it})\par
\eng{The problem? We must look into it.} « Le problème ? Nous devons l'examiner. » (\eng{look into} est inséparable : le pronom va après)\par
\eng{She looks after them every day.} « Elle s'occupe d'eux tous les jours. » (\eng{look after} est inséparable)
\end{exemplebox}

\begin{propriete}[Séparable ou inséparable ? Le test du pronom]
\begin{itemize}
\item \textbf{Phrasal verb séparable} (souvent transitif, particule = adverbe) : \eng{carry out}, \eng{give up}, \eng{turn down}, \eng{put off}. Nom après ou au milieu : \eng{carry out the plan} ou \eng{carry the plan out}. Pronom \textbf{toujours au milieu} : \eng{carry it out}.
\item \textbf{Phrasal verb inséparable} (particule = préposition) : \eng{look after}, \eng{look for}, \eng{stand up for}, \eng{speak out against}. Nom et pronom \textbf{toujours après} : \eng{look after the children}, \eng{look after them}.
\end{itemize}
\end{propriete}

\textbf{Piège 3 : calquer le « pour » français.}

En français, le but s'exprime par « pour + infinitif » : \emph{il a adhéré pour aider}. Le francophone traduit donc \eng{for to help}. Faux ! En anglais, \eng{for} + verbe au gérondif exprime l'\textbf{usage d'une chose} ; le but d'une \textbf{personne} s'exprime par l'infinitif seul (ou \eng{in order to}, \eng{so as to}).

\begin{exemplebox}[But : pas de « for to »]
FAUX : \eng{He joined for to help.}\par
JUSTE : \eng{He joined to help.} « Il a adhéré pour aider. »\par
JUSTE : \eng{She went to Bamenda in order to register to vote.} « Elle est allée à Bamenda pour s'inscrire sur les listes électorales. »\par
Attention à la nuance : \eng{A vote is a tool for changing society.} « Le vote est un outil pour changer la société. » (ici \eng{for + -ing} décrit la \textbf{fonction de l'outil}, pas le but d'une personne : c'est correct.)
\end{exemplebox}

\textbf{Piège 4 : inventer des phrasal verbs qui n'existent pas.}

Le francophone combine librement verbe + particule, comme il le ferait avec un préfixe en français. Résultat : des formes imaginaires ou des confusions de sens. \eng{Stand up against} existe (résister physiquement, s'opposer avec courage), mais il n'est \textbf{pas interchangeable} avec \eng{speak out against} (dénoncer publiquement, par la parole).

\begin{exemplebox}[Des verbes proches mais pas identiques]
\eng{Students spoke out against the new fees.} « Les étudiants ont dénoncé publiquement les nouveaux frais. » (paroles, pétitions, médias)\par
\eng{The village stood up against the illegal land sale.} « Le village s'est dressé contre la vente illégale des terres. » (opposition, résistance)\par
FAUX : \eng{They gave up against the decision.} — ce verbe n'existe pas ; dire \eng{They protested against the decision} ou \eng{They spoke out against it}.
\end{exemplebox}

\textbf{Piège 5 : traduire la particule littéralement.}

La particule d'un phrasal verb a souvent \textbf{perdu son sens spatial}. \eng{Give up} ne contient aucune idée de « haut », \eng{carry out} aucune idée de « dehors », \eng{look after} aucune idée de « après ». Chercher à traduire la particule séparément mène à des absurdités.

\begin{exemplebox}[Ne traduis jamais la particule seule]
\eng{Don't give up!} « N'abandonne pas ! » (et non « ne donne pas en haut »)\par
\eng{She looks after the children.} « Elle s'occupe des enfants. » (et non « elle regarde après les enfants »)\par
FAUX : \eng{She looks the children after.} — double erreur : calque du français « elle garde les enfants » avec particule rejetée à la fin, alors que \eng{look after} est inséparable.\par
JUSTE : \eng{She looks after the children.} / \eng{She looks after them.}
\end{exemplebox}

\begin{savaistu}[Une lettre change tout le sens]
\eng{To look after} signifie « s'occuper de » et \eng{to look for} signifie « chercher » : seule la particule diffère ! Compare : \eng{I'm looking for my keys.} « Je cherche mes clés. » et \eng{I'm looking after my neighbour's children.} « Je garde les enfants de ma voisine. » Autre paire célèbre : \eng{to look up} (chercher dans un dictionnaire, lever les yeux) et \eng{to look down on} (mépriser). Moralité : un phrasal verb s'apprend \textbf{en bloc}, comme un seul mot de vocabulaire, jamais verbe et particule séparément.
\end{savaistu}

\courssub{Stratégie de vérification : les deux questions}

\begin{methode}[Avant d'écrire ou de parler, pose-toi deux questions]
\begin{enumerate}
\item \textbf{Le verbe est-il séparable ?} Si oui et si le complément est un \textbf{pronom}, place-le au milieu : \eng{carry it out}, \eng{give it up}. Si tu hésites sur la séparabilité, utilise un \textbf{nom} après le verbe complet (\eng{carry out the project}) : c'est toujours correct.
\item \textbf{Le sujet du but est-il le même ?} Si le sujet de l'action et celui du but sont identiques, utilise l'infinitif : \eng{He joined to help.} Si les sujets sont différents, utilise \eng{so that + sujet + can/will} : \eng{He spoke slowly so that everyone could understand.} « Il parlait lentement pour que tout le monde puisse comprendre. »
\end{enumerate}
\textbf{Règle de sécurité anti-calque :} avant d'écrire, demande-toi si l'anglais utilise un phrasal verb \textbf{connu et vérifié}. Si tu hésites, n'invente rien : utilise le verbe simple (\eng{defend}, \eng{help}, \eng{protect}, \eng{continue}) plutôt qu'une combinaison imaginaire. Un anglais simple et correct vaut mieux qu'un phrasal verb fantaisiste.
\end{methode}

\begin{exoresolu}[Corrige les erreurs du candidat]
Un élève de Terminale écrit : \eng{(1) The association carried out it last year. (2) Many youths gave the struggle up against injustice for to defend their community. (3) She looks the children after to help their mother.} Corrige chaque phrase et justifie.
\tcblower
\textbf{(1)} \eng{The association carried it out last year.} — \eng{carry out} est séparable ; le pronom \eng{it} se place obligatoirement entre le verbe et la particule.\par
\textbf{(2)} \eng{Many youths gave up the struggle against injustice to defend their community.} « Beaucoup de jeunes ont abandonné la lutte contre l'injustice pour défendre leur communauté. » — Deux corrections : (a) avec un nom long (\eng{the struggle against injustice}), la particule \eng{up} se place mieux juste après le verbe ; (b) le but s'exprime par l'infinitif \eng{to defend}, jamais par \eng{for to}.\par
\textbf{(3)} \eng{She looks after the children to help their mother.} « Elle s'occupe des enfants pour aider leur mère. » — \eng{look after} est \textbf{inséparable} : le complément (nom ou pronom) suit toujours la particule.
\end{exoresolu}

\begin{exercice}[À toi de jouer]
Traduis en anglais en utilisant un phrasal verb et, quand il y a un but, la bonne expression du but :\par
1. Les citoyens ont défendu leurs droits pour protéger la démocratie.\par
2. Elle les a cherchés partout. (ses papiers d'identité)\par
3. Nous devons dénoncer la corruption pour que les choses changent.\par
4. Il a arrêté de fumer pour rester en bonne santé.\par
5. Ne traduis jamais un phrasal verb mot à mot !
\end{exercice}

\begin{corrige}[Exercice « À toi de jouer »]
1. \eng{The citizens stood up for their rights to protect democracy.} — \eng{stand up for} est inséparable ; but = infinitif \eng{to protect} (même sujet).\par
2. \eng{She looked for them everywhere.} — \eng{look for} est inséparable : le pronom \eng{them} va après la particule.\par
3. \eng{We must speak out against corruption so that things can change.} — sujets différents (\eng{we} / \eng{things}) : \eng{so that + can}.\par
4. \eng{He gave up smoking to stay healthy.} — \eng{give up + -ing} ; but = infinitif.\par
5. \eng{Never translate a phrasal verb word for word!} — \eng{word for word} = « mot à mot ».
\end{corrige}

\begin{aretenir}[L'essentiel de la section]
\begin{itemize}
\item Le français utilise un verbe unique ; l'anglais courant préfère un phrasal verb : \emph{abandonner} = \eng{give up}, \emph{défendre} = \eng{stand up for}, \emph{réaliser} = \eng{carry out}.
\item Un phrasal verb s'apprend \textbf{en bloc} : ne traduis jamais la particule littéralement.
\item Pronom + verbe séparable : le pronom va \textbf{au milieu} (\eng{carry it out}). Verbe inséparable : tout va après (\eng{look after them}).
\item Le but d'une personne = infinitif (\eng{to help}) ou \eng{in order to} / \eng{so that} ; jamais \eng{for to}.
\item En cas de doute, n'invente pas : utilise le verbe simple (\eng{defend}, \eng{protect}).
\end{itemize}
\end{aretenir}

\coursec{Entraînement guidé : de la phrase au paragraphe}

Après avoir comparé le français et l'anglais et repéré les pièges les plus fréquents des francophones (particule oubliée, pronom mal placé, \eng{for to} inventé), il est temps de passer à la pratique. Cette dernière section te conduit pas à pas, de la particule isolée jusqu'au paragraphe complet : c'est exactement le cheminement que l'examinateur attend de toi au Baccalauréat. Chaque exercice est précédé d'un rappel de méthode et suivi d'un corrigé détaillé.

\begin{popfigure}
\begin{center}
\begin{tikzpicture}[node distance=0.35cm]
\node[draw=popblue, fill=popblueL, rounded corners, minimum width=2.5cm, minimum height=0.9cm, align=center, font=\small\bfseries] (a) {1. Particule\\ \eng{up, out, for...}};
\node[draw=poporange, fill=poporangeL, rounded corners, minimum width=2.5cm, minimum height=0.9cm, align=center, font=\small\bfseries, right=of a] (b) {2. Pronom\\ \eng{carry it out}};
\node[draw=popgreen, fill=popgreenL, rounded corners, minimum width=2.5cm, minimum height=0.9cm, align=center, font=\small\bfseries, right=of b] (c) {3. But\\ \eng{to, in order to}};
\node[draw=poppurple, fill=poppurpleL, rounded corners, minimum width=2.5cm, minimum height=0.9cm, align=center, font=\small\bfseries, right=of c] (d) {4. Combinaison\\ 2 phrases $\rightarrow$ 1};
\node[draw=popgold, fill=popgoldL, rounded corners, minimum width=2.5cm, minimum height=0.9cm, align=center, font=\small\bfseries, right=of d] (e) {5. Paragraphe\\ production libre};
\draw[-{Stealth}, thick, popdark] (a) -- (b);
\draw[-{Stealth}, thick, popdark] (b) -- (c);
\draw[-{Stealth}, thick, popdark] (c) -- (d);
\draw[-{Stealth}, thick, popdark] (d) -- (e);
\end{tikzpicture}
\end{center}
\legende{Frise de progression : cinq étapes, de la particule isolée au paragraphe argumenté.}
\end{popfigure}

\courssub{Exercice 1 : compléter avec la particule correcte}

\begin{methode}[Comment choisir la bonne particule ?]
\begin{enumerate}
\item Lis toute la phrase et repère le sens voulu (défendre ? accomplir ? rechercher ? s'occuper de ?).
\item Identifie le verbe de base : certains verbes n'acceptent qu'une particule précise (\eng{stand up \textbf{for}}, \eng{look \textbf{after}}).
\item Vérifie dans ton tableau de phrasal verbs appris dans les sections précédentes.
\item Relis la phrase complète à voix haute : si le sens est bizarre, la particule est fausse.
\end{enumerate}
\end{methode}

\begin{exercice}[Exercice 1 — Fill in the blanks]
Complete each sentence with the correct particle: \textbf{up, out, for, after, about, forward}.
\begin{enumerate}
\item Citizens must stand \ldots\ldots{} \ldots\ldots{} their rights.
\item The committee carried \ldots\ldots{} a survey in Yaoundé schools.
\item Many young people are looking \ldots\ldots{} a better future abroad.
\item The NGO looks \ldots\ldots{} orphans in the North-West region.
\item The minister spoke \ldots\ldots{} the new law on television.
\item We look \ldots\ldots{} to hearing from the human rights commission.
\end{enumerate}
\end{exercice}

\begin{corrige}[Exercice 1]
\begin{enumerate}
\item \eng{stand \textbf{up for} their rights} — défendre leurs droits (phrasal verb à deux particules).
\item \eng{carried \textbf{out} a survey} — mené une enquête.
\item \eng{looking \textbf{for} a better future} — à la recherche d'un avenir meilleur.
\item \eng{looks \textbf{after} orphans} — s'occupe des orphelins.
\item \eng{spoke \textbf{about} the new law} — a parlé de la nouvelle loi.
\item \eng{look \textbf{forward to} hearing} — nous attendons avec plaisir (attention : \eng{to} est ici une préposition, suivie de la forme en \eng{-ing}).
\end{enumerate}
\end{corrige}

\courssub{Exercice 2 : replacer le pronom au bon endroit}

Rappel de la règle d'or : avec un phrasal verb \cle{séparable}, un \textbf{nom} peut se placer avant ou après la particule (\eng{carry out the plan} = \eng{carry the plan out}), mais un \textbf{pronom} (\eng{it, them, him}) se place \textbf{obligatoirement entre le verbe et la particule}.

\begin{exercice}[Exercice 2 — Put the pronoun in the right place]
Rewrite each sentence, replacing the underlined noun with a pronoun.
\begin{enumerate}
\item The government carried out \emph{the reforms}. $\rightarrow$ The government \ldots\ldots{}
\item She turned down \emph{the offer}. $\rightarrow$ She \ldots\ldots{}
\item They looked after \emph{the refugees}. $\rightarrow$ They \ldots\ldots{}
\item We must give up \emph{violence}. $\rightarrow$ We must \ldots\ldots{}
\end{enumerate}
\end{exercice}

\begin{corrige}[Exercice 2]
\begin{enumerate}
\item \eng{The government carried \textbf{them} out.} (réformes = pluriel)
\item \eng{She turned \textbf{it} down.} — Elle l'a refusée.
\item \eng{They looked after \textbf{them}.} — Attention : \eng{look after} est \textbf{inséparable} ! Le pronom reste après la particule.
\item \eng{We must give \textbf{it} up.} — Nous devons y renoncer.
\end{enumerate}
La phrase 3 est un piège classique : tous les phrasal verbs ne sont pas séparables. Quand la particule est une vraie préposition (\eng{look after, look for, stand up for}), le pronom suit toujours la particule.
\end{corrige}

\courssub{Exercice 3 : transformer « pour + infinitif » du français}

Le français utilise « pour + infinitif » ; l'anglais offre trois équivalents : \eng{to}, \eng{in order to}, \eng{so as to} + base verbale.

\begin{center}
\begin{tabularx}{\linewidth}{|X|X|X|}
\hline
\rowcolor{popblueL} Français & English & Remarque \\
\hline
Il travaille pour réussir. & \eng{He works \textbf{to} succeed.} & forme la plus courante \\
\hline
Elle manifeste afin de changer les choses. & \eng{She protests \textbf{in order to} change things.} & forme soutenue, idéale à l'écrit \\
\hline
Il part tôt pour ne pas être en retard. & \eng{He leaves early \textbf{so as not to} be late.} & négation : \eng{so as \textbf{not} to} \\
\hline
\end{tabularx}
\end{center}

\begin{exercice}[Exercice 3 — Translate into English]
Translate using \eng{to}, \eng{in order to} or \eng{so as to}.
\begin{enumerate}
\item Les élèves apprennent l'anglais pour comprendre le monde.
\item Elle a créé une association afin d'aider les veuves.
\item Il parle doucement pour ne pas réveiller le bébé.
\item Nous votons pour choisir nos dirigeants.
\end{enumerate}
\end{exercice}

\begin{corrige}[Exercice 3]
\begin{enumerate}
\item \eng{Students learn English \textbf{to} understand the world.}
\item \eng{She created an association \textbf{in order to} help widows.}
\item \eng{He speaks softly \textbf{so as not to} wake the baby.}
\item \eng{We vote \textbf{to} choose our leaders.}
\end{enumerate}
\end{corrige}

\courssub{Exercice 4 : combiner deux phrases en une seule}

\begin{methode}[Transformer deux phrases en une avec une expression de but]
\begin{enumerate}
\item \textbf{Vérifie que le sujet est le même} dans les deux phrases.
\item \textbf{Supprime le deuxième sujet} et le verbe conjugué qui le suit (\eng{she wanted to} disparaît).
\item \textbf{Relie avec \eng{to} / \eng{in order to} / \eng{so as to} + infinitif}, en intégrant si possible un phrasal verb.
\end{enumerate}
Modèle : \eng{She joined the NGO. She wanted to defend children.} $\rightarrow$ \eng{She joined the NGO \textbf{to stand up for} children.}
\end{methode}

\begin{exoresolu}[Exercice 4 — Combining sentences]
Énoncé : combine chaque paire de phrases en une seule phrase avec une expression de but.
\begin{enumerate}
\item They protest in the streets. They want to provoke change.
\item He studies law. He wants to defend the oppressed.
\end{enumerate}
\tcblower
Solution détaillée :
\begin{enumerate}
\item Le sujet est identique (\eng{they}). On supprime \eng{they want to} et on remplace \eng{provoke} par le phrasal verb \eng{bring about} : \eng{They protest \textbf{in order to bring about} change.} = « Ils manifestent afin de provoquer le changement. »
\item Sujet identique (\eng{he}). On relie avec \eng{to} et on utilise \eng{stand up for} : \eng{He studies law \textbf{to stand up for} the oppressed.} = « Il étudie le droit pour défendre les opprimés. »
\end{enumerate}
\end{exoresolu}

\begin{attention}[Un seul sujet avec « to » !]
Dans cette transformation, ne garde \textbf{jamais deux sujets} avec \eng{to} :\\
FAUX : \eng{She joined the NGO to she helps children.}\\
Si le sujet \textbf{change} entre les deux phrases, utilise \eng{so that} + sujet + verbe conjugué :\\
\eng{She joined the NGO \textbf{so that} the children could be protected.} = « Elle a rejoint l'ONG pour que les enfants soient protégés. »
\end{attention}

\courssub{Exercice 5 : corriger des phrases fautives typiques de francophones}

\begin{exercice}[Exercice 5 — Correct the mistakes]
Each sentence contains one typical French-speaker mistake. Correct it.
\begin{enumerate}
\item \eng{He went to Douala for to find a job.}
\item \eng{She carried out it successfully.}
\item \eng{They fight for defend their land.}
\item \eng{We must to stand up for justice.}
\item \eng{He studies hard for not fail the exam.}
\end{enumerate}
\end{exercice}

\begin{corrige}[Exercice 5]
\begin{enumerate}
\item \eng{He went to Douala \textbf{to} find a job.} — jamais \eng{for to}.
\item \eng{She carried \textbf{it out} successfully.} — le pronom se place entre le verbe et la particule.
\item \eng{They fight \textbf{to defend} their land.} — après \eng{for}, il faudrait un nom ; le but exige \eng{to} + infinitif.
\item \eng{We must stand up for justice.} — pas de \eng{to} après un modal (\eng{must}).
\item \eng{He studies hard \textbf{so as not to} fail the exam.} — but négatif : \eng{so as not to} (ou \eng{in order not to}).
\end{enumerate}
\end{corrige}

\courssub{Production finale : ton engagement citoyen en cinq phrases}

\begin{propriete}[Un bon paragraphe au Bac]
Un bon paragraphe au Baccalauréat \textbf{alterne} les phrasal verbs (qui donnent un ton naturel et idiomatique) et les expressions de but (qui rendent l'argumentation claire et logique). Vise 5 à 7 phrases, un connecteur par phrase (\eng{first, moreover, finally}), et aucune répétition de structure.
\end{propriete}

\begin{exemplebox}[Modèle de production]
\eng{Last year, I decided to join the environmental club of my school \textbf{in order to bring about} change in my neighbourhood. First, we \textbf{carried out} a cleaning campaign \textbf{to} protect the river near the market. Moreover, I always \textbf{stand up for} younger students when they are bullied. I also want to study law \textbf{so as to} defend people who cannot pay a lawyer. Finally, I believe every citizen should \textbf{give up} violence and \textbf{look forward to} a peaceful Cameroon.}\par
\smallskip
Traduction : « L'an dernier, j'ai décidé de rejoindre le club d'environnement de mon lycée afin de provoquer du changement dans mon quartier. D'abord, nous avons mené une campagne de nettoyage pour protéger la rivière près du marché. De plus, je défends toujours les élèves plus jeunes quand ils sont harcelés. Je veux aussi étudier le droit afin de défendre les gens qui ne peuvent pas payer un avocat. Enfin, je crois que chaque citoyen devrait renoncer à la violence et attendre avec espoir un Cameroun pacifique. »
\end{exemplebox}

\begin{experience}[À toi d'écrire !]
Rédige \textbf{5 phrases} sur un engagement citoyen personnel (réel ou imaginé : club, association, quartier, école). Contraintes :
\begin{itemize}
\item au moins \textbf{3 phrasal verbs} différents (\eng{stand up for, carry out, look after, bring about, give up, speak out...}) ;
\item au moins \textbf{2 expressions de but} (\eng{to, in order to, so as to, so that}) ;
\item au moins \textbf{2 connecteurs} (\eng{first, moreover, finally}).
\end{itemize}
Échange ensuite ta production avec ton voisin : il doit souligner tes phrasal verbs et encadrer tes expressions de but, puis vérifier qu'il n'y a ni double sujet après \eng{to}, ni \eng{for to}.
\end{experience}

\begin{aretenir}[L'essentiel de l'entraînement]
\begin{itemize}
\item La particule se choisit par le \textbf{sens} : \eng{stand up \textbf{for}} (défendre), \eng{carry \textbf{out}} (accomplir), \eng{look \textbf{after}} (s'occuper de).
\item Pronom + phrasal verb séparable : \textbf{toujours au milieu} (\eng{carry \textbf{it} out}).
\item But : \eng{to / in order to / so as to} + infinitif ; négation : \eng{so as \textbf{not} to}.
\item Deux phrases, un seul sujet $\rightarrow$ \eng{to} ; deux sujets différents $\rightarrow$ \eng{so that}.
\item Au Bac : alterne phrasal verbs et expressions de but pour un paragraphe naturel et clair.
\end{itemize}
\end{aretenir}

\newpage

\coursec{Activité d'intégration}

\courssub{Objectifs de l'activité}

À l'issue de cette activité d'intégration, tu seras capable de :

\begin{itemize}
\item réinvestir les \cle{phrasal verbs} étudiés dans la leçon (\eng{stand up for, look after, carry out, bring about, speak out against, put up with, call for, give up, take part in, clean up}) dans une production écrite personnelle et argumentée ;
\item exprimer le \cle{but} avec les structures \eng{to + base verbale}, \eng{in order to}, \eng{so as to} et \eng{so that + proposition}, en choisissant la structure adaptée au contexte ;
\item rédiger un texte collectif court, formel et engagé : une \emph{Citizens' Charter} (charte du citoyen) ;
\item présenter oralement un texte en deux minutes, avec une prononciation claire et un ton convaincant ;
\item évaluer la production des autres groupes selon des critères précis (correction grammaticale, richesse du lexique, force de conviction).
\end{itemize}

\courssub{Situation}

Ton école, \emph{Government Bilingual High School Mendong}, à Yaoundé, connaît quelques difficultés : déchets dans la cour, brimades (\eng{bullying}) contre les élèves de Seconde, retards fréquents, et peu de participation aux activités citoyennes. Le proviseur (\eng{the principal}) a lancé un concours : chaque classe de Terminale doit proposer une \emph{Citizens' Charter} — une charte du citoyen engagé — pour améliorer la vie de l'école. La meilleure charte sera affichée à l'entrée de l'établissement et lue lors de la cérémonie de fin d'année.

Voici le message officiel affiché au tableau d'information de l'école :

\begin{textebox}[Notice on the school board — GBHS Mendong, Yaoundé]
\textbf{CITIZENS' CHARTER COMPETITION}

Dear students,

Our school belongs to all of us. Too often, we put up with problems instead of solving them: litter in the courtyard, bullying of younger students, and silence when we should speak out against injustice.

The Students' Council invites every Form Five and Upper Sixth class to write a \textbf{Citizens' Charter}: six clear commitments to bring about positive change in our school and neighbourhood.

Each commitment must say \textbf{what} you will do and \textbf{why} you will do it.

The best charter will be displayed at the main entrance and read aloud at the end-of-year ceremony. Stand up for your school — take part in the competition!

Deadline: Friday, 4 p.m. — The Students' Council
\end{textebox}

\textbf{Glossaire :} \emph{litter} (n.m.) : déchets ; \emph{to display} (v.) : afficher ; \emph{aloud} (adv.) : à voix haute ; \emph{deadline} (n.f.) : date limite.

\courssub{Consignes}

Travail en \textbf{groupes de 4}. Durée totale : environ 50 minutes. Suis les étapes dans l'ordre.

\begin{enumerate}
\item \textbf{Compréhension du document (5 min).} Lis la notice ci-dessus. Releve : (a) trois problèmes mentionnés dans l'école ; (b) deux phrasal verbs utilisés dans le texte ; (c) ce que la charte gagnante deviendra.
\item \textbf{Remue-méninges lexical (5 min).} Dans ton groupe, listez au moins huit problèmes de votre école ou de votre quartier (ex. : \eng{noise, dirty toilets, exam cheating, water shortage, discrimination}). Pour chacun, associez un phrasal verb de la leçon qui permet d'agir : \eng{to clean up, to speak out against, to stand up for, to call for, to look after, to give up, to carry out, to bring about}.
\item \textbf{Rédaction de la charte (20 min).} Rédigez \textbf{six engagements} numérotés, commençant chacun par \eng{We will...} ou \eng{We promise to...}. Chaque engagement doit contenir \textbf{obligatoirement} : (a) au moins un phrasal verb correctement construit ; (b) une expression de but (\eng{to}, \eng{in order to}, \eng{so as to} ou \eng{so that}). Exemple de structure : \eng{We will + phrasal verb + complément + expression of purpose.}
\item \textbf{Relecture collective (5 min).} Vérifiez chaque phrase : le phrasal verb est-il bien formé (verbe + particule) ? La particule est-elle la bonne ? Le but est-il exprimé avec la bonne structure ? Corrigez les calques du français.
\item \textbf{Présentation orale (2 min par groupe).} Deux membres du groupe lisent la charte à voix haute, avec conviction. Les deux autres répondent aux questions de la classe.
\item \textbf{Vote et affichage (5 min).} La classe vote pour la charte la plus convaincante selon trois critères : correction des phrasal verbs, variété des expressions de but, force d'engagement. La charte gagnante est affichée au tableau.
\end{enumerate}

\begin{attention}[Critères du vote — sois exigeant]
Un engagement ne peut pas être validé si : le phrasal verb est mal formé (\eng{*stand up against bullied students} au lieu de \eng{stand up for bullied students}) ; le but est exprimé par \eng{*for to} ou \eng{*for + verbe en -ing} maladroit ; la phrase est un calque du français (\eng{*we will make down the rubbish}).
\end{attention}

\courssub{Ressources à mobiliser}

\begin{aretenir}[Boîte à outils de la charte]
\textbf{Phrasal verbs de la citoyenneté :} \eng{stand up for} (défendre), \eng{speak out against} (dénoncer), \eng{look after} (prendre soin de), \eng{carry out} (mener, réaliser), \eng{bring about} (provoquer, susciter), \eng{call for} (réclamer), \eng{give up} (abandonner), \eng{put up with} (tolérer), \eng{clean up} (nettoyer), \eng{take part in} (participer à).

\textbf{Expressions de but :}
\begin{itemize}
\item \eng{to + BV} : \eng{We will clean up the courtyard \textbf{to protect} our health.}
\item \eng{in order to / so as to + BV} (plus formel) : \eng{... \textbf{in order to bring about} change.}
\item \eng{so that + sujet + will/can + BV} (quand le sujet change) : \eng{... \textbf{so that everyone can feel} safe.}
\end{itemize}

\textbf{Formulations d'engagement :} \eng{We will... / We promise to... / We commit ourselves to + V-ing / We refuse to put up with...}
\end{aretenir}

\begin{propriete}[Rappel de la règle d'or]
Après \eng{to}, \eng{in order to} et \eng{so as to}, on utilise toujours la \cle{base verbale} : \eng{in order to \textbf{bring about}} (jamais \eng{*in order to bringing about}). Après \eng{so that}, on construit une \cle{proposition complète} avec son propre sujet : \eng{so that \textbf{our school becomes} cleaner}.
\end{propriete}

\courssub{Proposition de production et corrigé commenté}

\begin{exoresolu}[Modèle de Citizens' Charter — Groupe 3, Tle SH2]
\textbf{Énoncé.} Rédige, avec ton groupe, une charte de six engagements respectant les contraintes : un phrasal verb et une expression de but par engagement.

\tcblower

\textbf{Production modèle :}

\emph{The Citizens' Charter of Upper Sixth A2 — GBHS Mendong}

\emph{1. We will stand up for bullied students in order to bring about a safer school for everyone.}

\emph{2. We will carry out a clean-up campaign every Friday afternoon to look after our environment.}

\emph{3. We promise to speak out against exam cheating so that our certificates keep their value.}

\emph{4. We will call for a water tank in the courtyard so as to help students wash their hands.}

\emph{5. We refuse to put up with litter in the classrooms; we will give up plastic bags to keep our school clean.}

\emph{6. We commit ourselves to taking part in community work in Mendong neighbourhood in order to show that young citizens can bring about real change.}

\textbf{Commentaire des choix de langue :}
\begin{itemize}
\item \textbf{Engagement 1} : \eng{stand up for} = « défendre » (attention : \eng{for}, pas \eng{against}) ; le but est exprimé par \eng{in order to + BV}, registre formel adapté à une charte.
\item \textbf{Engagement 2} : \eng{carry out} = « mener à bien » ; le complément de temps \eng{every Friday afternoon} est placé avant l'expression de but pour garder une phrase claire ; \eng{to look after} enchaîne naturellement un second phrasal verb.
\item \textbf{Engagement 3} : \eng{so that + proposition} est obligatoire ici car le sujet change (\eng{our certificates}) ; \eng{keep} est au présent après \eng{so that} dans un contexte de vérité générale.
\item \textbf{Engagement 4} : \eng{call for} = « réclamer » ; \eng{so as to} varie la formulation et évite la répétition de \eng{in order to}.
\item \textbf{Engagement 5} : \eng{put up with} = « tolérer » ; la négation \eng{we refuse to put up with} montre qu'un phrasal verb peut suivre un autre verbe ; \eng{give up} + nom (\eng{plastic bags}) est correct, de même que \eng{give up + V-ing}.
\item \textbf{Engagement 6} : après \eng{commit ourselves to}, le verbe est en \emph{-ing} (\eng{taking part}) car ce \eng{to} est une préposition, pas une particule d'infinitif — piège classique évité.
\end{itemize}
\end{exoresolu}

\courssub{Bilan de l'activité}

Cette activité t'a permis de passer de la \cle{reconnaissance} des phrasal verbs (dans la notice de l'école) à leur \cle{production contrôlée} puis \cle{libre} dans un texte authentique et utile. Retiens les points suivants :

\begin{itemize}
\item Un phrasal verb est un couple inséparable \textbf{verbe + particule} dont le sens est souvent idiomatique : on ne le traduit pas mot à mot (\eng{stand up for} ne veut pas dire « se lever pour »).
\item Le choix de la particule change tout le sens : \eng{stand up \textbf{for}} (défendre) n'a rien à voir avec \eng{put up \textbf{with}} (tolérer).
\item Pour exprimer le but, trois outils : \eng{to + BV} (simple), \eng{in order to / so as to + BV} (formel), \eng{so that + proposition} (quand le sujet change).
\item Méfie-toi du \eng{to} prépositionnel suivi de \emph{-ing} : \eng{commit oneself to \textbf{taking} part}, \eng{look forward to \textbf{seeing}}.
\item Dans une charte, un discours ou une pétition — genres fréquents au chapitre \emph{Citizenship and Human Rights} — ces structures donnent à ton écrit un ton \cle{engagé, formel et convaincant}, exactement ce qu'attend l'épreuve de production écrite du baccalauréat.
\end{itemize}

\begin{methode}[Pour aller plus loin : auto-évaluation]
Relis ta charte et coche : (1) ai-je utilisé au moins quatre phrasal verbs différents ? (2) ai-je varié les expressions de but ? (3) chaque particule est-elle la bonne ? (4) ai-je évité les calques du français ? Si une réponse est « non », réécris l'engagement concerné avant l'affichage final.
\end{methode}

\newpage

\coursec{Méthodes et exercices résolus}

\coursec{Lesson 46 — Grammar: Review phrasal verbs in simple expressions of purpose}

\courssub{Observation : découvrir les phrasal verbs en contexte}

\begin{textebox}[Script d'écoute / Texte support — Citizenship in Action]
“Last year, the municipal council in Buea decided to \textbf{close down} the old community centre. Local residents did not \textbf{give in}. They \textbf{stood up for} their right to a public space and \textbf{carried on} organising peaceful sit-ins. Finally, the authorities \textbf{called off} the closure project. The centre was saved because citizens \textbf{spoke out against} the decision and \textbf{took part in} the consultation process.”
\end{textebox}

\begin{experience}[Activité de classe : Repérage et sens]
\begin{enumerate}
\item Écoutez ou lisez le texte. Soulignez les six \cle{phrasal verbs} (verbes à particule).
\item Associez chaque phrasal verb à son équivalent français parmi : \eng{abandonner / céder}, \eng{fermer}, \eng{défendre}, \eng{continuer}, \eng{annuler}, \eng{dénoncer publiquement}, \eng{participer à}.
\item Discutez en binôme : pourquoi le pronom \eng{it} ne peut-il pas être placé après \eng{called off} dans \eng{They called it off} ?
\end{enumerate}
\end{experience}

\courssub{Règle 1 : Formation et types de phrasal verbs}

\begin{definition}[Phrasal verb]
Un \cle{phrasal verb} est une unité lexicale formée d'un \textbf{verbe} + une \textbf{particule} (adverbe ou préposition) dont le sens global est souvent idiomatique (différent de la somme des parties). Ex. : \eng{give up} (abandonner) $\neq$ \eng{give} (donner) + \eng{up} (en haut).
\end{definition}

\begin{propriete}[Classification syntaxique — 4 types]
\begin{center}
\begin{tabularx}{\linewidth}{|X|X|X|X|}
\hline
\rowcolor{popblueL}
\textbf{Type} & \textbf{Structure} & \textbf{Place du COD (nom)} & \textbf{Place du COD (pronom)} \\
\hline
\textbf{Intransitif} & V + Particule & Pas de COD & Pas de COD \\
\hline
\textbf{Transitif séparable} & V + Particule + COD & V + COD + Part \textbf{OU} V + Part + COD & \textbf{V + Pron + Part} (obligatoire) \\
\hline
\textbf{Transitif inséparable} & V + Part + COD & V + Part + COD (seule place) & V + Part + Pron (seule place) \\
\hline
\textbf{À 3 éléments (V + Part + Prep)} & V + Part + Prep + COD & V + Part + Prep + COD & V + Part + Prep + Pron \\
\hline
\end{tabularx}
\end{center}

\end{propriete}

\begin{exemplebox}[Exemples par type]
\begin{itemize}
\item \textbf{Intransitif :} \eng{give in} (céder), \eng{turn up} (arriver/se présenter). \eng{The protesters refused to give in.}
\item \textbf{Séparable :} \eng{call off} (annuler), \eng{set up} (créer), \eng{put off} (reporter), \eng{clean up} (nettoyer). \eng{They called off the strike.} = \eng{They called the strike off.} $\rightarrow$ \eng{They called it off.}
\item \textbf{Inséparable :} \eng{look after} (s'occuper de), \eng{carry on} (continuer), \eng{stand for} (représenter/tolérer). \eng{She looks after the children.} (Jamais \eng{looks the children after}).
\item \textbf{À 3 éléments (inséparables) :} \eng{stand up for} (défendre), \eng{speak out against} (dénoncer), \eng{put up with} (tolérer), \eng{take part in} (participer à). \eng{We must stand up for our rights.}
\end{itemize}
\end{exemplebox}

\begin{methode}[Comment analyser un phrasal verb ? — Fiche méthode 1]
\begin{enumerate}
\item \textbf{Repérer la particule} (\eng{up, down, out, off, on, in, for, after, into, away, against}...).
\item \textbf{Tester la transitivité :} y a-t-il un complément d'objet (COD) possible ?
\item \textbf{Appliquer le test du pronom :} remplacez le COD par \eng{it/them}.
\begin{itemize}
\item Si le pronom va \textbf{au milieu} $\rightarrow$ \textbf{Séparable}.
\item Si le pronom reste \textbf{après} la particule/préposition $\rightarrow$ \textbf{Inséparable}.
\end{itemize}
\item \textbf{Identifier le sens contextuel} (polysémie fréquente : \eng{take up} = commencer / occuper / reprendre).
\item \textbf{Conjuguer le verbe seulement} : la particule est invariante. \eng{He \textbf{stood} up for his rights.}
\end{enumerate}
\end{methode}

\courssub{Règle 2 : Les phrasal verbs de la citoyenneté et des droits de l'homme}

\begin{propriete}[Lexique thématique obligatoire — Chapter 4]
\begin{center}
\begin{tabularx}{\linewidth}{|l|l|X|}
\hline
\rowcolor{popblueL}
\textbf{Phrasal Verb} & \textbf{Type} & \textbf{Sens français + Exemple} \\
\hline
\eng{stand up for} & 3 éléments & défendre \eng{-- Citizens must stand up for human rights.} \\
\eng{speak out against} & 3 éléments & dénoncer publiquement \eng{-- NGOs speak out against corruption.} \\
\eng{fight for} & Inséparable (prep) & se battre pour \eng{-- They fought for independence.} \\
\eng{take part in} & 3 éléments & participer à \eng{-- Youths should take part in elections.} \\
\eng{look after} & Inséparable & veiller sur / protéger \eng{-- The state must look after vulnerable people.} \\
\eng{give back} & Séparable & rendre / restituer / contribuer \eng{-- Volunteers give back to the community.} \\
\eng{set up} & Séparable & créer / fonder \eng{-- They set up an association.} \\
\eng{call off} & Séparable & annuler \eng{-- The demo was called off.} \\
\eng{put off} & Séparable & reporter \eng{-- The meeting was put off until Monday.} \\
\eng{put up with} & 3 éléments & tolérer / supporter \eng{-- We cannot put up with injustice.} \\
\eng{join in} & Intransitif & participer (se joindre à) \eng{-- Everyone joined in the cleanup.} \\
\eng{turn out} & Intransitif & se présenter / se révéler \eng{-- Many voters turned out.} \\
\hline
\end{tabularx}
\end{center}
\end{propriete}

\begin{savaistu}[Contexte camerounais : Bilingualism \& Civic Engagement]
Au Cameroun, le bilinguisme officiel (anglais/français) influence l'anglais camerounais (\eng{CamE}). Certains phrasal verbs y sont très productifs : \eng{“follow up”} (faire le suivi), \eng{“hand over”} (remettre officiellement), \eng{“sit in”} (faire un sit-in). Dans les provinces du Nord-Ouest et du Sud-Ouest (ex. Buea, Bamenda), l'anglais standard britannique reste la norme scolaire pour le Baccalauréat. Maîtriser ces phrasal verbs permet de rédiger des compositions sur \eng{civic responsibility}, \eng{good governance} ou \eng{human rights advocacy} avec la précision attendue.
\end{savaistu}

\courssub{Règle 3 : Exprimer le but (Expressions of Purpose)}

\begin{propriete}[Structures du but — Règle de choix]
Le choix dépend de l'\textbf{identité des sujets} des deux actions (action principale $\rightarrow$ but).
\begin{center}
\begin{tabularx}{\linewidth}{|X|X|X|}
\hline
\rowcolor{popblueL}
\textbf{Structure} & \textbf{Condition sujets} & \textbf{Exemple} \\
\hline
\textbf{to + BV} & \textbf{Même sujet} (standard) & \eng{She went to Yaoundé \textbf{to meet} the Minister.} \\
\textbf{in order to / so as to + BV} & \textbf{Même sujet} (formel, insistance) & \eng{They marched \textbf{in order to demand} justice.} \\
\textbf{so that + Sujet + can/could/will/would + BV} & \textbf{Sujets DIFFÉRENTS} & \eng{The law was passed \textbf{so that citizens could vote} freely.} \\
\textbf{for + Nom / -ing} & Fonction / Destination d'un objet & \eng{This hall is \textbf{for public meetings}.} / \eng{A tool \textbf{for voting}.} \\
\hline
\end{tabularx}
\end{center}

\end{propriete}

\begin{attention}[Négation du but — Piège majeur]
\begin{itemize}
\item \textbf{Correct (formel) :} \eng{in order \textbf{not to} + BV} / \eng{so as \textbf{not to} + BV}. \eng{He whispered \textbf{in order not to disturb} the meeting.}
\item \textbf{Incorrect / Maladroit :} \eng{not to} seul en début de proposition de but (\eng{He whispered not to disturb} $\leftarrow$ ambiguïté), \eng{for not to} (jamais), \eng{to not} (split infinitive à éviter au Bac).
\end{itemize}
\end{attention}

\begin{methode}[Choisir la bonne structure — Fiche méthode 3]
\begin{enumerate}
\item Identifier le sujet de l'action principale (S1) et le sujet de l'action de but (S2).
\item Si S1 = S2 $\rightarrow$ \eng{to} / \eng{in order to} / \eng{so as to} + \textbf{base verbale}.
\item Si S1 $\neq$ S2 $\rightarrow$ \eng{so that} + \textbf{S2 + modal (can/could/will/would)} + \textbf{base verbale}.
\item Si l'idée est la fonction d'un objet/lieu $\rightarrow$ \eng{for} + \textbf{nom / -ing}.
\item \textbf{Négation :} \eng{in order not to} / \eng{so as not to} (même sujet) ; \eng{so that} + S2 + \textbf{cannot/could not} (sujets différents).
\end{enumerate}
\end{methode}

\courssub{Comparaison français/anglais et pièges des francophones}

\begin{center}
\begin{tabularx}{\linewidth}{|X|X|X|}
\hline
\rowcolor{popblueL}
\textbf{Français (Calque fréquent)} & \textbf{Anglais correct} & \textbf{Explication / Règle} \\
\hline
« Il est venu pour m'aider. » $\rightarrow$ \eng{He came \textbf{for to help} me.} & \eng{He came \textbf{to help} me.} & \textbf{Jamais \eng{for to}.} \eng{to} + BV suffit (même sujet). \\
« Pour ne pas être en retard... » $\rightarrow$ \eng{\textbf{For not to} be late...} & \eng{\textbf{In order not to} be late...} / \eng{\textbf{So as not to} be late...} & \textbf{Jamais \eng{for not to}.} Négation = \eng{in order not to} / \eng{so as not to}. \\
« Les jeunes doivent participer à la vie locale. » $\rightarrow$ \eng{Youths must \textbf{participate to}...} & \eng{Youths must \textbf{take part in}...} / \eng{\textbf{participate in}...} & \textbf{Préposition \eng{in} après \eng{participate}.} \eng{Take part in} = phrasal verb 3 éléments. \\
« Il a dénoncé la corruption. » $\rightarrow$ \eng{He \textbf{denounced against} corruption.} & \eng{He \textbf{spoke out against} corruption.} & \eng{Denounce} est transitif direct (pas de \eng{against}). \eng{Speak out against} = phrasal verb. \\
« Elle a demandé de l'aide. » $\rightarrow$ \eng{She \textbf{demanded} help.} & \eng{She \textbf{asked for} help.} & \textbf{Faux ami :} \eng{demand} = exiger (autorité). \eng{Ask for} = demander. \\
« Nous ne pouvons pas tolérer cela. » $\rightarrow$ \eng{We cannot \textbf{support} this.} & \eng{We cannot \textbf{put up with} this.} & \eng{Support} = soutenir. \eng{Put up with} = tolérer (registre courant/standard). \\
\hline
\end{tabularx}
\end{center}

\begin{attention}[Les 5 erreurs qui coûtent des points au Bac — Synthèse]
\begin{enumerate}
\item \textbf{Pronom mal placé (Séparables).} \eng{He gave up it} $\times$ $\rightarrow$ \eng{He \textbf{gave it up}.} \checkmark (Règle absolue : Pronom = Milieu).
\item \textbf{Particule oubliée ou inventée (3 éléments).} \eng{Stand for their rights} $\times$ $\rightarrow$ \eng{\textbf{Stand up for} their rights.} \checkmark ; \eng{Participate to} $\times$ $\rightarrow$ \eng{\textbf{Take part in} / Participate in.} \checkmark
\item \textbf{Calque « pour » : \eng{for to} / \eng{for not to}.} \eng{He did it for to help} $\times$ $\rightarrow$ \eng{He did it \textbf{to help} / \textbf{in order to help}.} \checkmark ; \eng{She left for not to see him} $\times$ $\rightarrow$ \eng{She left \textbf{in order not to see} him.} \checkmark
\item \textbf{Faux amis lexicaux.} \eng{Demand} (exiger) vs \eng{Ask for} (demander) ; \eng{Assist} (aider) vs \eng{Attend} (assister à) ; \eng{Actually} (en fait) vs \eng{Currently} (actuellement).
\item \textbf{Morphologie : Particule conjuguée / Oubli du -ing après préposition.} \eng{She \textbf{carried on} work} $\times$ $\rightarrow$ \eng{She \textbf{carried on working}.} \checkmark ; \eng{Interested in \textbf{set up}} $\times$ $\rightarrow$ \eng{Interested in \textbf{setting up}.} \checkmark ; \eng{They \textbf{setted} up} $\times$ $\rightarrow$ \eng{They \textbf{set} up.} \checkmark (verbe irrégulier).
\end{enumerate}
\end{attention}

\courssub{Entraînement guidé : de la phrase au paragraphe}

\begin{exoresolu}[Analyse complète — Fiche méthode 1 appliquée]
\textbf{Énoncé.} Read the sentence: \eng{“When the authorities tried to \textbf{close down} the community radio station, the journalists of Buea did not \textbf{give in}. They \textbf{stood up for} press freedom and \textbf{carried on} broadcasting until the decision was \textbf{called off}.”}
\begin{enumerate}
\item Identify the five phrasal verbs.
\item Give the French meaning for each.
\item Classify each: Separable / Inseparable / Intransitive / 3-part.
\item Rewrite with a pronoun: \eng{The authorities called off the decision.} $\rightarrow$ \eng{The authorities \dots}
\end{enumerate}
\tcblower
\textbf{Solution détaillée.}
\textbf{1.} \eng{close down, give in, stood up for, carried on, called off}.
\textbf{2.} \eng{close down} = fermer ; \eng{give in} = céder ; \eng{stand up for} = défendre ; \eng{carry on} = continuer ; \eng{call off} = annuler.
\textbf{3.}
\begin{itemize}
\item \eng{close down} : Transitif \textbf{Séparable} (\eng{close it down}).
\item \eng{give in} : \textbf{Intransitif} (pas de COD).
\item \eng{stand up for} : \textbf{3 éléments (Inséparable)} (\eng{stand up for it}).
\item \eng{carry on} : \textbf{Inséparable} (suivi de -ing, pas de COD nominal direct).
\item \eng{call off} : Transitif \textbf{Séparable} (\eng{call it off}).
\end{itemize}
\textbf{4.} \eng{The authorities \textbf{called it off}.} (Séparable + pronom $\rightarrow$ milieu obligatoire).
\end{exoresolu}

\begin{exoresolu}[Construction et transformation — Fiche méthode 2]
\textbf{Énoncé.}
\begin{enumerate}
\item Complete: \eng{Citizens must stand \dots\ for their rights and speak \dots\ against injustice.}
\item Put in correct form: \eng{The NGO is interested in (set up) \dots\ a legal aid centre in Bamenda.}
\item Rewrite with \eng{it}: \eng{The mayor turned down the proposal.}
\item Translate: « Les jeunes devraient participer à la vie de leur communauté. »
\end{enumerate}
\tcblower
\textbf{Solution détaillée.}
\textbf{1.} \eng{stand \textbf{up} for} / \eng{speak \textbf{out} against}. (Particules fixes des phrasal verbs à 3 éléments).
\textbf{2.} \eng{\textbf{setting up}}. (Préposition \eng{in} + V-\textbf{ing} obligatoire).
\textbf{3.} \eng{The mayor turned \textbf{it} down.} (\eng{turn down} séparable + pronom $\rightarrow$ milieu).
\textbf{4.} \eng{Young people should \textbf{take part in} the life of their community.} (\eng{Participate in} possible mais \eng{take part in} phrasal verb attendu). \textbf{Piège :} \eng{participate to} $\times$.
\end{exoresolu}

\begin{exoresolu}[Expressions of Purpose — Fiche méthode 3]
\textbf{Énoncé.}
\begin{enumerate}
\item Join using \eng{to}: \eng{The students visited the court. They wanted to observe a trial.}
\item Join using \eng{so that}: \eng{The mayor spoke slowly. Everyone could follow his speech.}
\item Rewrite with \eng{in order not to}: \eng{She whispered because she did not want to disturb the meeting.}
\item Translate: « Il a créé une association pour défendre les droits des enfants. »
\end{enumerate}
\tcblower
\textbf{Solution détaillée.}
\textbf{1.} \eng{The students visited the court \textbf{to observe} a trial.} (S1=S2=\eng{students} ; suppression de \eng{wanted to}).
\textbf{2.} \eng{The mayor spoke slowly \textbf{so that everyone could follow} his speech.} (S1=\eng{mayor} $\neq$ S2=\eng{everyone} ; passé $\rightarrow$ \eng{could}).
\textbf{3.} \eng{She whispered \textbf{in order not to disturb} the meeting.} (Négation formelle du but, même sujet).
\textbf{4.} \eng{He \textbf{set up} an association \textbf{to defend} children's rights.} (\eng{set up} irrégulier : passé = \eng{set} ; même sujet $\rightarrow$ \eng{to defend} ; \eng{children's rights} = génitif).
\end{exoresolu}

\begin{exoresolu}[Combinaison Phrasal Verbs + But — Fiche méthode 4]
\textbf{Énoncé.} Write a short paragraph (5--6 lines) about a citizenship action in your town, using at least three phrasal verbs and two different expressions of purpose.
\tcblower
\textbf{Modèle de production (Niveau Bac Excellent).}
\eng{Last month, students from Government High School Bamenda \textbf{set up} a hygiene club \textbf{in order to} \textbf{clean up} the neighbourhood market. Every Saturday, they \textbf{turn up} early \textbf{to} collect waste and sensitise traders. They also organise talks \textbf{so that} the population \textbf{can understand} the health risks of pollution. Their action has encouraged many residents to \textbf{join in} and \textbf{give back} to their community.}

\textbf{Analyse linguistique :}
\begin{itemize}
\item \textbf{Phrasal verbs (5) :} \eng{set up} (créer, séparable), \eng{clean up} (nettoyer, séparable), \eng{turn up} (arriver, intransitif), \eng{join in} (participer, intransitif), \eng{give back} (contribuer, séparable).
\item \textbf{But varié (3 structures) :} \eng{in order to} + BV (formel, S1=S2), \eng{to} + BV (standard, S1=S2), \eng{so that} + \eng{population can understand} (S1 $\neq$ S2).
\item \textbf{Grammaire :} \eng{has encouraged} (present perfect, conséquence présente), \eng{can} (présent après présent \eng{organise}).
\end{itemize}
\end{exoresolu}

\begin{exoresolu}[Traduction guidée — Fiche méthode 5]
\textbf{Énoncé.} Translate into English:
\begin{enumerate}
\item « Les citoyens doivent dénoncer la corruption. »
\item « Ils ont reporté la réunion à vendredi. »
\item « Nous ne pouvons pas tolérer cette injustice. »
\item « Elle a demandé de l'aide à un avocat. »
\end{enumerate}
\tcblower
\textbf{Solution détaillée.}
\textbf{1.} \eng{Citizens must \textbf{speak out against} corruption.} (\eng{Denounce} transitif direct, pas \eng{against} ; \eng{speak out against} = phrasal verb 3 éléments).
\textbf{2.} \eng{They \textbf{put off} the meeting \textbf{until} Friday.} (\eng{Put off} séparable ; \eng{until} pour la date limite, pas \eng{to}).
\textbf{3.} \eng{We cannot \textbf{put up with} this injustice.} (\eng{Put up with} = tolérer, 3 éléments inséparable ; \eng{tolerate} correct mais moins ciblé leçon).
\textbf{4.} \eng{She \textbf{asked a lawyer for} help.} (Structure \eng{ask someone for something} ; \textbf{Jamais} \eng{demand help to a lawyer} — \eng{demand} = exiger).
\end{exoresolu}

\begin{exercice}[Entraînement personnel — Devoir / Classe]
\begin{enumerate}
\item \textbf{Classify} as Separable (Sep), Inseparable (Insep), Intransitive (Intr), or 3-part (3P): \eng{look after, give up, turn out, put up with, call off, stand up for, join in, set up}.
\item \textbf{Complete} with the correct particle(s): \eng{a) The government wants to \_\_\_\_ down the illegal market. b) Don't \_\_\_\_ in to pressure. c) We must \_\_\_\_ up \_\_\_\_ our leaders. d) They \_\_\_\_ on working despite the rain. e) The match was \_\_\_\_ off due to rain.}
\item \textbf{Rewrite} placing the pronoun correctly: \eng{a) She turned down the offer. (it) b) They set up the committee. (it) c) He puts up with the noise. (it) d) We called off the strike. (it)}
\item \textbf{Join} the sentences using the \textbf{correct expression of purpose}:
\begin{itemize}
\item The NGO organised a workshop. \textit{They wanted to train volunteers.} (\eng{to})
\item The teacher explained the law. \textit{The students could know their rights.} (\eng{so that})
\item He left early. \textit{He didn't want to miss the bus.} (\eng{in order not to})
\item This building is used. \textit{Public meetings take place here.} (\eng{for})
\end{itemize}
\item \textbf{Translate} into English:
\begin{itemize}
\item « L'association a été créée pour aider les orphelins. »
\item « Nous devons nous battre pour la justice. »
\item « Ils ont annulé la manifestation pour éviter la violence. »
\item « Le ministre a promis de veiller sur les personnes âgées. »
\end{itemize}
\end{enumerate}
\end{exercice}

\begin{corrige}[Exercice 46]
\textbf{1. Classification :}
\eng{look after} (Insep), \eng{give up} (Sep), \eng{turn out} (Intr), \eng{put up with} (3P), \eng{call off} (Sep), \eng{stand up for} (3P), \eng{join in} (Intr), \eng{set up} (Sep).

\textbf{2. Complétion :}
a) \textbf{close} down \quad b) \textbf{give} in \quad c) \textbf{stand} up \textbf{for} \quad d) \textbf{carry} on \quad e) \textbf{call} off

\textbf{3. Pronoms :}
a) She turned \textbf{it} down. \quad b) They set \textbf{it} up. \quad c) He puts up with \textbf{it}. (3P, pronom après préposition) \quad d) We called \textbf{it} off.

\textbf{4. But :}
a) The NGO organised a workshop \textbf{to train} volunteers. (Même sujet)
b) The teacher explained the law \textbf{so that the students could know} their rights. (Sujets différents, passé $\rightarrow$ could)
c) He left early \textbf{in order not to miss} the bus. (Négation, même sujet)
d) This building is used \textbf{for public meetings} / \textbf{for holding public meetings}. (Fonction lieu)

\textbf{5. Traduction :}
a) The association \textbf{was set up} \textbf{to help} orphans. (Passif + but même sujet agent implicite).
b) We must \textbf{fight for} justice. (\eng{Fight for} inséparable).
c) They \textbf{called off} the demonstration \textbf{in order to avoid} violence / \textbf{so as to avoid} violence. (Même sujet, formel).
d) The minister promised to \textbf{look after} the elderly. (\eng{Look after} inséparable ; \eng{the elderly} = les personnes âgées).
\end{corrige}

\begin{aretenir}[Synthèse Bac — Lesson 46]
\begin{itemize}
\item \textbf{Phrasal Verb} = Verbe + Particule (sens idiomatique). Apprendre \textbf{verbe + particule(s) + type (Sep/Insep/3P/Intr)}.
\item \textbf{Règle d'or pronom :} Séparable $\rightarrow$ Pronom \textbf{AU MILIEU} (\eng{give it up}). Inséparable/3P $\rightarrow$ Pronom \textbf{APRÈS} (\eng{look after them, stand up for it}).
\item \textbf{But (Purpose) :} S1=S2 $\rightarrow$ \eng{to / in order to / so as to} + BV. S1$\neq$S2 $\rightarrow$ \eng{so that} + S2 + \eng{can/could}. Négation $\rightarrow$ \eng{in order not to / so as not to}. Fonction $\rightarrow$ \eng{for} + N/-ing.
\item \textbf{Citoyenneté :} Maîtriser la liste thématique (\eng{stand up for, speak out against, take part in, set up, put up with, give back...}).
\item \textbf{Anti-calques :} \eng{participate in / take part in} (pas \eng{to}), \eng{ask for} (pas \eng{demand}), \eng{put up with} (pas \eng{support}), \eng{until} date (pas \eng{to}).
\end{itemize}
\end{aretenir}

\newpage

\coursec{Exercices}

\courssub{Je vérifie mes connaissances}

Avant de commencer, rappelle-toi des deux points essentiels de la leçon : un \cle{phrasal verb} (verbe à particule) associe un verbe et une particule (\eng{up, out, for, after...}) dont le sens est souvent différent du verbe seul ; et pour exprimer le \cle{but}, l'anglais utilise \eng{to}, \eng{in order to} ou \eng{so as to} + base verbale. Les exercices qui suivent vérifient ces deux acquis.

\begin{exercice}[QCM : la bonne particule]
Choose the correct particle to complete each sentence.
\begin{enumerate}
\item The volunteers decided to carry \_\_\_\_ a campaign against corruption. \hfill (a) on \quad (b) out \quad (c) up \quad (d) off
\item We must never give \_\_\_\_ the fight for justice. \hfill (a) in \quad (b) out \quad (c) up \quad (d) away
\item The lawyer agreed to look \_\_\_\_ the case of the detained journalist. \hfill (a) after \quad (b) into \quad (c) for \quad (d) over
\item Many citizens stood \_\_\_\_ during the meeting to defend their rights. \hfill (a) up \quad (b) down \quad (c) off \quad (d) away
\end{enumerate}
\end{exercice}

\begin{exercice}[Vrai ou faux ?]
Say whether each sentence is TRUE or FALSE, and correct the false ones.
\begin{enumerate}
\item In a phrasal verb, the particle can change the meaning of the verb completely.
\item We can say \eng{They carried out it successfully.}
\item \eng{In order to} and \eng{so as to} are used to express purpose.
\item The negative form of \eng{in order to} is \eng{in order don't}.
\end{enumerate}
\end{exercice}

\begin{exercice}[Vocabulaire : phrasal verbs de la citoyenneté]
Match each phrasal verb with its French equivalent.
\begin{center}
\begin{tabularx}{\linewidth}{|X|X|}
\hline
\rowcolor{popblueL} \textbf{Phrasal verb} & \textbf{Traduction française} \\
\hline
1. to stand up for & a. abandonner, renoncer \\
\hline
2. to give up & b. défendre (une cause, une personne) \\
\hline
3. to look after & c. mener, réaliser (une action) \\
\hline
4. to carry out & d. s'occuper de, prendre soin de \\
\hline
5. to find out & e. découvrir, se renseigner \\
\hline
\end{tabularx}
\end{center}
\end{exercice}

\begin{exercice}[À compléter : verbe ou particule ?]
Complete each sentence with the correct form of the verb in brackets or with the missing particle.
\begin{enumerate}
\item The students \_\_\_\_ (organise) a debate on human rights last week.
\item If you don't know your rights, you should \_\_\_\_ out more about the law.
\item She \_\_\_\_ (look) after the orphans since 2020.
\item The government promised to put \_\_\_\_ new proposals to protect women.
\end{enumerate}
\end{exercice}

\courssub{Je m'entraîne}

\begin{exercice}[Traduction vers l'anglais]
Traduire en anglais les phrases suivantes.
\begin{enumerate}
\item Les citoyens doivent défendre leurs droits afin de provoquer un vrai changement, et ils ne doivent jamais abandonner.
\item Les jeunes de Bamenda ont organisé une marche pour sensibiliser la population.
\item Elle s'occupe des enfants des rues afin qu'ils puissent retourner à l'école.
\end{enumerate}
\end{exercice}

\begin{exercice}[Transformation : exprimer le but]
Rewrite each pair of sentences as one sentence using \eng{to}, \eng{in order to} or \eng{so as to}.
\begin{enumerate}
\item Amina joined the NGO. She wanted to defend children's rights.
\item The mayor went to the village. He wanted to listen to the inhabitants' complaints.
\item They left early. They did not want to miss the public meeting. (use a negative form)
\item I am learning English. I want to work for an international organisation.
\end{enumerate}
\end{exercice}

Lis d'abord le texte suivant, puis fais l'exercice de complétion.

\begin{textebox}[A citizens' meeting in Douala]
Last Saturday, the inhabitants of Bonaberi decided to stand up for their rights. They carried out a peaceful march to ask for clean water. A young leader put forward the idea of creating a neighbourhood committee. The committee will look after the poorest families and find out why the water project was abandoned. Nobody wants to give up until the authorities take the situation seriously.
\end{textebox}

\begin{exercice}[Texte à trous]
Copy the text above and fill in the blanks with the correct particle: \emph{up, out, forward, after}. (Attention : certaines particules sont utilisées deux fois.)
\end{exercice}

\begin{exercice}[Appariement : début et fin de phrase]
Match each beginning (1--5) with the correct ending (a--e) to make meaningful sentences.
\begin{center}
\begin{tabularx}{\linewidth}{|X|X|}
\hline
\rowcolor{popblueL} \textbf{Beginning} & \textbf{Ending} \\
\hline
1. The activists set up a committee & a. be heard by everyone in the hall. \\
\hline
2. She went to the police station & b. miss this opportunity to change things. \\
\hline
3. We must act now so as not to & c. who arrived from the border last month. \\
\hline
4. He spoke loudly in order to & d. to report a case of discrimination. \\
\hline
5. They look after the refugees & e. to protect the environment of the city. \\
\hline
\end{tabularx}
\end{center}
\end{exercice}

\courssub{Je me prépare au Bac}

Lis attentivement le texte suivant, puis réponds aux questions de l'exercice type Bac.

\begin{textebox}[Youth and citizenship in Cameroon]
In many towns of Cameroon, young people are no longer waiting for change: they are creating it. In Yaoundé, a group of students has set up an association called “Voices of Tomorrow”. Their aim is to carry out civic education programmes in secondary schools, so as to help teenagers find out about their rights and duties as citizens.

Every weekend, the members visit a different school. They organise debates, put forward new ideas and encourage students to stand up for justice and equality. “Many young people think citizenship only means voting,” explains Clarisse, the president of the association. “We want to show them that it also means looking after their community and never giving up when things become difficult.”

The association also works with local radio stations in order to reach young people in rural areas. Thanks to these broadcasts, thousands of listeners have learnt about human rights. The members know the road is long, but they are determined: they will not give up until every young Cameroonian knows that citizenship is built day after day.
\end{textebox}

\begin{exercice}[Compréhension et langue — type Bac]
Read the text above carefully and answer the questions.

\textbf{A. Comprehension questions.}
\begin{enumerate}
\item What is the name of the association created by the students, and what is its aim?
\item According to Clarisse, what does citizenship mean besides voting? Quote from the text.
\item How does the association reach young people in rural areas?
\item Say whether the following statement is TRUE or FALSE and justify with a quotation: \emph{The members of the association are ready to stop their activities soon.}
\end{enumerate}

\textbf{B. Language work.}
\begin{enumerate}
\item Find in the text three phrasal verbs and give their meaning in French.
\item Rewrite using \eng{so as to}: \emph{The members visit schools every weekend. They want to help teenagers discover their rights.}
\item Find and correct the mistake: \emph{The students want to find about out their rights.}
\end{enumerate}
\end{exercice}

\begin{exercice}[Traduction et correction d'erreurs — type Bac]
\textbf{A. Traduire en anglais.}
\begin{enumerate}
\item Les citoyens doivent défendre leurs droits afin de provoquer un vrai changement, et ils ne doivent jamais abandonner.
\item Pour lutter contre la corruption, les journalistes ont mené une longue enquête.
\end{enumerate}

\textbf{B. Find and correct the mistake in each sentence.}
\begin{enumerate}
\item She looks the children after.
\item He joined for to help the victims.
\item They carried out it successfully.
\item We must act now for not to miss this opportunity.
\end{enumerate}
\end{exercice}

\begin{exercice}[Rédaction guidée — type Bac]
Write a paragraph of about \textbf{80 words} about what young Cameroonians can do to promote human rights in their community.

Your paragraph must contain:
\begin{itemize}
\item at least \textbf{three phrasal verbs} (for example: \emph{stand up for, carry out, give up, look after, set up, find out});
\item at least \textbf{two expressions of purpose} (\emph{to, in order to, so as to}, including one negative form);
\item a clear conclusion expressing your personal opinion.
\end{itemize}
\end{exercice}

\begin{methode}[Réussir le paragraphe type Bac]
\begin{enumerate}
\item \textbf{Planifie} : note trois actions concrètes (ex. créer un club, organiser une campagne, aider les plus pauvres) et associe à chacune un phrasal verb.
\item \textbf{Rédige} : une phrase d'introduction, deux ou trois phrases d'actions avec \eng{to / in order to / so as to}, dont une forme négative (\eng{so as not to}), puis une conclusion personnelle (\eng{In my opinion...}).
\item \textbf{Vérifie} : compte tes phrasal verbs (au moins 3) et tes expressions de but (au moins 2) ; contrôle l'ordre des mots (pronom entre verbe et particule) et le nombre de mots (environ 80).
\end{enumerate}
\end{methode}

\begin{exemplebox}[Modèle de paragraphe]
\eng{Young Cameroonians can do a lot to promote human rights in their community. First, they can set up youth clubs in order to discuss problems such as discrimination and corruption. They can also carry out awareness campaigns in schools and markets, so as to help people find out about their rights. Besides, they should look after vulnerable people, like orphans and street children, and they must act now so as not to lose the future of their country. In my opinion, young people must never give up, because real change always begins with citizens who stand up for justice.}

\emph{Traduction :} « Les jeunes Camerounais peuvent faire beaucoup pour promouvoir les droits de l'homme dans leur communauté. D'abord, ils peuvent créer des clubs de jeunes afin de discuter de problèmes comme la discrimination et la corruption. Ils peuvent aussi mener des campagnes de sensibilisation dans les écoles et les marchés, pour aider les gens à découvrir leurs droits. En outre, ils devraient s'occuper des personnes vulnérables, comme les orphelins et les enfants des rues, et ils doivent agir maintenant pour ne pas perdre l'avenir de leur pays. À mon avis, les jeunes ne doivent jamais abandonner, car le vrai changement commence toujours avec des citoyens qui défendent la justice. »
\end{exemplebox}

\begin{attention}[Les trois pièges à éviter au Bac]
\begin{itemize}
\item \textbf{Le pronom mal placé} : on écrit \eng{They carried it out}, jamais \eng{They carried out it}. Avec un nom, les deux ordres sont possibles : \eng{They carried out the plan} ou \eng{They carried the plan out}.
\item \textbf{Le but avec \eng{for}} : \eng{for} + nom exprime le but (\eng{He came for the meeting}), mais devant un verbe il faut \eng{to / in order to / so as to}, jamais \eng{for to}.
\item \textbf{La négation} : \eng{not} se place devant \eng{to} : \eng{so as not to}, \eng{in order not to} — jamais \eng{to not} ni \eng{don't to}.
\end{itemize}
\end{attention}

\newpage

\coursec{Corrigés des exercices}

\begin{corrige}[Exercice 1 — QCM : la bonne particule]
\begin{enumerate}
\item (b) \eng{out} — \eng{to carry out a campaign} = mener, réaliser une campagne. Phrase complète : \eng{The volunteers decided to carry out a campaign against corruption.} (« Les volontaires ont décidé de mener une campagne contre la corruption. »)
\item (c) \eng{up} — \eng{to give up} = abandonner, renoncer. \eng{We must never give up the fight for justice.} (« Nous ne devons jamais abandonner le combat pour la justice. »)
\item (b) \eng{into} — \eng{to look into a case} = enquêter sur une affaire. \eng{The lawyer agreed to look into the case of the detained journalist.} (« L'avocat a accepté d'enquêter sur l'affaire du journaliste détenu. »)
\item (a) \eng{up} — \eng{to stand up} = se lever. \eng{Many citizens stood up during the meeting to defend their rights.} (« Beaucoup de citoyens se sont levés pendant la réunion pour défendre leurs droits. »)
\end{enumerate}
\textbf{Point de vigilance :} ne pas confondre \eng{carry on} (continuer), \eng{carry out} (mener à bien) et \eng{carry off} (emporter, remporter). La particule change tout le sens.
\end{corrige}

\begin{corrige}[Exercice 2 — Vrai ou faux ?]
\begin{enumerate}
\item TRUE. La particule peut modifier radicalement le sens : \eng{to look} (regarder), \eng{to look after} (s'occuper de), \eng{to look for} (chercher), \eng{to look into} (enquêter sur).
\item FALSE. Avec un pronom, celui-ci se place obligatoirement \emph{entre} le verbe et la particule : \eng{They carried it out successfully.} (« Ils l'ont menée à bien. »)
\item TRUE. \eng{In order to} et \eng{so as to} expriment le but : \eng{She studies hard in order to pass the Bac.} (« Elle travaille dur afin de réussir le Bac. »)
\item FALSE. La négation se fait avec \eng{not} placé devant \eng{to} : \eng{in order not to}, \eng{so as not to}. Exemple : \eng{He left early so as not to be late.} (« Il est parti tôt pour ne pas être en retard. »)
\end{enumerate}
\end{corrige}

\begin{corrige}[Exercice 3 — Vocabulaire : phrasal verbs de la citoyenneté]
1 -- b (\eng{to stand up for} = défendre une cause, une personne) ; 2 -- a (\eng{to give up} = abandonner) ; 3 -- d (\eng{to look after} = s'occuper de) ; 4 -- c (\eng{to carry out} = mener, réaliser) ; 5 -- e (\eng{to find out} = découvrir, se renseigner).

\textbf{Point de vigilance :} \eng{to stand up for} est un phrasal verb à \emph{deux} particules (\eng{up} + \eng{for}) ; il est inséparable : \eng{She stood up for her friend.} (« Elle a défendu son amie. »)
\end{corrige}

\begin{corrige}[Exercice 4 — À compléter : verbe ou particule ?]
\begin{enumerate}
\item \eng{organised} — le repère temporel \eng{last week} impose le prétérit simple : \eng{The students organised a debate on human rights last week.}
\item \eng{find} — \eng{to find out} = se renseigner : \eng{If you don't know your rights, you should find out more about the law.}
\item \eng{has looked} — \eng{since 2020} impose le present perfect : \eng{She has looked after the orphans since 2020.} (« Elle s'occupe des orphelins depuis 2020. »)
\item \eng{forward} — \eng{to put forward} = proposer, avancer : \eng{The government promised to put forward new proposals to protect women.}
\end{enumerate}
\end{corrige}

\begin{corrige}[Exercice 5 — Traduction vers l'anglais]
\begin{enumerate}
\item \eng{Citizens must stand up for their rights in order to bring about real change, and they must never give up.} — « afin de » = \eng{in order to} ; « provoquer » = \eng{to bring about} (phrasal verb) ; « abandonner » = \eng{to give up}.
\item \eng{The young people of Bamenda organised a march to raise awareness among the population.} (ou : \eng{... to sensitise the population.}) — « pour sensibiliser » = \eng{to} + base verbale.
\item \eng{She looks after street children so that they can go back to school.} — « afin qu'ils puissent » = \eng{so that} + proposition : on ne peut pas utiliser \eng{in order to} car le sujet de la seconde action (\eng{they}) est différent du premier (\eng{she}).
\end{enumerate}
\textbf{Point de vigilance :} \eng{in order to / so as to} n'est possible que si le sujet des deux actions est le même ; sinon, utiliser \eng{so that} + sujet + \eng{can/could}.
\end{corrige}

\begin{corrige}[Exercice 6 — Transformation : exprimer le but]
\begin{enumerate}
\item \eng{Amina joined the NGO to defend children's rights.} (ou \eng{in order to defend children's rights.})
\item \eng{The mayor went to the village in order to listen to the inhabitants' complaints.}
\item \eng{They left early so as not to miss the public meeting.} (ou \eng{in order not to miss...}) — forme négative : \eng{not} devant \eng{to}.
\item \eng{I am learning English to work for an international organisation.} (ou \eng{so as to work...})
\end{enumerate}
\textbf{Point de vigilance :} jamais \eng{for to} + verbe ; le but devant un verbe s'exprime avec \eng{to}, \eng{in order to} ou \eng{so as to}.
\end{corrige}

\begin{corrige}[Exercice 7 — Texte à trous]
(1) \eng{up} — \eng{stand up for their rights} (défendre leurs droits) ; (2) \eng{out} — \eng{carried out a peaceful march} (mené une marche pacifique) ; (3) \eng{forward} — \eng{put forward the idea} (proposé l'idée) ; (4) \eng{after} — \eng{look after the poorest families} (s'occuper des familles les plus pauvres) ; (5) \eng{out} — \eng{find out why} (découvrir pourquoi) ; (6) \eng{up} — \eng{give up} (abandonner).

Texte complet : \eng{Last Saturday, the inhabitants of Bonaberi decided to stand up for their rights. They carried out a peaceful march to ask for clean water. A young leader put forward the idea of creating a neighbourhood committee. The committee will look after the poorest families and find out why the water project was abandoned. Nobody wants to give up until the authorities take the situation seriously.}
\end{corrige}

\begin{corrige}[Exercice 8 — Appariement : début et fin de phrase]
1 -- e : \eng{The activists set up a committee to protect the environment of the city.} (\eng{to set up} = créer, fonder)
2 -- d : \eng{She went to the police station to report a case of discrimination.}
3 -- b : \eng{We must act now so as not to miss this opportunity to change things.} — remarque la forme négative du but : \eng{so as not to miss} (« pour ne pas manquer »).
4 -- a : \eng{He spoke loudly in order to be heard by everyone in the hall.} — but + passif : \eng{in order to be heard}.
5 -- c : \eng{They look after the refugees who arrived from the border last month.}
\end{corrige}

\begin{corrige}[Exercice 9 — Compréhension et langue — type Bac]
\textbf{A. Comprehension.}
\begin{enumerate}
\item The association is called “Voices of Tomorrow”. Its aim is to carry out civic education programmes in secondary schools, so as to help teenagers find out about their rights and duties as citizens.
\item According to Clarisse, citizenship “also means looking after their community and never giving up when things become difficult.”
\item The association works with local radio stations in order to reach young people in rural areas.
\item FALSE. The text says: “they will not give up until every young Cameroonian knows that citizenship is built day after day.”
\end{enumerate}
\textbf{B. Language work.}
\begin{enumerate}
\item Trois phrasal verbs (au choix) : \eng{to set up} (créer, fonder), \eng{to carry out} (mener, réaliser), \eng{to find out} (découvrir, se renseigner), \eng{to stand up for} (défendre), \eng{to look after} (s'occuper de), \eng{to give up} (abandonner).
\item \eng{The members visit schools every weekend so as to help teenagers discover their rights.}
\item Erreur : l'ordre des mots. Correction : \eng{The students want to find out about their rights.} — \eng{to find out about} = se renseigner sur ; la particule \eng{out} reste collée au verbe, et la préposition \eng{about} précède le complément.
\end{enumerate}
\textbf{Barème indicatif (sur 10) :} A. 4 questions $\times$ 1,5 pt = 6 pts (réponse correcte 1 pt + citation ou justification 0,5 pt) ; B. 1 : 1,5 pt (0,5 pt par phrasal verb correctement traduit) ; B. 2 : 1,5 pt ; B. 3 : 1 pt.
\textbf{Points de vigilance :} en compréhension, répondre par des phrases complètes en anglais et citer le texte entre guillemets pour justifier ; ne pas recopier tout un paragraphe.
\end{corrige}

\begin{corrige}[Exercice 10 — Traduction et correction d'erreurs — type Bac]
\textbf{A. Traduction.}
\begin{enumerate}
\item \eng{Citizens must stand up for their rights in order to bring about real change, and they must never give up.}
\item \eng{To fight against corruption, the journalists carried out a long investigation.}
\end{enumerate}
\textbf{B. Correction.}
\begin{enumerate}
\item \eng{She looks after the children.} — \eng{to look after} est un phrasal verb \emph{inséparable} : avec un nom comme complément, on ne peut pas placer celui-ci entre le verbe et la particule.
\item \eng{He joined to help the victims.} (ou \eng{in order to help the victims.}) — le but ne s'exprime jamais avec \eng{for to}.
\item \eng{They carried it out successfully.} — avec un pronom, celui-ci se place obligatoirement entre le verbe et la particule.
\item \eng{We must act now so as not to miss this opportunity.} (ou \eng{in order not to miss...}) — la négation du but : \eng{not} devant \eng{to}.
\end{enumerate}
\textbf{Barème indicatif (sur 10) :} A. 2 phrases $\times$ 2,5 pts = 5 pts ; B. 4 phrases $\times$ 1,25 pt = 5 pts (0,5 pt pour l'erreur repérée + 0,75 pt pour la correction).
\textbf{Points de vigilance :} distinguer les phrasal verbs séparables (\eng{carry out} : \eng{carry it out}) des inséparables (\eng{look after} : jamais \eng{look the children after}).
\end{corrige}

\begin{corrige}[Exercice 11 — Rédaction guidée — type Bac]
\textbf{Barème indicatif (sur 10) :} respect de la consigne et du nombre de mots : 2 pts ; au moins 3 phrasal verbs correctement employés : 3 pts ; au moins 2 expressions de but dont une négative : 2 pts ; correction grammaticale et orthographe : 2 pts ; clarté de la conclusion personnelle : 1 pt.

\textbf{Proposition de rédaction modèle :}

\eng{Young Cameroonians can do a lot to promote human rights in their community. First, they can set up youth clubs in order to discuss problems such as discrimination and corruption. They can also carry out awareness campaigns in schools and markets, so as to help people find out about their rights. Besides, they should look after vulnerable people, like orphans and street children, and they must act now so as not to lose the future of their country. In my opinion, young people must never give up, because real change always begins with citizens who stand up for justice.}

\emph{Traduction :} « Les jeunes Camerounais peuvent faire beaucoup pour promouvoir les droits de l'homme dans leur communauté. D'abord, ils peuvent créer des clubs de jeunes afin de discuter de problèmes comme la discrimination et la corruption. Ils peuvent aussi mener des campagnes de sensibilisation dans les écoles et les marchés, pour aider les gens à découvrir leurs droits. En outre, ils devraient s'occuper des personnes vulnérables, comme les orphelins et les enfants des rues, et ils doivent agir maintenant pour ne pas perdre l'avenir de leur pays. À mon avis, les jeunes ne doivent jamais abandonner, car le vrai changement commence toujours avec des citoyens qui défendent la justice. »

\textbf{Points de vigilance :}
\begin{itemize}
\item Vérifier que chaque phrasal verb est bien formé : \eng{set up}, \eng{carry out}, \eng{find out}, \eng{look after}, \eng{give up}, \eng{stand up for}.
\item La forme négative du but : \eng{so as not to lose}, jamais \eng{to not lose} ni \eng{for not to lose}.
\item Introduire l'opinion personnelle avec \eng{In my opinion...} ou \eng{I think that...}, et non avec une traduction mot à mot du français.
\end{itemize}
\end{corrige}

\newpage

\coursec{Fiche bilan}

\begin{popfigure}
\begin{tikzpicture}[mindmap, grow cyclic, every node/.style={concept, text=white, font=\small\bfseries},
root concept/.append style={concept color=popblue, font=\normalsize\bfseries, minimum size=3.2cm},
level 1/.append style={level distance=3.6cm, sibling angle=51},
level 1 concept/.append style={minimum size=2.4cm}]
\node[root concept]{Phrasal verbs \& purpose}
child[concept color=poporange]{node[align=center]{Formation\\ verb + particle}}
child[concept color=popgreen]{node[align=center]{Separable\\ vs inseparable}}
child[concept color=poppurple]{node[align=center]{Citizenship\\ stand up for, vote}}
child[concept color=poppink]{node[align=center]{Purpose\\ to / in order to}}
child[concept color=popgold]{node[align=center]{Purpose\\ so as to / so that}}
child[concept color=popblue!70]{node[align=center]{Pièges\\ des francophones}}
child[concept color=popgreen!70]{node[align=center]{Méthode\\ phrase $\rightarrow$ texte}};
\end{tikzpicture}
\legende{Carte mentale de la leçon : les phrasal verbs et l'expression du but.}
\end{popfigure}

\begin{bilan}
\textbf{1. Lexique clé de la citoyenneté et des droits}
\begin{itemize}
\item \eng{to stand up for one's rights} = défendre ses droits ; \eng{to speak out against injustice} = dénoncer l'injustice
\item \eng{to vote for a candidate} = voter pour un candidat ; \eng{to take part in elections} = participer aux élections
\item \eng{to fight against corruption} = lutter contre la corruption ; \eng{to look after the community} = prendre soin de la communauté
\item \eng{to bring about change} = provoquer le changement ; \eng{to put an end to discrimination} = mettre fin à la discrimination
\item \eng{a citizen} (un citoyen), \eng{human rights} (les droits de l'homme), \eng{freedom of speech} (la liberté d'expression)
\end{itemize}

\textbf{2. Règles de grammaire à connaître par cœur}
\begin{center}
\begin{tabularx}{\linewidth}{|X|X|X|}
\hline
\rowcolor{popblueL} Règle & Exemple & Traduction \\
\hline
Phrasal verb = verbe + particule (adverbe ou préposition) & \eng{Citizens stood up for their rights.} & Les citoyens ont défendu leurs droits. \\
\hline
Séparable : le complément nominal peut se placer avant ou après la particule & \eng{They called off the meeting. / They called the meeting off.} & Ils ont annulé la réunion. \\
\hline
Avec un pronom, il est obligatoirement entre le verbe et la particule & \eng{They called it off.} (jamais \emph{called off it}) & Ils l'ont annulée. \\
\hline
Inséparable : la particule reste collée au verbe & \eng{We must fight against injustice.} & Nous devons lutter contre l'injustice. \\
\hline
But affirmatif : \eng{to} / \eng{in order to} / \eng{so as to} + base verbale & \eng{She went to the polling station to vote.} & Elle est allée au bureau de vote pour voter. \\
\hline
But négatif : \eng{in order not to} / \eng{so as not to} & \eng{He spoke quietly so as not to disturb anyone.} & Il parlait doucement pour ne déranger personne. \\
\hline
But avec sujet différent : \eng{so that} + sujet + modal & \eng{The teacher explained so that everyone could understand.} & Le professeur a expliqué pour que tout le monde comprenne. \\
\hline
\end{tabularx}
\end{center}

\textbf{3. Méthodes express}
\begin{itemize}
\item Identifier un phrasal verb : chercher verbe + particule dont le sens global diffère du verbe seul (\eng{look up} $\neq$ \eng{look}).
\item Exprimer le but : question « pourquoi ? » $\rightarrow$ \eng{to} + BV (même sujet) ; sujets différents $\rightarrow$ \eng{so that} + phrase.
\item Transformation : « Elle étudie pour réussir » $\rightarrow$ \eng{She studies in order to pass.} (jamais \emph{for to pass}).
\item Au passé, seul le verbe se conjugue ; la particule ne change jamais : \eng{stood up for}, \eng{called off}.
\end{itemize}
\end{bilan}

\begin{aretenir}[Checklist avant le Bac]
\begin{itemize}
\item[$\square$] Je sais reconnaître un phrasal verb dans un texte et en deviner le sens grâce au contexte.
\item[$\square$] Je connais au moins dix phrasal verbs liés à la citoyenneté et aux droits de l'homme.
\item[$\square$] Je distingue les phrasal verbs séparables des inséparables.
\item[$\square$] Je place correctement le pronom complément : \eng{They called it off.} (jamais après la particule).
\item[$\square$] Je conjugue le verbe au temps voulu sans toucher à la particule : \eng{stand up for} $\rightarrow$ \eng{stood up for}.
\item[$\square$] J'exprime le but affirmatif avec \eng{to}, \eng{in order to} ou \eng{so as to} + base verbale.
\item[$\square$] J'exprime le but négatif avec \eng{in order not to} / \eng{so as not to}.
\item[$\square$] J'utilise \eng{so that} + sujet + \eng{can/could/will/would} quand les sujets sont différents.
\item[$\square$] J'évite le piège français \emph{« for to + verbe »} : on dit \eng{He came to help}, jamais \emph{for to help}.
\item[$\square$] Je ne confonds pas \eng{for} + nom (but devant un nom : \eng{for the election}) et \eng{to} + verbe.
\item[$\square$] Je sais rédiger un court paragraphe sur la citoyenneté en employant au moins trois phrasal verbs et une expression de but.
\item[$\square$] Je relève et corrige les erreurs de particules dans mes propres productions écrites.
\end{itemize}
\end{aretenir}

\findecours

\end{document}
